%% Volume I, Chapter 25: Linear Systems of Differential Equations
\chapter{Linear Systems of Differential Equations}
\label{ch:linear-ode-systems}

\abstract{Coupled springs, mixing tanks, reacting chemicals, predators and their prey: wherever several quantities change together and act on one another, the laws of change form a system of differential equations. Near an equilibrium such a system behaves like a linear one, $\mathbf y'=A\mathbf y$, and everything about a linear system is written in the eigenvalues of its matrix. Their real parts decide whether a disturbance dies away or grows, their imaginary parts how fast the system turns or swings, and the eigenvectors give the directions in which it moves. This chapter puts the linear algebra of Part~III to work: diagonalisation, the Jordan form, the matrix exponential, the spectral theorem and the Sylvester equation all return as statements about motion.}

\tikzset{ch25arrow/.style={-{Latex[length=1.7mm,width=1.3mm]}, black, line width=0.5pt}, ch25orbit/.style={black, line width=0.5pt}}

The two previous chapters met systems of differential equations several times without studying them for their own sake. \Cref{thm:linear-system} solved the system $\mathbf y'=A\mathbf y$ with a constant matrix $A$ by the matrix exponential, $\mathbf y(t)=e^{tA}\mathbf y(0)$, and \cref{prop:exp-jordan} showed that the entries of $e^{tA}$ are combinations of the functions $t^je^{\mu t}$ for the eigenvalues $\mu$ of $A$. \Cref{thm:linear-existence} proved that every system $\mathbf y'=A(t)\mathbf y+\mathbf b(t)$ with continuous coefficients has exactly one solution with a given initial value, on the whole interval on which the coefficients are continuous, and \cref{ch:linear-higher-order} used this theorem for equations of higher order through their companion systems. \Cref{ch:first-order-equations} described the solutions of a single autonomous equation $y'=f(y)$ by its phase line and decided the stability of an equilibrium by the sign of $f'$. This chapter brings these threads together.

The first two sections develop the general theory of linear systems: the structure of the set of solutions (\cref{sec:ode-systems}) and the fundamental matrices, whose determinant satisfies Liouville's formula (\cref{sec:fundamental-matrices}). For constant coefficients, \cref{sec:eigenvalue-method} finds the solutions from the eigenvalues and eigenvectors of $A$, and from its Jordan chains when $A$ is not diagonalisable. \Cref{sec:phase-plane} draws the solutions of systems of two equations as curves in the plane and classifies the possible pictures by the trace and the determinant of $A$. \Cref{sec:linear-stability} decides when all solutions tend to zero, which happens exactly when every eigenvalue has negative real part, and constructs quadratic functions that decrease along the solutions. \Cref{sec:nonhomogeneous-systems} solves the inhomogeneous system by Duhamel's formula and studies its response to bounded and periodic inputs. \Cref{sec:normal-modes} treats the vibrations of mechanical systems $M\mathbf y''+K\mathbf y=\mathbf 0$ through the generalised eigenvalue problem. \Cref{sec:nonlinear-systems} turns to nonlinear systems, which near an equilibrium behave like their linearisation; there we draw the phase portraits of the pendulum, of a predator--prey model and of two competing species. \Cref{sec:systems-numerics} looks at Euler's method for systems, its region of stability and the difficulty of stiff systems.

We keep the conventions of Part~IV (\cref{ch:first-order-equations}): the independent variable is $t$, derivatives are written with primes, and the unknown is a vector-valued function, written in bold. \Cref{ch:canonical-forms,ch:linear-higher-order} called it $\mathbf x$; in this chapter we call it $\mathbf y$, with the components $y_1,\ldots,y_n$, as in the systems $\mathbf y'=\mathbf f(t,\mathbf y)$ of \cref{sec:existence-uniqueness}. Scalars are real or complex, $\mathbb F$ denotes either field, $\|\cdot\|$ is the Euclidean norm of $\mathbb F^n$ and $\|A\|_2$ the norm of a matrix induced by it. Throughout, $I$ denotes an interval with more than one point; the identity matrix is $I_n$, or $I$ when no confusion with the interval is possible.

\section{Systems of First-Order Equations}
\label{sec:ode-systems}

\subsection{Examples}

Systems of differential equations arise whenever the rate of change of one quantity depends on others. A few examples show the variety, and they will return later in the chapter.

Two tanks, each holding one unit of volume of brine, are connected by two pipes (\cref{fig:systems-examples}A). Fresh water flows into the first tank at the rate $q$, the mixture flows from the first tank into the second at the rate $p+q$ and back from the second into the first at the rate $p$, and the mixture leaves the second tank at the rate $q$, so that both volumes stay constant. If $y_1(t)$ and $y_2(t)$ are the amounts of salt in the two tanks and the tanks are stirred so well that the concentrations are uniform, then salt leaves a tank at the rate of the outflow times the concentration, and
\[
y_1'=-(p+q)\,y_1+p\,y_2,\qquad y_2'=(p+q)\,y_1-(p+q)\,y_2 .
\]
For $p=q=1$ this is $\mathbf y'=A\mathbf y$ with
\begin{equation}
A=\begin{bmatrix}-2&1\\2&-2\end{bmatrix}.
\label{eq:tank-matrix}
\end{equation}
Neither equation can be solved alone, because each tank receives salt from the other; this is what makes the equations a system.

Two bodies of mass $m$ slide on a horizontal track, joined to each other and to two walls by three springs with the constant $k$ (\cref{fig:systems-examples}B). If $y_1$ and $y_2$ are the displacements of the bodies from their positions of rest, the left spring is stretched by $y_1$, the middle one by $y_2-y_1$ and the right one by $-y_2$, and Hooke's law and Newton's law give
\[
my_1''=-ky_1+k(y_2-y_1),\qquad my_2''=-k(y_2-y_1)-ky_2 ,
\]
that is, $M\mathbf y''+K\mathbf y=\mathbf 0$ with the mass matrix $M=mI_2$ and the stiffness matrix $K=kT_2$, the chain of masses of \cref{sec:quadratic-forms} with $n=2$. With the velocities $y_3\coloneqq y_1'$ and $y_4\coloneqq y_2'$ it becomes a system of four equations of the first order. Electrical networks with several loops lead to systems in the same way, through Kirchhoff's laws applied to each loop, and the analogy of \cref{tab:analogy} translates every mechanical statement into an electrical one~\cite[Chap.~7]{boyce2017}.

A population of prey $u(t)$, say hares, grows at the rate $au$ in the absence of predators, and a population of predators $w(t)$, say lynxes, dies out at the rate $cw$ in the absence of prey. Encounters, whose number per unit of time is taken proportional to $uw$, reduce the prey and feed the predators:
\begin{equation}
u'=au-buw,\qquad w'=-cw+duw ,
\label{eq:lotka-volterra}
\end{equation}
with positive constants $a$, $b$, $c$ and $d$. These are the \term{Lotka--Volterra equations}~\cite{lotka1925,volterra1926}. They are nonlinear, but the populations $u^*=c/d$ and $w^*=a/b$ make both right sides vanish, and for the deviations $\xi\coloneqq u-u^*$ and $\eta\coloneqq w-w^*$ the equations read $\xi'=-bu^*\eta-b\xi\eta$ and $\eta'=dw^*\xi+d\xi\eta$. When the deviations are small, the products $\xi\eta$ are much smaller still, and dropping them leaves the linear system
\[
\begin{bmatrix}\xi\\\eta\end{bmatrix}'=\begin{bmatrix}0&-bc/d\\ad/b&0\end{bmatrix}\begin{bmatrix}\xi\\\eta\end{bmatrix},
\]
the \term{linearisation} of the system at its equilibrium. \Cref{sec:nonlinear-systems} studies how far it describes the nonlinear system.

\begin{figure}
\sdpanel{A}{\begin{tikzpicture}[x=1cm,y=1cm, font=\scriptsize]
\fill[sdfill] (0,0) rectangle (1.4,1.1); \draw[line width=0.6pt] (0,1.6) -- (0,0) -- (1.4,0) -- (1.4,1.6);
\fill[sdfill] (2.6,0) rectangle (4.0,1.1); \draw[line width=0.6pt] (2.6,1.6) -- (2.6,0) -- (4.0,0) -- (4.0,1.6);
\draw[-{Latex[length=1.6mm]}, line width=0.5pt] (-0.8,1.3) -- (0.2,1.3) node[pos=0, left] {$q$};
\draw[-{Latex[length=1.6mm]}, line width=0.5pt] (1.2,0.3) -- (2.8,0.3) node[midway, below] {$p+q$};
\draw[-{Latex[length=1.6mm]}, line width=0.5pt] (2.8,0.8) -- (1.2,0.8) node[midway, above] {$p$};
\draw[-{Latex[length=1.6mm]}, line width=0.5pt] (3.8,0.3) -- (4.8,0.3) node[pos=1, right] {$q$};
\node at (0.7,0.55) {$y_1$}; \node at (3.3,0.55) {$y_2$};
\end{tikzpicture}}\hfill
\sdpanel{B}{\begin{tikzpicture}[x=1cm,y=1cm, font=\scriptsize]
\fill[pattern=north east lines] (-0.25,0) rectangle (0,1); \draw[line width=0.6pt] (0,0) -- (0,1);
\fill[pattern=north east lines] (5.2,0) rectangle (5.45,1); \draw[line width=0.6pt] (5.2,0) -- (5.2,1);
\draw[line width=0.5pt] (0,0.5) -- (0.2,0.5) -- (0.3,0.7) -- (0.5,0.3) -- (0.7,0.7) -- (0.9,0.3) -- (1.1,0.7) -- (1.2,0.5) -- (1.4,0.5);
\draw[line width=0.6pt, fill=white] (1.4,0.15) rectangle (2.1,0.85); \node at (1.75,0.5) {$m$};
\draw[line width=0.5pt] (2.1,0.5) -- (2.3,0.5) -- (2.4,0.7) -- (2.6,0.3) -- (2.8,0.7) -- (3.0,0.3) -- (3.2,0.7) -- (3.3,0.5) -- (3.5,0.5);
\draw[line width=0.6pt, fill=white] (3.5,0.15) rectangle (4.2,0.85); \node at (3.85,0.5) {$m$};
\draw[line width=0.5pt] (4.2,0.5) -- (4.35,0.5) -- (4.45,0.7) -- (4.6,0.3) -- (4.75,0.7) -- (4.9,0.3) -- (5.0,0.5) -- (5.2,0.5);
\draw[-{Latex[length=1.4mm]}, line width=0.4pt] (1.75,1.05) -- (2.25,1.05) node[right] {$y_1$};
\draw[-{Latex[length=1.4mm]}, line width=0.4pt] (3.85,1.05) -- (4.35,1.05) node[right] {$y_2$};
\node at (0.7,0.05) {$k$}; \node at (2.8,0.05) {$k$}; \node at (4.7,0.05) {$k$};
\end{tikzpicture}}
\caption{Two systems. (A) Two tanks of brine with the amounts of salt $y_1$ and $y_2$; fresh water enters at the rate $q$, and the mixture is pumped between the tanks and leaves at the rates shown. (B) Two equal masses joined to each other and to two walls by three equal springs; $y_1$ and $y_2$ are the displacements from the positions of rest.\label{fig:systems-examples}}
\end{figure}

\subsection{Linear Systems}

\begin{definition}[Linear system]
\label{def:linear-ode-system}
Let $A\colon I\to\mathbb F^{n\times n}$ and $\mathbf b\colon I\to\mathbb F^n$ be continuous. The system
\begin{equation}
\mathbf y'=A(t)\,\mathbf y+\mathbf b(t)
\label{eq:linear-ode-system}
\end{equation}
is a \term{linear system} of $n$ differential equations of the first order, with the \term{coefficient matrix} $A$ and the \term{forcing term} or input $\mathbf b$. It is \term{homogeneous} if $\mathbf b\equiv\mathbf 0$, and it has \term{constant coefficients} if $A$ does not depend on $t$. A \term{solution} on $I$ is a continuously differentiable function $\mathbf y\colon I\to\mathbb F^n$ that satisfies \eqref{eq:linear-ode-system} at every $t\in I$.
\end{definition}

Written out, the $i$th equation is $y_i'=\sum_ja_{ij}(t)y_j+b_i(t)$: the rate of change of each component is a linear combination of all components, plus an input. Vector- and matrix-valued functions are differentiated and integrated entry by entry, as in \cref{sec:matrix-exponential}, and complex-valued entries are handled by \cref{lem:complex-calculus}. The product rule $(F\mathbf y)'=F'\mathbf y+F\mathbf y'$ holds for a differentiable matrix-valued $F$, with the order of the factors kept.

Every linear equation of order $k$ is equivalent to a linear system: by \cref{prop:companion-system}, $y$ solves $y^{(k)}+a_{k-1}(t)y^{(k-1)}+\cdots+a_0(t)y=g(t)$ if and only if $(y,y',\ldots,y^{(k-1)})$ solves the companion system with the matrix \eqref{eq:companion-system}. The system of the two masses is of this kind, an equation of the second order for a vector, and becomes a system of the first order in the same way. The converse is not true: not every system is equivalent to a single equation of order $n$ for one of its components. For constant coefficients, a system $\mathbf y'=A\mathbf y$ can be transformed by a constant change of variables $\mathbf y=S\mathbf z$ into a companion system $\mathbf z'=C\mathbf z$ only if $A$ is similar to a companion matrix, and by \cref{prop:nonderogatory} this happens exactly when every eigenvalue of $A$ has a single Jordan block. The system $\mathbf y'=\lambda I_2\mathbf y$, which consists of two unrelated copies of $y'=\lambda y$, is the simplest system that is not of this form. Systems are therefore the more general object, and the theory of this chapter contains that of \cref{ch:linear-higher-order} as a special case.

\subsection{Existence and Uniqueness}

The basic facts about linear systems were proved in \cref{ch:linear-higher-order}, and we recall them. By \cref{thm:linear-existence}, for every $t_0\in I$ and $\mathbf y_0\in\mathbb F^n$ the system \eqref{eq:linear-ode-system} has exactly one solution on $I$ with $\mathbf y(t_0)=\mathbf y_0$, and it is real if $A$, $\mathbf b$ and $\mathbf y_0$ are real. The solution exists on the whole interval $I$, however long: a linear system does not blow up. The proof used the norm $\|\cdot\|_\infty$; for estimates of the size of solutions the Euclidean norm is more natural, and we need the triangle inequality for integrals in this norm, which \cref{prop:mvt-integral-rn} proved for real vectors.

\begin{lemma}[Norms of integrals]
\label{lem:integral-norm}
Let $a<b$.
\begin{enumerate}
\item If $\mathbf v\colon[a,b]\to\mathbb F^n$ is continuous, then $\bigl\|\int_a^b\mathbf v\bigr\|\le\int_a^b\|\mathbf v\|$.
\item If $M\colon[a,b]\to\mathbb F^{n\times n}$ is continuous, then $t\mapsto\|M(t)\|_2$ is continuous and $\bigl\|\int_a^bM\bigr\|_2\le\int_a^b\|M\|_2$.
\end{enumerate}
\end{lemma}

\begin{proof}
(i) The function $t\mapsto\|\mathbf v(t)\|$ is continuous, as a composition of continuous functions. Let $\mathbf c\coloneqq\int_a^b\mathbf v$; we may assume $\mathbf c\ne\mathbf 0$. By linearity of the integral, applied entry by entry, $\|\mathbf c\|^2=\mathbf c^{\mathsf H}\mathbf c=\int_a^b\mathbf c^{\mathsf H}\mathbf v(t)\,dt$. This number is real, so it equals the integral of the real part of $\mathbf c^{\mathsf H}\mathbf v(t)$, which is at most $|\mathbf c^{\mathsf H}\mathbf v(t)|\le\|\mathbf c\|\,\|\mathbf v(t)\|$ by the Cauchy--Schwarz inequality (\cref{thm:cauchy-schwarz-general}). By the monotonicity of the integral (\cref{thm:integral-monotone}), $\|\mathbf c\|^2\le\|\mathbf c\|\int_a^b\|\mathbf v\|$, and dividing by $\|\mathbf c\|$ gives the claim. (ii) Since $\bigl|\|M(t)\|_2-\|M(s)\|_2\bigr|\le\|M(t)-M(s)\|_2\le\|M(t)-M(s)\|_{\mathrm F}$ (\cref{prop:matrix-norms}), and the Frobenius norm of $M(t)-M(s)$ tends to zero as $s\to t$, the norm is continuous. For a unit vector $\mathbf x$, the vector $\bigl(\int_a^bM\bigr)\mathbf x$ is $\int_a^bM(t)\mathbf x\,dt$, entry by entry, so by (i) its norm is at most $\int_a^b\|M(t)\mathbf x\|\,dt\le\int_a^b\|M\|_2$.
\end{proof}

For nonlinear systems $\mathbf y'=\mathbf f(t,\mathbf y)$ we use the results of \cref{sec:existence-uniqueness} in the form described at its end: with the Euclidean norm in place of absolute values, they hold for systems as stated there~\cite[Chap.~II]{hartman2002}. In particular, if $\mathbf f$ and its Jacobian matrix with respect to $\mathbf y$ are continuous on an open set $D\subseteq\mathbb R\times\mathbb R^n$, then every initial value problem has a unique maximal solution, two solutions that agree at one time agree on their common interval, and a maximal solution leaves every compact subset of $D$ as $t$ approaches an end of its interval (\cref{cor:local-uniqueness,thm:maximal-solution}). Volume~II proves these theorems in general.

\subsection{The Structure of the Solution Set}

For a homogeneous system, the solutions form a vector space whose dimension is the number of equations, and the initial value at any time is a complete set of coordinates on it.

\begin{theorem}[Solution space of a linear system]
\label{thm:system-solution-space}
Let $A\colon I\to\mathbb F^{n\times n}$ be continuous, and let $V$ be the set of solutions of $\mathbf y'=A(t)\mathbf y$ on $I$.
\begin{enumerate}
\item $V$ is a vector space over $\mathbb F$, and for every $t_0\in I$ the map $E_{t_0}\colon V\to\mathbb F^n$, $\mathbf y\mapsto\mathbf y(t_0)$, is an isomorphism. In particular $\dim V=n$.
\item If $A$ is real, the real solutions form a real vector space of dimension $n$, and every basis of it is a basis of the space of complex solutions.
\item If $\mathbf b\colon I\to\mathbb F^n$ is continuous and $\mathbf y_{\mathrm p}$ is one solution of $\mathbf y'=A(t)\mathbf y+\mathbf b(t)$, then the solutions of this system are exactly the functions $\mathbf y_{\mathrm p}+\mathbf y_{\mathrm h}$ with $\mathbf y_{\mathrm h}\in V$.
\end{enumerate}
\end{theorem}

\begin{proof}
(i) The map $\mathbf y\mapsto\mathbf y'-A\mathbf y$ from continuously differentiable functions to continuous ones is linear, and $V$ is its kernel, hence a vector space (\cref{prop:kernel-image}). The evaluation $E_{t_0}$ is linear. It is surjective, since by \cref{thm:linear-existence} every vector $\mathbf y_0$ is the value at $t_0$ of a solution, and injective, since a solution with $\mathbf y(t_0)=\mathbf 0$ has the same initial value as the zero solution and is therefore zero by the uniqueness part of the same theorem. So $E_{t_0}$ is an isomorphism, and $\dim V=n$ (\cref{thm:isomorphic-dimension}). (ii) If $A$ is real, $E_{t_0}$ maps the real solutions injectively onto $\mathbb R^n$, by the last statement of \cref{thm:linear-existence}, so they form a real space of dimension $n$. A basis of it is mapped to $n$ real vectors that are independent over $\mathbb R$; the real matrix formed by them has a nonzero determinant (\cref{thm:det-invertible}), so they are independent over $\mathbb C$ as well, and since $E_{t_0}$ is an isomorphism of the complex spaces, their preimages form a basis of $V$. (iii) This is \cref{prop:preimage-flat} for the linear map $\mathbf y\mapsto\mathbf y'-A\mathbf y$, whose preimage of $\mathbf b$ is the set of solutions of the inhomogeneous system.
\end{proof}

The theorem is the system form of \cref{thm:solution-space}, and its proof is the same; for the companion system of an equation of order $k$ it is that theorem, since the solutions of the system and of the equation correspond one to one. Part (iii) is the principle that \cref{ch:linear-systems,ch:linear-maps} established for linear equations of every kind: the general solution of an inhomogeneous linear problem is one particular solution plus the general solution of the homogeneous problem. For the tank system \eqref{eq:tank-matrix} the theorem says that the solutions form a plane in the space of all pairs of functions, and that the two initial amounts of salt determine the solution; for the two masses, which form a system of four equations, the two initial positions and the two initial velocities do.

\section{Fundamental Matrices}
\label{sec:fundamental-matrices}

\subsection{Fundamental Matrices and the Transition Matrix}

A basis of the solution space can be stored as the columns of a matrix, and the whole theory of the homogeneous system can then be written as statements about one matrix-valued function.

\begin{definition}[Fundamental matrix]
\label{def:fundamental-matrix}
Let $A\colon I\to\mathbb F^{n\times n}$ be continuous. A \term{matrix solution} of $\mathbf y'=A(t)\mathbf y$ is a continuously differentiable $Y\colon I\to\mathbb F^{n\times n}$ with $Y'=A(t)Y$ on $I$, that is, a matrix whose columns are solutions. A \term{fundamental matrix} is a matrix solution whose columns form a basis of the solution space, a \term{fundamental set of solutions}.
\end{definition}

\begin{proposition}[Fundamental matrices]
\label{prop:fundamental-matrices}
Let $Y$ be a matrix solution of $\mathbf y'=A(t)\mathbf y$ on $I$.
\begin{enumerate}
\item The following are equivalent: $Y$ is a fundamental matrix; $Y(t_0)$ is invertible for some $t_0\in I$; $Y(t)$ is invertible for every $t\in I$.
\item If $Y$ is a fundamental matrix, the solutions are exactly the functions $Y(t)\mathbf c$ with $\mathbf c\in\mathbb F^n$, and the solution with $\mathbf y(t_0)=\mathbf y_0$ has $\mathbf c=Y(t_0)^{-1}\mathbf y_0$.
\item If $Y$ is a fundamental matrix, the fundamental matrices are exactly the matrices $YC$ with an invertible constant $C\in\mathbb F^{n\times n}$.
\end{enumerate}
\end{proposition}

\begin{proof}
(i) The columns of $Y$ are solutions, and by \cref{thm:system-solution-space}\,(i) they form a basis of the solution space if and only if their values at any one time $t$ form a basis of $\mathbb F^n$, that is, if and only if $Y(t)$ is invertible; the isomorphism $E_t$ exists for every $t$. (ii) $Y(t)\mathbf c=\sum_jc_j\mathbf y_j(t)$ is a combination of the columns, and every solution is one, since the columns are a basis. The value at $t_0$ is $Y(t_0)\mathbf c$. (iii) $(YC)'=Y'C=AYC$, so $YC$ is a matrix solution, and it is invertible at every $t$ if $C$ is. Conversely, if $Z$ is a fundamental matrix, then by (ii) each column of $Z$ is $Y\mathbf c_j$ for a constant vector $\mathbf c_j$, so $Z=YC$ with $C\coloneqq Y(t_0)^{-1}Z(t_0)$, which is invertible.
\end{proof}

Part (i) generalises \cref{thm:wronskian-solutions}: the determinant of a matrix solution vanishes everywhere or nowhere. For the companion system of an equation of order $k$, the matrix $Y$ of \eqref{eq:wronskian-matrix} formed from $k$ solutions and their derivatives is a matrix solution, and its determinant is their Wronskian. Among all fundamental matrices one is distinguished at each time $s$: the one that equals the identity there.

\begin{definition}[Transition matrix]
\label{def:transition-matrix}
For $t,s\in I$, the \term{transition matrix} $\Phi(t,s)$ of $\mathbf y'=A(t)\mathbf y$ is the value at $t$ of the matrix solution $Y$ with $Y(s)=I_n$.
\end{definition}

The matrix solution with $Y(s)=I_n$ exists and is unique, by \cref{thm:linear-existence} applied to each column, whose initial value is a column of the identity, and it is a fundamental matrix by \cref{prop:fundamental-matrices}\,(i). Its $j$th column is the solution that starts at $\mathbf e_j$ at the time $s$.

\begin{proposition}[Properties of the transition matrix]
\label{prop:transition-matrix}
Let $A\colon I\to\mathbb F^{n\times n}$ be continuous, and let $Y$ be any fundamental matrix of $\mathbf y'=A(t)\mathbf y$. For all $t,s,r\in I$:
\begin{enumerate}
\item $\Phi(t,s)=Y(t)Y(s)^{-1}$, and the solution with $\mathbf y(s)=\mathbf y_0$ is $\mathbf y(t)=\Phi(t,s)\mathbf y_0$;
\item $\Phi(t,s)\Phi(s,r)=\Phi(t,r)$ and $\Phi(t,t)=I_n$;
\item $\Phi(t,s)$ is invertible, with $\Phi(t,s)^{-1}=\Phi(s,t)$;
\item if $A$ is constant, $\Phi(t,s)=e^{(t-s)A}$.
\end{enumerate}
\end{proposition}

\begin{proof}
(i) $t\mapsto Y(t)Y(s)^{-1}$ is a matrix solution, by \cref{prop:fundamental-matrices}\,(iii), with the value $I_n$ at $t=s$, so it is $\Phi(\cdot,s)$ by uniqueness. The function $\Phi(t,s)\mathbf y_0$ is a combination of the columns, hence a solution, with the value $\mathbf y_0$ at $s$. (ii) By (i), $\Phi(t,s)\Phi(s,r)=Y(t)Y(s)^{-1}Y(s)Y(r)^{-1}=Y(t)Y(r)^{-1}=\Phi(t,r)$, and $\Phi(t,t)=Y(t)Y(t)^{-1}=I_n$. (iii) follows from (ii) with $r=t$. (iv) By \cref{thm:exp-derivative} and the chain rule, $t\mapsto e^{(t-s)A}$ is a matrix solution, and its value at $t=s$ is $e^0=I_n$ (\cref{prop:exp-basic}).
\end{proof}

The transition matrix carries the state of the system from the time $s$ to the time $t$: whatever the state $\mathbf y(s)$ is, the state at the time $t$ is $\Phi(t,s)\mathbf y(s)$, a linear function of it. Part (ii) says that evolving from $r$ to $s$ and then from $s$ to $t$ is the same as evolving from $r$ to $t$ directly, and part (iii) that the evolution can be undone: the state at an earlier time is determined by the state at a later one. For constant coefficients these are the exponential laws $e^{(t-s)A}e^{(s-r)A}=e^{(t-r)A}$ and $(e^{tA})^{-1}=e^{-tA}$ of \cref{cor:exp-laws}, and the transition matrix depends only on the elapsed time $t-s$: a system with constant coefficients behaves in the same way whenever it is started.

\subsection{Liouville's Formula}

The determinant of a matrix solution satisfies a linear equation of the first order, which generalises Abel's formula.

\begin{theorem}[Liouville's formula]
\label{thm:liouville-formula}
Let $Y$ be a matrix solution of $\mathbf y'=A(t)\mathbf y$ on $I$, with $A$ continuous, and let $t_0\in I$. Then $(\det Y)'=\operatorname{tr}A(t)\cdot\det Y$ on $I$, and
\[
\det Y(t)=\det Y(t_0)\exp\Bigl(\int_{t_0}^t\operatorname{tr}A(s)\,ds\Bigr)\qquad(t\in I).
\]
\end{theorem}

\begin{proof}
By \cref{thm:jacobi-formula}, $(\det Y)'=\operatorname{tr}(\operatorname{adj}Y\cdot Y')$; the theorem was stated for real entries, but its proof uses only the product rule, which holds for complex-valued functions (\cref{lem:complex-calculus}), as in the proof of \cref{thm:abel}. Since $Y'=AY$, and the trace of a product does not change when the factors are exchanged,
\[
(\det Y)'=\operatorname{tr}\bigl(\operatorname{adj}Y\cdot AY\bigr)=\operatorname{tr}\bigl(A\,Y\operatorname{adj}Y\bigr)=\operatorname{tr}\bigl(A\cdot\det Y\,I_n\bigr)=\operatorname{tr}A\cdot\det Y ,
\]
by \cref{thm:adjugate}. With $P(t)\coloneqq\int_{t_0}^t\operatorname{tr}A(s)\,ds$, the function $e^{-P}\det Y$ has the derivative $e^{-P}\bigl((\det Y)'-\operatorname{tr}A\cdot\det Y\bigr)=0$, computed as in the proof of \cref{thm:abel}, so it is constant and equal to its value $\det Y(t_0)$ at $t_0$ (\cref{lem:complex-calculus}).
\end{proof}

The proof via the adjugate needs no invertibility, so the theorem applies to every matrix solution, and since the exponential has no zero, it gives a second proof of \cref{prop:fundamental-matrices}\,(i): the determinant is zero everywhere or nowhere. For the companion system of $y^{(k)}+a_{k-1}y^{(k-1)}+\cdots+a_0y=0$ the trace of \eqref{eq:companion-system} is $-a_{k-1}$, and the formula is Abel's formula (\cref{thm:abel}) for the Wronskian. For a constant $A$ and $Y(t)=e^{tA}$, it is the identity $\det e^{tA}=e^{t\operatorname{tr}A}$ of \cref{prop:exp-basic}\,(iii).

The formula has a geometric meaning. Let $A$ be real. By \cref{thm:det-volume}, $|\det\Phi(t,s)|$ is the factor by which the linear map $\Phi(t,s)$ changes volumes: a region $R$ of initial values at the time $s$, a cloud of possible initial states, is carried by the solutions into the region $\Phi(t,s)R$ at the time $t$, whose volume is $\exp(\int_s^t\operatorname{tr}A)$ times that of $R$. The trace of $A$ is the divergence of the vector field $\mathbf y\mapsto A\mathbf y$ (\cref{sec:div-curl}), and the formula says that the volume of a cloud of states grows or shrinks at the rate of the divergence. For the tanks \eqref{eq:tank-matrix}, $\operatorname{tr}A=-4$, and every region of initial amounts of salt shrinks in area like $e^{-4t}$; for the harmonic oscillator $y''+\omega^2y=0$, written as a system, the trace is zero and areas in the phase plane are preserved, as they are by every motion of a mechanical system without friction, according to a theorem of Liouville in mechanics~\cite[Sect.~16]{arnold1989}.

\subsection{Variable Coefficients}

For constant coefficients, $\Phi(t,s)=e^{(t-s)A}$. It is natural to guess that for variable coefficients the transition matrix is the exponential of $\int_s^tA$, by analogy with the scalar equation $y'=a(t)y$, whose solutions are $\exp(\int_s^ta)\,y(s)$ (\cref{thm:linear-homogeneous}). The guess is right when the values of $A$ commute with each other, and wrong in general.

\begin{proposition}[Commuting coefficients]
\label{prop:commuting-coefficients}
Let $A\colon I\to\mathbb F^{n\times n}$ be continuous, and suppose that $A(t)A(r)=A(r)A(t)$ for all $t,r\in I$. Then
\[
\Phi(t,s)=\exp\Bigl(\int_s^tA(r)\,dr\Bigr)\qquad(t,s\in I).
\]
\end{proposition}

\begin{proof}
Fix $s$ and let $B(t)\coloneqq\int_s^tA(r)\,dr$. Multiplication by a constant matrix commutes with integration, entry by entry, so $A(t)B(t)=\int_s^tA(t)A(r)\,dr=\int_s^tA(r)A(t)\,dr=B(t)A(t)$, and in the same way $B(t)$ and $B(t+h)$ commute, both being integrals of the commuting family $A(r)$. Let $\Delta\coloneqq B(t+h)-B(t)=\int_t^{t+h}A$, which commutes with $B(t)$. By \cref{cor:exp-laws}\,(ii),
\[
\frac{e^{B(t+h)}-e^{B(t)}}h=e^{B(t)}\,\frac{e^\Delta-I}h=e^{B(t)}\Bigl(\frac\Delta h+\frac{e^\Delta-I-\Delta}h\Bigr).
\]
As $h\to0$, $\Delta/h\to A(t)$ entry by entry, by the fundamental theorem of calculus (\cref{thm:ftc1}), applied to the real and imaginary parts of each entry. In particular $\|\Delta\|_2\le\|\Delta\|_{\mathrm F}\le C|h|$ for small $h$, with a constant $C$, and the series of the exponential gives $\|e^\Delta-I-\Delta\|_2\le\sum_{k\ge2}\|\Delta\|_2^k/k!\le\|\Delta\|_2^2e^{\|\Delta\|_2}\le C^2h^2e^{C|h|}$, so the last term tends to zero. Hence $t\mapsto e^{B(t)}$ is differentiable with the derivative $e^{B(t)}A(t)=A(t)e^{B(t)}$, since $A(t)$ commutes with $B(t)$ and hence with every power of it. It is a matrix solution with the value $I_n$ at $t=s$, so it is $\Phi(\cdot,s)$.
\end{proof}

The hypothesis holds for $1\times1$ systems, for diagonal $A(t)$, and for $A(t)=a(t)A_0$ with a scalar function $a$ and a constant matrix $A_0$; in the last case $\Phi(t,s)=\exp\bigl((\int_s^ta)\,A_0\bigr)$, a change of the time scale in the system $\mathbf y'=A_0\mathbf y$. Without the hypothesis the formula fails. For
\[
A(t)=\begin{bmatrix}1&t\\0&0\end{bmatrix},\qquad A(t)A(r)=\begin{bmatrix}1&r\\0&0\end{bmatrix}\ne\begin{bmatrix}1&t\\0&0\end{bmatrix}=A(r)A(t)\quad(t\ne r),
\]
the solution with $\mathbf y(0)=\mathbf e_2$ has $y_2\equiv1$ and $y_1'=y_1+t$, $y_1(0)=0$, so $y_1(t)=e^t-1-t$, by \cref{thm:linear-first-order}. On the other hand,
\[
M\coloneqq\int_0^tA(r)\,dr=\begin{bmatrix}t&t^2/2\\0&0\end{bmatrix}
\]
satisfies $M^2=tM$, so $M^k=t^{k-1}M$ for $k\ge1$ and $e^M=I+((e^t-1)/t)M$ for $t\ne0$, whose second column is $(t(e^t-1)/2,1)$. At $t=1$ the true solution has $y_1=e-2=0.718$, the exponential formula gives $(e-1)/2=0.859$. For systems with variable coefficients there is in general no formula for the solutions in terms of integrals of the coefficients, just as there is none for equations of the second order (\cref{sec:variable-coefficients}); \cref{ch:series-solutions} finds solutions as power series when the coefficients are themselves power series.

\begin{note}{Periodic coefficients and Floquet theory}
If $A(t+T)=A(t)$ for all $t$, the transition matrix satisfies $\Phi(t+T,s+T)=\Phi(t,s)$, since both are matrix solutions in $t$ with the value $I_n$ at the same point, and so $\Phi(kT,0)=\Phi(T,0)^k$ by \cref{prop:transition-matrix}\,(ii). The solutions at the times $kT$ are therefore the iterates of the \emph{monodromy matrix} $\Phi(T,0)$. Since $\Phi(t,0)=\Phi(t-kT,0)\,\Phi(T,0)^k$ for $kT\le t<(k+1)T$, and $\Phi(\cdot,0)$ is bounded on $[0,T]$, \cref{thm:powers-size} shows that all solutions tend to zero if and only if every eigenvalue of $\Phi(T,0)$, every \emph{Floquet multiplier}, has modulus less than $1$. Floquet showed in 1883 that $\Phi(t,0)=P(t)e^{tB}$ with a $T$-periodic matrix $P(t)$ and a constant matrix $B$, so that a periodic system becomes one with constant coefficients after a periodic change of variables~\cite{floquet1883}; Hartman~\cite[Chap.~IV]{hartman2002} gives the proof. The eigenvalues of the matrices $A(t)$ themselves do not decide stability, even when they do not depend on $t$. For the $\pi$-periodic matrix
\[
A(t)=\begin{bmatrix}-1+\frac32\cos^2t&1-\frac32\sin t\cos t\\[2pt]-1-\frac32\sin t\cos t&-1+\frac32\sin^2t\end{bmatrix},
\]
the trace is $-1/2$ and the determinant $1/2$ for every $t$, so the eigenvalues $(-1\pm i\sqrt7)/4$ have negative real parts for every $t$; yet the function $e^{t/2}(-\cos t,\sin t)$, which grows without bound, is a solution, as one verifies by differentiation. The example is due to Markus and Yamabe~\cite[p.~310]{markus1960}.
\end{note}

\section{The Eigenvalue Method}
\label{sec:eigenvalue-method}

For a system with constant coefficients, $\mathbf y'=A\mathbf y$, the solutions are known: $\mathbf y(t)=e^{tA}\mathbf y(0)$ by \cref{thm:linear-system}. But the formula hides their shape, and the series $\sum_kt^kA^k/k!$ is not a practical way to see it. The shape is revealed by the eigenvalues and eigenvectors of $A$, as the shape of the powers $A^k\mathbf x_0$ was in \cref{sec:matrix-powers}. This section develops the method for diagonalisable matrices, for real matrices with complex eigenvalues and for defective matrices, and collects the ways of computing $e^{tA}$ by hand.

\subsection{Exponential Solutions}

For the scalar equation $p(D)y=0$ of \cref{sec:constant-coefficients}, the trial function $e^{\lambda t}$ led to the characteristic equation. For a system, the natural trial is an exponential times a constant vector.

\begin{proposition}[Exponential solutions]
\label{prop:exponential-solutions}
Let $A\in\mathbb F^{n\times n}$, $\lambda\in\mathbb C$ and $\mathbf v\in\mathbb C^n$ with $\mathbf v\ne\mathbf 0$. The function $\mathbf y(t)=e^{\lambda t}\mathbf v$ is a solution of $\mathbf y'=A\mathbf y$ if and only if $A\mathbf v=\lambda\mathbf v$.
\end{proposition}

\begin{proof}
By \cref{lem:complex-calculus}\,(ii), applied to each entry, $\mathbf y'(t)=\lambda e^{\lambda t}\mathbf v$, while $A\mathbf y(t)=e^{\lambda t}A\mathbf v$. Since $e^{\lambda t}\ne0$, the two agree at some, and then at every, $t$ if and only if $\lambda\mathbf v=A\mathbf v$.
\end{proof}

An eigenvector thus spans a line of initial values along which the system acts like the scalar equation $y'=\lambda y$. For a real eigenvalue and a real eigenvector, a solution that starts on the line $\{c\mathbf v:c\in\mathbb R\}$ stays on it for all time, and moves along it towards the origin if $\lambda<0$ and away from it if $\lambda>0$; these \term{straight-line solutions} are the skeleton of the pictures in \cref{sec:phase-plane}. When there are enough eigenvectors, the straight-line solutions give everything.

\begin{theorem}[Diagonalisable systems]
\label{thm:diagonalisable-systems}
Let $A\in\mathbb F^{n\times n}$ be diagonalisable, $A=S\Lambda S^{-1}$, with the eigenvectors $\mathbf v_1,\ldots,\mathbf v_n$ as the columns of $S$ and the eigenvalues $\Lambda=\operatorname{diag}(\lambda_1,\ldots,\lambda_n)$. Then the functions $e^{\lambda_1t}\mathbf v_1,\ldots,e^{\lambda_nt}\mathbf v_n$ form a fundamental set of solutions of $\mathbf y'=A\mathbf y$, and
\[
e^{tA}=Se^{t\Lambda}S^{-1}=S\operatorname{diag}\bigl(e^{\lambda_1t},\ldots,e^{\lambda_nt}\bigr)S^{-1}.
\]
The solution with $\mathbf y(0)=\mathbf y_0$ is $\mathbf y(t)=\sum_ic_ie^{\lambda_it}\mathbf v_i$ with $(c_1,\ldots,c_n)=S^{-1}\mathbf y_0$.
\end{theorem}

\begin{proof}
Each function is a solution by \cref{prop:exponential-solutions}, so $Y(t)\coloneqq Se^{t\Lambda}$, whose columns they are, is a matrix solution; $Y(0)=S$ is invertible, so $Y$ is a fundamental matrix (\cref{prop:fundamental-matrices}). By \cref{prop:transition-matrix}, $e^{tA}=\Phi(t,0)=Y(t)Y(0)^{-1}=Se^{t\Lambda}S^{-1}$. The last statement is $\mathbf y(t)=e^{tA}\mathbf y_0=Se^{t\Lambda}\mathbf c$ with $\mathbf c=S^{-1}\mathbf y_0$.
\end{proof}

The formula for $e^{tA}$ is also \cref{prop:exp-basic}\,(ii) with $X=S$. The theorem has a simple interpretation. In the coordinates $\mathbf z\coloneqq S^{-1}\mathbf y$ of the eigenvector basis, the system becomes $\mathbf z'=S^{-1}AS\mathbf z=\Lambda\mathbf z$, that is, $z_i'=\lambda_iz_i$ for each $i$: the change of variables \term{decouples} the system into $n$ scalar equations, each of which describes one \term{mode} that evolves on its own, multiplied by $e^{\lambda_it}$. This is the continuous counterpart of the expansion \eqref{eq:power-expansion} of the powers $A^k\mathbf x_0=\sum_ic_i\lambda_i^k\mathbf x_i$, with $e^{\lambda_it}$ in place of $\lambda_i^k$; the factor $e^{\lambda t}$ for the time $t$ is the $t$th power of $e^\lambda$, the factor of one unit of time.

For the tanks \eqref{eq:tank-matrix}, the characteristic polynomial $(\lambda+2)^2-2$ has the roots $\lambda_{1,2}=-2\pm\sqrt2$, and the first row of $A-\lambda I$, which is $(-2-\lambda,1)$, gives the eigenvectors $\mathbf v_{1,2}=(1,\pm\sqrt2)$. If at first all the salt, one unit, is in the first tank, then $\mathbf y_0=(1,0)=(\mathbf v_1+\mathbf v_2)/2$, and
\[
\mathbf y(t)=\frac12e^{(-2+\sqrt2)t}\begin{bmatrix}1\\\sqrt2\end{bmatrix}+\frac12e^{(-2-\sqrt2)t}\begin{bmatrix}1\\-\sqrt2\end{bmatrix}=e^{-2t}\begin{bmatrix}\cosh\sqrt2\,t\\\sqrt2\sinh\sqrt2\,t\end{bmatrix}.
\]
The salt is washed out, and both amounts decay. The two modes decay at very different rates, $2-\sqrt2=0.586$ and $2+\sqrt2=3.414$. The fast mode, with the eigenvector $(1,-\sqrt2)$ of mixed signs, describes the redistribution of salt between the tanks, and it has practically disappeared after one unit of time, when $e^{-3.414}=0.033$; afterwards the slow mode dominates, and the amounts keep the fixed ratio $y_2/y_1=\sqrt2\tanh\sqrt2\,t\to\sqrt2$ while they decay together like $e^{-0.586t}$.

The theorem also settles the question left open in \cref{ch:linear-higher-order}: a repeated eigenvalue does not by itself produce factors $t^j$. For $A=\lambda I_n$ every vector is an eigenvector, and the solutions are $e^{\lambda t}\mathbf y_0$. For the symmetric matrix
\[
A=\begin{bmatrix}1&2&2\\2&1&2\\2&2&1\end{bmatrix},
\]
the eigenvalue $5$ has the eigenvector $(1,1,1)$, and the eigenvalue $-1$ has the two-dimensional eigenspace of the vectors with $y_1+y_2+y_3=0$; $A$ is diagonalisable by the spectral theorem (\cref{thm:spectral}), and every solution is $e^{5t}c(1,1,1)+e^{-t}\mathbf w$ with a vector $\mathbf w$ whose components add up to zero. The double eigenvalue $-1$ contributes no solution $te^{-t}$. The factors $t^j$ of \cref{thm:constant-coefficients} arose because the companion matrix of an equation has a single Jordan block for each eigenvalue (\cref{prop:nonderogatory}); they come from Jordan blocks, not from repeated eigenvalues.

\subsection{Complex Eigenvalues of Real Matrices}

A real matrix may have nonreal eigenvalues, and then the exponential solutions are complex, although the system and the solutions of interest are real. As for equations with real coefficients (\cref{prop:linear-structure}\,(iv)), the real and imaginary parts of a complex solution are real solutions.

\begin{proposition}[Real solutions from a complex eigenvalue]
\label{prop:real-solutions}
Let $A\in\mathbb R^{n\times n}$.
\begin{enumerate}
\item If $\mathbf y$ is a complex solution of $\mathbf y'=A\mathbf y$, then $\operatorname{Re}\mathbf y$ and $\operatorname{Im}\mathbf y$ are real solutions.
\item Let $\lambda=a+ib$ with $b\ne0$ be an eigenvalue of $A$ with the eigenvector $\mathbf w=\mathbf p+i\mathbf q$, $\mathbf p,\mathbf q\in\mathbb R^n$. Then
\begin{equation}
\mathbf y_1(t)=e^{at}\bigl(\cos(bt)\,\mathbf p-\sin(bt)\,\mathbf q\bigr),\qquad\mathbf y_2(t)=e^{at}\bigl(\sin(bt)\,\mathbf p+\cos(bt)\,\mathbf q\bigr)
\label{eq:real-complex-solutions}
\end{equation}
are real solutions, and they are linearly independent.
\end{enumerate}
\end{proposition}

\begin{proof}
(i) Differentiation acts on real and imaginary parts separately, and multiplication by the real matrix $A$ preserves them, so $\mathbf y'=A\mathbf y$ splits into $(\operatorname{Re}\mathbf y)'=A\operatorname{Re}\mathbf y$ and $(\operatorname{Im}\mathbf y)'=A\operatorname{Im}\mathbf y$. (ii) By Euler's formula (\cref{cor:euler-consequences}),
\[
e^{\lambda t}\mathbf w=e^{at}(\cos bt+i\sin bt)(\mathbf p+i\mathbf q)=\mathbf y_1(t)+i\,\mathbf y_2(t),
\]
which is a solution by \cref{prop:exponential-solutions}, so $\mathbf y_1$ and $\mathbf y_2$ are solutions by (i). Their values at $t=0$ are $\mathbf p$ and $\mathbf q$. Conjugating $A\mathbf w=\lambda\mathbf w$ shows that $\bar{\mathbf w}=\mathbf p+i(-\mathbf q)$ is an eigenvector for $\bar\lambda=a-ib$, and by \cref{prop:real-2x2-form} its real and imaginary parts $\mathbf p$ and $-\mathbf q$ are linearly independent. So the solutions are independent, by \cref{thm:system-solution-space}\,(i).
\end{proof}

In the plane spanned by $\mathbf p$ and $\mathbf q$, the solutions \eqref{eq:real-complex-solutions} rotate and change their length: in the coordinates of the basis $\mathbf p,\mathbf q$, the solution $\mathbf y_1$ has the coordinates $e^{at}(\cos bt,-\sin bt)$, a point that turns about the origin with the angular velocity $b$ while its distance from the origin changes by the factor $e^{at}$. The real part $a$ of the eigenvalue is the rate of growth or decay, the imaginary part $b$ the angular frequency, exactly as for the real solutions $e^{\gamma t}\cos\beta t$ of a single equation (\cref{cor:real-fundamental}).

\begin{corollary}[Real fundamental set]
\label{cor:real-fundamental-system}
Let $A\in\mathbb R^{n\times n}$ be diagonalisable over $\mathbb C$, with the real eigenvalues $\lambda_1,\ldots,\lambda_r$ and the pairs of nonreal eigenvalues $a_j\pm ib_j$, $b_j>0$, $j=1,\ldots,s$, counted with multiplicity. Choose real eigenvectors $\mathbf v_1,\ldots,\mathbf v_r$ and eigenvectors $\mathbf w_j=\mathbf p_j+i\mathbf q_j$ for $a_j+ib_j$ such that $\mathbf v_1,\ldots,\mathbf v_r,\bar{\mathbf w}_1,\ldots,\bar{\mathbf w}_s$ together with the vectors $\mathbf w_j$ form a basis of $\mathbb C^n$. Then the $n$ real functions $e^{\lambda_it}\mathbf v_i$ and the pairs \eqref{eq:real-complex-solutions} formed from the $\mathbf w_j$ are a fundamental set of solutions, real and complex.
\end{corollary}

\begin{proof}
All are real solutions, by \cref{prop:exponential-solutions,prop:real-solutions}. Their values at $t=0$ are the vectors $\mathbf v_i$, $\mathbf p_j$ and $\mathbf q_j$. Since $\mathbf w_j=\mathbf p_j+i\mathbf q_j$ and $\bar{\mathbf w}_j=\mathbf p_j-i\mathbf q_j$, the vectors $\mathbf p_j$ and $\mathbf q_j$ span the same complex space as $\mathbf w_j$ and $\bar{\mathbf w}_j$, so the $n$ values span $\mathbb C^n$ and are a basis of it, and the solutions form a basis by \cref{thm:system-solution-space}. As in the proof of \cref{thm:real-block-diagonal}, the values are $n$ real vectors independent over $\mathbb C$, hence over $\mathbb R$, and the solutions are also a basis of the real solutions.
\end{proof}

Such eigenvectors exist: by \cref{thm:diagonalisable,thm:eigenvalues-exist}, a basis of each eigenspace $E_{\bar\mu}$ conjugates to a basis of $E_\mu$, as in the proof of \cref{thm:real-block-diagonal}. In matrix form, the corollary is the real block diagonal form of that theorem: if $S^{-1}AS=D$ is block diagonal with the diagonal entries $\lambda_i$ and the blocks $C_{a_j,b_j}$, then $e^{tA}=Se^{tD}S^{-1}$, and $e^{tD}$ has the diagonal entries $e^{\lambda_it}$ and the blocks $e^{tC_{a,b}}=e^{at}R_{bt}$ computed in \cref{sec:matrix-exponential}, rotations multiplied by exponentials.

For the matrix
\begin{equation}
A=\begin{bmatrix}-1&-4\\1&-1\end{bmatrix},
\label{eq:spiral-example}
\end{equation}
the characteristic polynomial is $(\lambda+1)^2+4$, with the roots $-1\pm2i$. The second row of $A-(-1+2i)I$ is $(1,-2i)$, so $\mathbf w=(2i,1)$ is an eigenvector for $-1+2i$, with $\mathbf p=(0,1)$ and $\mathbf q=(2,0)$, and \eqref{eq:real-complex-solutions} gives the real solutions
\[
\mathbf y_1(t)=e^{-t}\begin{bmatrix}-2\sin2t\\\cos2t\end{bmatrix},\qquad\mathbf y_2(t)=e^{-t}\begin{bmatrix}2\cos2t\\\sin2t\end{bmatrix}.
\]
Without the factor $e^{-t}$ they would run around the ellipse $y_1^2+4y_2^2=4$ with the period $\pi$; with it, they spiral into the origin, the distance shrinking by the factor $e^{-\pi}=0.043$ in every turn.

\subsection{Defective Matrices}

When $A$ is not diagonalisable, there are too few eigenvectors, and the straight-line solutions do not span the solution space. Jordan chains supply the missing solutions, as generalised eigenvectors supply the missing basis vectors in \cref{sec:jordan-form}. We number a chain from its bottom, as in the proof of \cref{thm:jordan-form}: vectors $\mathbf v_1,\ldots,\mathbf v_k$ with
\[
(A-\lambda I)\mathbf v_1=\mathbf 0,\quad\mathbf v_1\ne\mathbf 0,\qquad(A-\lambda I)\mathbf v_j=\mathbf v_{j-1}\quad(j=2,\ldots,k).
\]

\begin{proposition}[Solutions from a Jordan chain]
\label{prop:chain-solutions}
Let $\mathbf v_1,\ldots,\mathbf v_k$ be a Jordan chain of $A$ for the eigenvalue $\lambda$, numbered from the bottom. Then for $j=1,\ldots,k$
\[
\mathbf y_j(t)\coloneqq e^{\lambda t}\Bigl(\mathbf v_j+t\,\mathbf v_{j-1}+\frac{t^2}{2!}\mathbf v_{j-2}+\cdots+\frac{t^{j-1}}{(j-1)!}\mathbf v_1\Bigr)=e^{\lambda t}\sum_{i=0}^{j-1}\frac{t^i}{i!}\mathbf v_{j-i}
\]
is a solution of $\mathbf y'=A\mathbf y$, with $\mathbf y_j(0)=\mathbf v_j$.
\end{proposition}

\begin{proof}
Write $\mathbf y_j=e^{\lambda t}\mathbf p_j(t)$ with the vector polynomial $\mathbf p_j(t)\coloneqq\sum_{i<j}(t^i/i!)\mathbf v_{j-i}$. Then $\mathbf y_j'=\lambda\mathbf y_j+e^{\lambda t}\mathbf p_j'$, so $\mathbf y_j$ is a solution if and only if $\mathbf p_j'=(A-\lambda I)\mathbf p_j$. Differentiating term by term,
\[
\mathbf p_j'(t)=\sum_{i=1}^{j-1}\frac{t^{i-1}}{(i-1)!}\mathbf v_{j-i}=\sum_{i=0}^{j-2}\frac{t^i}{i!}\mathbf v_{j-1-i},
\]
and applying $A-\lambda I$, which maps $\mathbf v_{j-i}$ to $\mathbf v_{j-i-1}$ for $j-i\ge2$ and $\mathbf v_1$ to $\mathbf 0$, gives
\[
(A-\lambda I)\mathbf p_j(t)=\sum_{i=0}^{j-2}\frac{t^i}{i!}\mathbf v_{j-1-i} ,
\]
the same vector polynomial. The value at $0$ is $\mathbf p_j(0)=\mathbf v_j$.
\end{proof}

An eigenvalue $\mu$ is called \term{semisimple} if its algebraic and geometric multiplicities agree, $\operatorname{am}(\mu)=\operatorname{gm}(\mu)$, or equivalently, by \cref{thm:jordan-form}, if all its Jordan blocks have size $1$, so that its index, the size of its largest Jordan block, is $\operatorname{ind}(\mu)=1$ (\cref{sec:jordan-form}).

\begin{theorem}[General solution]
\label{thm:general-solution-systems}
Let $A\in\mathbb C^{n\times n}$ have the pairwise different eigenvalues $\mu_1,\ldots,\mu_k$.
\begin{enumerate}
\item The solutions of \cref{prop:chain-solutions} formed from the chains of a Jordan basis of $A$ are a fundamental set of solutions of $\mathbf y'=A\mathbf y$.
\item Let $\mathbf y_0=\mathbf u_1+\cdots+\mathbf u_k$ with $\mathbf u_i$ in the generalised eigenspace $G_{\mu_i}(A)$. Then
\begin{equation}
e^{tA}\mathbf y_0=\sum_{i=1}^ke^{\mu_it}\sum_{j=0}^{\operatorname{ind}(\mu_i)-1}\frac{t^j}{j!}(A-\mu_iI)^j\mathbf u_i .
\label{eq:general-solution}
\end{equation}
In particular, every component of every solution has the form $\sum_ip_i(t)e^{\mu_it}$ with polynomials $p_i$ of degree less than $\operatorname{ind}(\mu_i)$, and the factors $t^j$ with $j\ge1$ occur only for eigenvalues that are not semisimple.
\end{enumerate}
\end{theorem}

\begin{proof}
(i) By \cref{thm:jordan-form}, $\mathbb C^n$ has a basis consisting of Jordan chains of $A$. Each chain gives solutions by \cref{prop:chain-solutions}, whose values at $0$ are the vectors of the chain; together these values form the Jordan basis, so the solutions are a fundamental set by \cref{thm:system-solution-space}. (ii) The decomposition of $\mathbf y_0$ exists by \cref{thm:generalised-eigenspaces}\,(iii). For each $i$, the matrices $\mu_itI$ and $t(A-\mu_iI)$ commute, so $e^{tA}=e^{\mu_it}e^{t(A-\mu_iI)}$ by \cref{cor:exp-laws}\,(ii). By \cref{thm:jordan-minimal}, $G_{\mu_i}(A)=\ker(A-\mu_iI)^{\operatorname{ind}(\mu_i)}$, so the series of $e^{t(A-\mu_iI)}\mathbf u_i$ stops after the term with $j=\operatorname{ind}(\mu_i)-1$. Adding over $i$ gives \eqref{eq:general-solution}. An eigenvalue is semisimple exactly when its index is $1$, and then its sum has only the term $j=0$.
\end{proof}

Part (ii) is \cref{prop:exp-jordan} made explicit, and it is a practical formula: it needs the generalised eigenspaces and the powers of $A-\mu_iI$, but no Jordan basis.

The defective matrix \eqref{eq:schur-example} of \cref{ch:canonical-forms},
\[
A=\begin{bmatrix}3&1\\-1&1\end{bmatrix},\qquad N\coloneqq A-2I=\begin{bmatrix}1&1\\-1&-1\end{bmatrix},\qquad N^2=0,
\]
has the double eigenvalue $2$ with the single eigenvector $(1,-1)$. The chain $\mathbf v_1=(1,-1)=N\mathbf e_1$, $\mathbf v_2=\mathbf e_1$ gives the fundamental set
\[
\mathbf y_1(t)=e^{2t}\begin{bmatrix}1\\-1\end{bmatrix},\qquad\mathbf y_2(t)=e^{2t}\Bigl(\begin{bmatrix}1\\0\end{bmatrix}+t\begin{bmatrix}1\\-1\end{bmatrix}\Bigr),
\]
and \eqref{eq:general-solution}, with the whole space as the single generalised eigenspace, gives $e^{tA}=e^{2t}(I+tN)$, the matrix found in \cref{sec:matrix-exponential}. For the matrix of order $4$ of the worked example in \cref{sec:jordan-form}, with the eigenvalue $2$ of index $2$ and the simple eigenvalue $-1$, the chain $(0,1,1,1)$, $\mathbf e_1$, the eigenvector $\mathbf e_2$ for $2$ and the eigenvector $(-1,1,1,0)$ for $-1$ give the general solution
\[
\mathbf y(t)=c_1e^{2t}\begin{bmatrix}0\\1\\1\\1\end{bmatrix}+c_2e^{2t}\Bigl(\begin{bmatrix}1\\0\\0\\0\end{bmatrix}+t\begin{bmatrix}0\\1\\1\\1\end{bmatrix}\Bigr)+c_3e^{2t}\begin{bmatrix}0\\1\\0\\0\end{bmatrix}+c_4e^{-t}\begin{bmatrix}-1\\1\\1\\0\end{bmatrix}.
\]
Only the chain of length $2$ produces a factor $t$; the eigenvalue $2$ also has a block of size $1$, whose solution $e^{2t}\mathbf e_2$ has none.

\subsection{Symmetric Matrices}

For a real symmetric matrix the method is at its simplest. By the spectral theorem (\cref{thm:spectral}), $A=Q\Lambda Q^{\mathsf T}$ with an orthogonal $Q$ whose columns $\mathbf q_i$ are orthonormal eigenvectors, and \cref{thm:diagonalisable-systems} gives
\[
\mathbf y(t)=\sum_{i=1}^ne^{\lambda_it}\langle\mathbf q_i,\mathbf y_0\rangle\,\mathbf q_i,\qquad\|\mathbf y(t)\|^2=\sum_{i=1}^ne^{2\lambda_it}\langle\mathbf q_i,\mathbf y_0\rangle^2,
\]
by Pythagoras. The eigenvalues are real, so the solutions do not oscillate, and $e^{\lambda_1t}\|\mathbf y_0\|\le\|\mathbf y(t)\|\le e^{\lambda_nt}\|\mathbf y_0\|$ for $t\ge0$, with the smallest and the largest eigenvalue: the length of every solution is controlled by the extreme eigenvalues alone, a fact that \cref{sec:linear-stability} shows to fail for nonsymmetric matrices. The system is the \term{gradient flow} of the quadratic function $f(\mathbf y)\coloneqq\langle\mathbf y,A\mathbf y\rangle/2$, since $\nabla f(\mathbf y)=A\mathbf y$ for symmetric $A$: the solutions move in the direction of steepest ascent of $f$, and $f(\mathbf y(t))$ increases, with $(f\circ\mathbf y)'=\langle\nabla f(\mathbf y),\mathbf y'\rangle=\|A\mathbf y\|^2$.

Three bodies in a row exchange heat with their neighbours, at a rate proportional to the difference of the temperatures, and with no loss to the surroundings. With the unit constant of proportionality the temperatures satisfy $\mathbf y'=-L\mathbf y$ with
\[
L=\begin{bmatrix}1&-1&0\\-1&2&-1\\0&-1&1\end{bmatrix},
\]
which has the eigenvalues $0$, $1$, $3$ with the orthonormal eigenvectors $(1,1,1)/\sqrt3$, $(1,0,-1)/\sqrt2$ and $(1,-2,1)/\sqrt6$. So
\[
\mathbf y(t)=\bar y\begin{bmatrix}1\\1\\1\end{bmatrix}+\frac{y_1(0)-y_3(0)}2\,e^{-t}\begin{bmatrix}1\\0\\-1\end{bmatrix}+\frac{y_1(0)-2y_2(0)+y_3(0)}6\,e^{-3t}\begin{bmatrix}1\\-2\\1\end{bmatrix},
\]
where $\bar y$ is the mean of the initial temperatures. All temperatures approach the common value $\bar y$, the sum of the temperatures being conserved because the columns of $L$ add up to zero; the difference between the two ends decays at the rate $1$, the bending of the profile at the rate $3$. The matrix $L$ is the Laplacian of a path, of the kind that \cref{sec:structured-eigen} met for graphs.

\subsection{Computing the Exponential by Hand}

For small matrices, Cayley--Hamilton gives closed formulas for $e^{tA}$ that avoid eigenvectors altogether.

\begin{proposition}[Exponential of a $2\times2$ matrix]
\label{prop:exp-2x2}
Let $A\in\mathbb C^{2\times2}$ have the eigenvalues $\lambda_1$ and $\lambda_2$.
\begin{enumerate}
\item If $\lambda_1\ne\lambda_2$, then
\[
e^{tA}=\frac{e^{\lambda_1t}(A-\lambda_2I)-e^{\lambda_2t}(A-\lambda_1I)}{\lambda_1-\lambda_2}.
\]
\item If $\lambda_1=\lambda_2=\lambda$, then $e^{tA}=e^{\lambda t}\bigl(I+t(A-\lambda I)\bigr)$.
\item If $A$ is real with the eigenvalues $a\pm ib$, $b>0$, then
\[
e^{tA}=e^{at}\Bigl(\cos(bt)\,I+\frac{\sin bt}b\,(A-aI)\Bigr).
\]
\end{enumerate}
\end{proposition}

\begin{proof}
By the Cayley--Hamilton theorem (\cref{thm:cayley-hamilton}), $(A-\lambda_1I)(A-\lambda_2I)=0$. (i) Let $M(t)$ be the right side. Then $M(0)=((A-\lambda_2I)-(A-\lambda_1I))/(\lambda_1-\lambda_2)=I$. Since $A(A-\lambda_2I)=(A-\lambda_1I)(A-\lambda_2I)+\lambda_1(A-\lambda_2I)=\lambda_1(A-\lambda_2I)$ and, in the same way, $A(A-\lambda_1I)=\lambda_2(A-\lambda_1I)$,
\[
AM(t)=\frac{\lambda_1e^{\lambda_1t}(A-\lambda_2I)-\lambda_2e^{\lambda_2t}(A-\lambda_1I)}{\lambda_1-\lambda_2}=M'(t).
\]
So $M$ is the matrix solution with $M(0)=I$, which is $e^{tA}$ by \cref{thm:linear-system}, applied to each column. (ii) Now $(A-\lambda I)^2=0$, and $e^{tA}=e^{\lambda t}e^{t(A-\lambda I)}=e^{\lambda t}(I+t(A-\lambda I))$ by \cref{cor:exp-laws}\,(ii), the series of the second factor stopping after two terms. (iii) Apply (i) with $\lambda_{1,2}=a\pm ib$, so that $\lambda_1-\lambda_2=2ib$ and $A-\lambda_{2,1}I=(A-aI)\pm ibI$:
\[
\begin{aligned}
e^{tA}&=\frac{e^{at}}{2ib}\Bigl(e^{ibt}\bigl((A-aI)+ibI\bigr)-e^{-ibt}\bigl((A-aI)-ibI\bigr)\Bigr)\\
&=e^{at}\Bigl(\frac{e^{ibt}-e^{-ibt}}{2ib}(A-aI)+\frac{e^{ibt}+e^{-ibt}}2I\Bigr),
\end{aligned}
\]
which is the claim, by \cref{cor:euler-consequences}.
\end{proof}

Formula (i) is Sylvester's formula, the case $n=2$ of the Lagrange form of $f(A)$ in \cref{thm:function-polynomial}, and the same method works for every order: $e^{tA}$ is the polynomial in $A$ of degree less than $n$ that interpolates $e^{\lambda t}$ on the spectrum (\cref{sec:matrix-functions}). For the spiral \eqref{eq:spiral-example}, (iii) with $a=-1$ and $b=2$ gives
\[
e^{tA}=e^{-t}\Bigl(\cos2t\,I+\frac{\sin2t}2\begin{bmatrix}0&-4\\1&0\end{bmatrix}\Bigr)=e^{-t}\begin{bmatrix}\cos2t&-2\sin2t\\\sin(2t)/2&\cos2t\end{bmatrix},
\]
whose first column is $\mathbf y_2/2$ and whose second column is $\mathbf y_1$, with the solutions $\mathbf y_1$ and $\mathbf y_2$ found above. In summary, there are three ways to the solutions by hand: eigenvectors and Jordan chains (\cref{thm:diagonalisable-systems,thm:general-solution-systems}), the generalised eigenspaces of \eqref{eq:general-solution}, and interpolation formulas such as \cref{prop:exp-2x2}. In floating-point arithmetic none of them is the method of choice; \cref{sec:computing-matrix-functions} showed that eigenvectors can be ill-conditioned and that $e^{tA}$ is best computed by scaling and squaring, and \cref{sec:systems-numerics} uses the computed exponential to step through time.

\section{The Phase Plane}
\label{sec:phase-plane}

\subsection{Orbits and Phase Portraits}

A solution of a system of two equations is a pair of functions, and it can be drawn in two ways: as two graphs over the time axis, or as a single curve $t\mapsto(y_1(t),y_2(t))$ in the plane, traced as time runs. The second way loses the time scale but shows the interplay of the two components, and the pictures of many curves together reveal the structure of all solutions at once. \Cref{sec:exact} used it for the conservation of energy, and \cref{sec:higher-numerics} for numerical methods.

For an \term{autonomous} system $\mathbf y'=\mathbf f(\mathbf y)$, in which the right side does not depend on $t$, the plane of the states $\mathbf y$ is the \term{phase plane} (the \term{phase space} for $n>2$), the curve traced by a solution is its \term{orbit} or trajectory, and a picture of representative orbits, with arrows that show the direction of motion, is a \term{phase portrait}. A point $\mathbf y^*$ with $\mathbf f(\mathbf y^*)=\mathbf 0$ is an \term{equilibrium}: the constant function $\mathbf y^*$ is a solution, and its orbit is a single point. The phase portrait is the system's analogue of the phase line of \cref{sec:autonomous}, and the following lemma makes it a meaningful picture.

\begin{lemma}[Orbits of autonomous systems]
\label{lem:orbits}
Let $U\subseteq\mathbb R^n$ be open, and let $\mathbf f\colon U\to\mathbb R^n$ be continuously differentiable.
\begin{enumerate}
\item If $\mathbf y$ is a solution of $\mathbf y'=\mathbf f(\mathbf y)$ on an interval $J$, then $t\mapsto\mathbf y(t-c)$ is a solution on $J+c$ for every $c\in\mathbb R$.
\item The orbits of two maximal solutions are either equal or disjoint. In particular, orbits do not cross, and a nonconstant solution never reaches an equilibrium in finite time.
\item If a maximal solution satisfies $\mathbf y(t_1)=\mathbf y(t_2)$ for some $t_1<t_2$, it is defined on $\mathbb R$ and periodic with the period $t_2-t_1$, and its orbit is a closed curve.
\end{enumerate}
\end{lemma}

\begin{proof}
(i) By the chain rule, the derivative of $\mathbf y(t-c)$ at $t$ is $\mathbf y'(t-c)=\mathbf f(\mathbf y(t-c))$. (ii) Let $\mathbf y$ and $\mathbf z$ be maximal solutions with $\mathbf y(t_1)=\mathbf z(t_2)$. By (i), $\mathbf w(t)\coloneqq\mathbf z(t+t_2-t_1)$ is a maximal solution with $\mathbf w(t_1)=\mathbf y(t_1)$, and by the uniqueness of maximal solutions (\cref{thm:maximal-solution}, in the form for systems) $\mathbf w=\mathbf y$; so the orbit of $\mathbf z$, which is that of $\mathbf w$, is the orbit of $\mathbf y$. An equilibrium is the orbit of a constant solution, so a solution that reaches it is constant. (iii) With $T\coloneqq t_2-t_1$, the function $\mathbf y(t+T)$ is a maximal solution with the same value at $t_1$ as $\mathbf y$, so it is $\mathbf y$; in particular its interval $(\omega_--T,\omega_+-T)$ equals $(\omega_-,\omega_+)$, which is possible only if both ends are infinite.
\end{proof}

For a linear system $\mathbf y'=A\mathbf y$ with a real $2\times2$ matrix $A$, the origin is an equilibrium, and it is the only one if $\det A\ne0$. Linearity adds two symmetries: if $\mathbf y(t)$ is a solution, so are $c\,\mathbf y(t)$ for every real $c$, so the portrait looks the same at every scale, and $\mathbf y(-t)$ is a solution of $\mathbf y'=-A\mathbf y$, so the portrait of $-A$ is that of $A$ with all arrows reversed.

\subsection{Similar Systems and the Real Canonical Forms}

If $A=SBS^{-1}$ with a real invertible matrix $S$, then $\mathbf y=S\mathbf z$ solves $\mathbf y'=A\mathbf y$ if and only if $\mathbf z$ solves $\mathbf z'=B\mathbf z$, since $\mathbf y'=S\mathbf z'$ and $A\mathbf y=SB\mathbf z$. The phase portrait of $A$ is therefore the image of the phase portrait of $B$ under the linear map $S$. A linear map of the plane distorts lengths and angles, turning circles into ellipses, but it maps lines through the origin to lines through the origin, preserves tangency and the direction of motion, and keeps a curve that winds around the origin winding around it. So it suffices to understand the portraits of a few simple matrices.

\begin{proposition}[Real canonical forms of order $2$]
\label{prop:real-forms-2x2}
Every real $2\times2$ matrix $A$ is similar, by a real invertible matrix, to exactly one of the following:
\begin{enumerate}
\item $\operatorname{diag}(\lambda_1,\lambda_2)$ with real $\lambda_1\ge\lambda_2$, if $A$ has real eigenvalues and is diagonalisable;
\item the Jordan block $J_2(\lambda)$, if $A$ has a double real eigenvalue $\lambda$ and is not diagonalisable;
\item $C_{a,b}$ with $b>0$, if $A$ has the nonreal eigenvalues $a\pm ib$.
\end{enumerate}
\end{proposition}

\begin{proof}
If the eigenvalues are real and $A$ is diagonalisable, the eigenvectors can be chosen real, since $A-\lambda I$ is real, and form the columns of $S$. If they are real and $A$ is not diagonalisable, the eigenvalue $\lambda$ is double, $(A-\lambda I)^2=0$ by the Cayley--Hamilton theorem (\cref{thm:cayley-hamilton}), and $A-\lambda I\ne0$; for a real $\mathbf v_2$ with $\mathbf v_1\coloneqq(A-\lambda I)\mathbf v_2\ne\mathbf 0$, the real vectors $\mathbf v_1$, $\mathbf v_2$ form a Jordan chain, they are independent because $\mathbf v_1$ lies in the kernel of $A-\lambda I$ and $\mathbf v_2$ does not, and the matrix with the columns $\mathbf v_1$, $\mathbf v_2$ transforms $A$ into $J_2(\lambda)$. If they are nonreal, \cref{prop:real-2x2-form}, applied with an eigenvector for $a-ib$, gives the similarity to $C_{a,b}$. The three forms are not similar to each other, since similar matrices have the same eigenvalues and either both or neither are diagonalisable, and within each type the eigenvalues determine the matrix, with the order $\lambda_1\ge\lambda_2$ and the sign $b>0$ fixed.
\end{proof}

\subsection{Classification of Equilibria}

\begin{theorem}[Phase portraits of linear systems]
\label{thm:phase-classification}
Let $A\in\mathbb R^{2\times2}$ with $\det A\ne0$, and let $\lambda_1$, $\lambda_2$ be its eigenvalues. Exactly one of the following holds.
\begin{enumerate}
\item \emph{Saddle}: $\lambda_2<0<\lambda_1$. The lines spanned by the eigenvectors are unions of orbits; solutions on the line of $\lambda_2$ tend to $\mathbf 0$, those on the line of $\lambda_1$ tend to infinity, and every other solution tends to infinity as $t\to\pm\infty$, its direction approaching that of the eigenvector for $\lambda_1$ as $t\to\infty$ and that for $\lambda_2$ as $t\to-\infty$.
\item \emph{Node}: $\lambda_1$ and $\lambda_2$ are real, different and of the same sign. If they are negative, every solution tends to $\mathbf 0$, and every solution that does not lie on the eigenvector line of the eigenvalue of larger modulus approaches $\mathbf 0$ tangentially to the eigenvector line of the eigenvalue of smaller modulus. If they are positive, the same holds as $t\to-\infty$.
\item \emph{Star}: $A=\lambda I$. Every orbit other than the origin is an open ray from the origin.
\item \emph{Degenerate node}: $\lambda_1=\lambda_2=\lambda$ and $A\ne\lambda I$. All orbits other than the origin approach the origin as $t\to\infty$ if $\lambda<0$, and as $t\to-\infty$ if $\lambda>0$, tangentially to the single eigenvector line.
\item \emph{Spiral}: $\lambda_{1,2}=a\pm ib$ with $a\ne0$ and $b>0$. Every solution other than $\mathbf 0$ winds around the origin infinitely often, tending to $\mathbf 0$ as $t\to\infty$ if $a<0$ and as $t\to-\infty$ if $a>0$.
\item \emph{Centre}: $\lambda_{1,2}=\pm ib$ with $b>0$. Every orbit other than the origin is an ellipse centred at the origin, and every solution is periodic with the period $2\pi/b$.
\end{enumerate}
\end{theorem}

\begin{proof}
By \cref{prop:real-forms-2x2}, $A=SBS^{-1}$ with a real $S$ and one of the three canonical forms $B$, and $\mathbf y=S\mathbf z$ with $\mathbf z'=B\mathbf z$. Statements about limits, about lines through the origin and about directions transfer from $\mathbf z$ to $\mathbf y$: if $\mathbf z(t)/\|\mathbf z(t)\|\to\mathbf d$, then $\mathbf y(t)/\|\mathbf y(t)\|\to S\mathbf d/\|S\mathbf d\|$, by continuity, and $S$ maps the eigenvectors of $B$ to those of $A$. So we may assume $A=B$; we write $c_1$ and $c_2$ for the components of $\mathbf z(0)$.

(i) and (ii): $B=\operatorname{diag}(\lambda_1,\lambda_2)$, the solutions are $\mathbf z(t)=(c_1e^{\lambda_1t},c_2e^{\lambda_2t})$, and the coordinate axes, spanned by the eigenvectors, are unions of orbits. In case (i), with $c_1c_2\ne0$, the first component tends to infinity and the second to $0$ as $t\to\infty$, so $\|\mathbf z\|\to\infty$ and $\mathbf z/\|\mathbf z\|\to(\operatorname{sign}c_1,0)$; as $t\to-\infty$ the roles are exchanged. In case (ii) with negative eigenvalues, let $|\lambda_2|>|\lambda_1|$, so that $\lambda_2<\lambda_1<0$. Every solution tends to $\mathbf 0$. If $c_1\ne0$, the ratio $z_2/z_1=(c_2/c_1)e^{(\lambda_2-\lambda_1)t}$ tends to $0$, so the direction of $\mathbf z(t)$ approaches that of $\pm\mathbf e_1$, the eigenvector of the eigenvalue of smaller modulus; if $c_1=0$, the solution lies on the other axis. Positive eigenvalues are the case of $-B$ with time reversed.

(iii) $\mathbf z(t)=e^{\lambda t}\mathbf z(0)$ runs along the ray through $\mathbf z(0)$.

(iv) $B=J_2(\lambda)$, and by \cref{prop:exp-2x2}\,(ii) $\mathbf z(t)=e^{\lambda t}(c_1+c_2t,c_2)$. If $c_2=0$, the solution lies on the eigenvector line $\mathbb R\mathbf e_1$. Otherwise $z_1/z_2=c_1/c_2+t$ tends to $\pm\infty$ as $t\to\pm\infty$, so the direction approaches that of $\pm\mathbf e_1$ at both ends, and $\|\mathbf z(t)\|\le e^{\lambda t}(|c_1|+|c_2|(1+|t|))$ tends to $0$ as $t\to\infty$ if $\lambda<0$, by \cref{prop:exp-growth}.

(v) and (vi): $B=C_{a,b}$, and $\mathbf z(t)=e^{tC_{a,b}}\mathbf z(0)=e^{at}R_{bt}\mathbf z(0)$ (\cref{sec:matrix-exponential}): the point $\mathbf z(0)$ is rotated by the angle $bt$ and its length multiplied by $e^{at}$. For $a\ne0$ the orbit is a spiral that winds around the origin once in each time interval of length $2\pi/b$, and $\|\mathbf z(t)\|=e^{at}\|\mathbf z(0)\|$ tends to $0$ in the stated direction of time. For $a=0$ the orbit is the circle of radius $r\coloneqq\|\mathbf z(0)\|$, run through with the period $2\pi/b$, and its image under $S$ is the curve $\langle\mathbf y,G\mathbf y\rangle=r^2$ with the positive definite matrix $G\coloneqq S^{-\mathsf T}S^{-1}$, an ellipse centred at the origin (\cref{prop:ellipsoid}).
\end{proof}

\Cref{fig:phase-portraits} shows the six types for matrices that are not in canonical form, so that the eigenvector lines are oblique. In the saddle the two eigenvector lines, called the \term{stable} and the \term{unstable line}, divide the plane into four regions, and the other orbits come in along the stable line and leave along the unstable one, like hyperbolas; in the canonical coordinates they satisfy $|z_2|=C|z_1|^{\lambda_2/\lambda_1}$. In the node, the eigenvector of the eigenvalue of smaller modulus is the \term{slow direction}: the fast component dies out first, and the solutions arrive along the slow direction, as the salt in the tanks settled into the ratio $\sqrt2$. The names refer to the behaviour near the equilibrium as seen in the picture; the stability of the equilibrium, defined in \cref{sec:linear-stability}, depends on the direction of the arrows.

\begin{figure}
\sdpanel{A}{\begin{tikzpicture}
\begin{axis}[sdplot, width=0.31\textwidth, height=0.31\textwidth, xmin=-1, xmax=1, ymin=-1, ymax=1, clip=true, xtick={-1,0,1}, ytick={-1,0,1}, xlabel={$y_1$}, ylabel={$y_2$}]
\addplot[ch25orbit] coordinates {(0,0) (0.671,1.342)};
\draw[ch25arrow] (axis cs:0.271,0.541) -- (axis cs:0.275,0.550);
\addplot[ch25orbit] coordinates {(0,0) (-0.671,-1.342)};
\draw[ch25arrow] (axis cs:-0.271,-0.541) -- (axis cs:-0.275,-0.550);
\addplot[ch25orbit] coordinates {(0,0) (1.061,-1.061)};
\draw[ch25arrow] (axis cs:0.557,-0.557) -- (axis cs:0.550,-0.550);
\addplot[ch25orbit] coordinates {(0,0) (-1.061,1.061)};
\draw[ch25arrow] (axis cs:-0.557,0.557) -- (axis cs:-0.550,0.550);
\addplot[ch25orbit] coordinates {(1.041,-1.010) (0.992,-0.957) (0.944,-0.905) (0.896,-0.853) (0.848,-0.800) (0.802,-0.747) (0.756,-0.694) (0.710,-0.639) (0.666,-0.584) (0.623,-0.527) (0.582,-0.469) (0.544,-0.410) (0.510,-0.348) (0.479,-0.284) (0.455,-0.217) (0.437,-0.149) (0.426,-0.078) (0.423,-0.008) (0.425,0.064) (0.432,0.135) (0.444,0.206) (0.459,0.277) (0.476,0.346) (0.495,0.414) (0.516,0.481) (0.538,0.550) (0.562,0.617) (0.586,0.685) (0.612,0.755) (0.639,0.824) (0.667,0.893) (0.694,0.962) (0.722,1.029) (0.728,1.044)};
\draw[ch25arrow] (axis cs:0.426,-0.068) -- (axis cs:0.425,-0.060);
\addplot[ch25orbit] coordinates {(1.041,-0.968) (0.994,-0.913) (0.948,-0.858) (0.903,-0.803) (0.859,-0.748) (0.816,-0.690) (0.775,-0.633) (0.736,-0.573) (0.699,-0.512) (0.665,-0.451) (0.635,-0.387) (0.609,-0.322) (0.587,-0.255) (0.571,-0.185) (0.560,-0.115) (0.555,-0.044) (0.554,0.026) (0.558,0.096) (0.565,0.166) (0.575,0.236) (0.588,0.305) (0.604,0.376) (0.621,0.444) (0.640,0.514) (0.661,0.584) (0.684,0.655) (0.707,0.725) (0.731,0.796) (0.756,0.865) (0.781,0.934) (0.807,1.000) (0.823,1.042)};
\draw[ch25arrow] (axis cs:0.554,-0.022) -- (axis cs:0.554,-0.014);
\addplot[ch25orbit] coordinates {(1.040,-0.914) (0.996,-0.857) (0.953,-0.798) (0.912,-0.739) (0.872,-0.679) (0.834,-0.618) (0.799,-0.555) (0.767,-0.490) (0.738,-0.425) (0.714,-0.360) (0.693,-0.291) (0.677,-0.221) (0.666,-0.152) (0.659,-0.082) (0.657,-0.011) (0.658,0.061) (0.663,0.132) (0.671,0.204) (0.681,0.274) (0.694,0.345) (0.709,0.416) (0.727,0.488) (0.744,0.555) (0.764,0.627) (0.786,0.698) (0.808,0.767) (0.830,0.836) (0.853,0.902) (0.877,0.971) (0.901,1.038) (0.903,1.044)};
\draw[ch25arrow] (axis cs:0.657,0.025) -- (axis cs:0.657,0.033);
\addplot[ch25orbit] coordinates {(1.008,-1.040) (0.955,-0.991) (0.903,-0.943) (0.849,-0.894) (0.795,-0.846) (0.741,-0.799) (0.687,-0.754) (0.631,-0.709) (0.575,-0.666) (0.517,-0.626) (0.456,-0.590) (0.392,-0.559) (0.324,-0.538) (0.253,-0.533) (0.184,-0.546) (0.120,-0.575) (0.062,-0.615) (0.008,-0.663) (-0.041,-0.715) (-0.087,-0.769) (-0.132,-0.825) (-0.175,-0.883) (-0.217,-0.943) (-0.258,-1.004) (-0.282,-1.042)};
\draw[ch25arrow] (axis cs:0.343,-0.544) -- (axis cs:0.335,-0.542);
\addplot[ch25orbit] coordinates {(0.965,-1.042) (0.910,-0.995) (0.855,-0.951) (0.799,-0.908) (0.743,-0.866) (0.684,-0.825) (0.624,-0.788) (0.562,-0.755) (0.497,-0.727) (0.430,-0.707) (0.360,-0.698) (0.290,-0.703) (0.222,-0.722) (0.158,-0.752) (0.100,-0.791) (0.044,-0.837) (-0.007,-0.886) (-0.056,-0.939) (-0.103,-0.993) (-0.141,-1.041)};
\draw[ch25arrow] (axis cs:0.406,-0.704) -- (axis cs:0.398,-0.703);
\addplot[ch25orbit] coordinates {(0.901,-1.041) (0.844,-1.001) (0.785,-0.962) (0.725,-0.925) (0.661,-0.892) (0.596,-0.864) (0.528,-0.842) (0.458,-0.830) (0.387,-0.828) (0.317,-0.839) (0.249,-0.862) (0.186,-0.893) (0.125,-0.933) (0.070,-0.977) (0.018,-1.024) (0.002,-1.041)};
\draw[ch25arrow] (axis cs:0.454,-0.829) -- (axis cs:0.446,-0.829);
\addplot[ch25orbit] coordinates {(-1.008,1.040) (-0.955,0.991) (-0.903,0.943) (-0.849,0.894) (-0.795,0.846) (-0.741,0.799) (-0.687,0.754) (-0.631,0.709) (-0.575,0.666) (-0.517,0.626) (-0.456,0.590) (-0.392,0.559) (-0.324,0.538) (-0.253,0.533) (-0.184,0.546) (-0.120,0.575) (-0.062,0.615) (-0.008,0.663) (0.041,0.715) (0.087,0.769) (0.132,0.825) (0.175,0.883) (0.217,0.943) (0.258,1.004) (0.282,1.042)};
\draw[ch25arrow] (axis cs:-0.343,0.544) -- (axis cs:-0.335,0.542);
\addplot[ch25orbit] coordinates {(-0.965,1.042) (-0.910,0.995) (-0.855,0.951) (-0.799,0.908) (-0.743,0.866) (-0.684,0.825) (-0.624,0.788) (-0.562,0.755) (-0.497,0.727) (-0.430,0.707) (-0.360,0.698) (-0.290,0.703) (-0.222,0.722) (-0.158,0.752) (-0.100,0.791) (-0.044,0.837) (0.007,0.886) (0.056,0.939) (0.103,0.993) (0.141,1.041)};
\draw[ch25arrow] (axis cs:-0.406,0.704) -- (axis cs:-0.398,0.703);
\addplot[ch25orbit] coordinates {(-0.901,1.041) (-0.844,1.001) (-0.785,0.962) (-0.725,0.925) (-0.661,0.892) (-0.596,0.864) (-0.528,0.842) (-0.458,0.830) (-0.387,0.828) (-0.317,0.839) (-0.249,0.862) (-0.186,0.893) (-0.125,0.933) (-0.070,0.977) (-0.018,1.024) (-0.002,1.041)};
\draw[ch25arrow] (axis cs:-0.454,0.829) -- (axis cs:-0.446,0.829);
\addplot[ch25orbit] coordinates {(-1.041,1.010) (-0.992,0.957) (-0.944,0.905) (-0.896,0.853) (-0.848,0.800) (-0.802,0.747) (-0.756,0.694) (-0.710,0.639) (-0.666,0.584) (-0.623,0.527) (-0.582,0.469) (-0.544,0.410) (-0.510,0.348) (-0.479,0.284) (-0.455,0.217) (-0.437,0.149) (-0.426,0.078) (-0.423,0.008) (-0.425,-0.064) (-0.432,-0.135) (-0.444,-0.206) (-0.459,-0.277) (-0.476,-0.346) (-0.495,-0.414) (-0.516,-0.481) (-0.538,-0.550) (-0.562,-0.617) (-0.586,-0.685) (-0.612,-0.755) (-0.639,-0.824) (-0.667,-0.893) (-0.694,-0.962) (-0.722,-1.029) (-0.728,-1.044)};
\draw[ch25arrow] (axis cs:-0.426,0.068) -- (axis cs:-0.425,0.060);
\addplot[ch25orbit] coordinates {(-1.041,0.968) (-0.994,0.913) (-0.948,0.858) (-0.903,0.803) (-0.859,0.748) (-0.816,0.690) (-0.775,0.633) (-0.736,0.573) (-0.699,0.512) (-0.665,0.451) (-0.635,0.387) (-0.609,0.322) (-0.587,0.255) (-0.571,0.185) (-0.560,0.115) (-0.555,0.044) (-0.554,-0.026) (-0.558,-0.096) (-0.565,-0.166) (-0.575,-0.236) (-0.588,-0.305) (-0.604,-0.376) (-0.621,-0.444) (-0.640,-0.514) (-0.661,-0.584) (-0.684,-0.655) (-0.707,-0.725) (-0.731,-0.796) (-0.756,-0.865) (-0.781,-0.934) (-0.807,-1.000) (-0.823,-1.042)};
\draw[ch25arrow] (axis cs:-0.554,0.022) -- (axis cs:-0.554,0.014);
\addplot[ch25orbit] coordinates {(-1.040,0.914) (-0.996,0.857) (-0.953,0.798) (-0.912,0.739) (-0.872,0.679) (-0.834,0.618) (-0.799,0.555) (-0.767,0.490) (-0.738,0.425) (-0.714,0.360) (-0.693,0.291) (-0.677,0.221) (-0.666,0.152) (-0.659,0.082) (-0.657,0.011) (-0.658,-0.061) (-0.663,-0.132) (-0.671,-0.204) (-0.681,-0.274) (-0.694,-0.345) (-0.709,-0.416) (-0.727,-0.488) (-0.744,-0.555) (-0.764,-0.627) (-0.786,-0.698) (-0.808,-0.767) (-0.830,-0.836) (-0.853,-0.902) (-0.877,-0.971) (-0.901,-1.038) (-0.903,-1.044)};
\draw[ch25arrow] (axis cs:-0.657,-0.025) -- (axis cs:-0.657,-0.033);
\fill[black] (axis cs:0,0) circle (1.3pt);
\end{axis}
\end{tikzpicture}}\hfill
\sdpanel{B}{\begin{tikzpicture}
\begin{axis}[sdplot, width=0.31\textwidth, height=0.31\textwidth, xmin=-1, xmax=1, ymin=-1, ymax=1, clip=true, xtick={-1,0,1}, ytick={-1,0,1}, xlabel={$y_1$}, ylabel={$y_2$}]
\addplot[ch25orbit] coordinates {(0,0) (1.342,0.671)};
\draw[ch25arrow] (axis cs:0.559,0.279) -- (axis cs:0.550,0.275);
\addplot[ch25orbit] coordinates {(0,0) (-1.342,-0.671)};
\draw[ch25arrow] (axis cs:-0.559,-0.279) -- (axis cs:-0.550,-0.275);
\addplot[ch25orbit] coordinates {(0,0) (0.000,1.500)};
\draw[ch25arrow] (axis cs:0.000,0.560) -- (axis cs:0.000,0.550);
\addplot[ch25orbit] coordinates {(0,0) (-0.000,-1.500)};
\draw[ch25arrow] (axis cs:-0.000,-0.560) -- (axis cs:-0.000,-0.550);
\addplot[ch25orbit] coordinates {(-1.016,1.041) (-0.998,0.968) (-0.980,0.900) (-0.961,0.829) (-0.940,0.756) (-0.919,0.688) (-0.897,0.619) (-0.874,0.550) (-0.850,0.483) (-0.825,0.417) (-0.798,0.350) (-0.768,0.285) (-0.736,0.221) (-0.702,0.159) (-0.664,0.100) (-0.621,0.043) (-0.573,-0.008) (-0.519,-0.053) (-0.457,-0.087) (-0.390,-0.107) (-0.320,-0.112) (-0.250,-0.102) (-0.183,-0.082) (-0.117,-0.056) (-0.054,-0.027) (-0.027,-0.013)};
\draw[ch25arrow] (axis cs:-0.843,0.464) -- (axis cs:-0.840,0.457);
\addplot[ch25orbit] coordinates {(-1.041,0.426) (-1.010,0.359) (-0.978,0.296) (-0.943,0.233) (-0.906,0.172) (-0.866,0.112) (-0.822,0.055) (-0.774,0.002) (-0.722,-0.045) (-0.664,-0.087) (-0.600,-0.119) (-0.532,-0.140) (-0.463,-0.148) (-0.393,-0.146) (-0.324,-0.133) (-0.256,-0.114) (-0.190,-0.089) (-0.126,-0.061) (-0.063,-0.031) (-0.027,-0.013)};
\draw[ch25arrow] (axis cs:-0.807,0.039) -- (axis cs:-0.802,0.033);
\addplot[ch25orbit] coordinates {(-1.041,-0.532) (-0.978,-0.499) (-0.916,-0.466) (-0.852,-0.433) (-0.790,-0.400) (-0.727,-0.368) (-0.665,-0.336) (-0.602,-0.303) (-0.539,-0.271) (-0.476,-0.239) (-0.413,-0.207) (-0.350,-0.175) (-0.287,-0.144) (-0.224,-0.112) (-0.162,-0.081) (-0.099,-0.050) (-0.036,-0.018) (-0.027,-0.013)};
\draw[ch25arrow] (axis cs:-0.695,-0.351) -- (axis cs:-0.688,-0.347);
\addplot[ch25orbit] coordinates {(-0.410,-1.043) (-0.398,-0.970) (-0.387,-0.898) (-0.374,-0.827) (-0.361,-0.755) (-0.347,-0.683) (-0.332,-0.612) (-0.316,-0.542) (-0.299,-0.474) (-0.280,-0.406) (-0.258,-0.339) (-0.234,-0.272) (-0.205,-0.208) (-0.171,-0.146) (-0.129,-0.090) (-0.076,-0.044) (-0.027,-0.014)};
\draw[ch25arrow] (axis cs:-0.342,-0.661) -- (axis cs:-0.341,-0.653);
\addplot[ch25orbit] coordinates {(0.407,-1.042) (0.398,-0.967) (0.390,-0.897) (0.380,-0.825) (0.371,-0.754) (0.360,-0.683) (0.349,-0.610) (0.337,-0.540) (0.325,-0.470) (0.311,-0.399) (0.296,-0.329) (0.279,-0.259) (0.259,-0.191) (0.236,-0.124) (0.207,-0.060) (0.167,-0.002) (0.105,0.031) (0.036,0.017) (0.027,0.013)};
\draw[ch25arrow] (axis cs:0.351,-0.624) -- (axis cs:0.350,-0.616);
\addplot[ch25orbit] coordinates {(1.016,-1.041) (0.998,-0.968) (0.980,-0.900) (0.961,-0.829) (0.940,-0.756) (0.919,-0.688) (0.897,-0.619) (0.874,-0.550) (0.850,-0.483) (0.825,-0.417) (0.798,-0.350) (0.768,-0.285) (0.736,-0.221) (0.702,-0.159) (0.664,-0.100) (0.621,-0.043) (0.573,0.008) (0.519,0.053) (0.457,0.087) (0.390,0.107) (0.320,0.112) (0.250,0.102) (0.183,0.082) (0.117,0.056) (0.054,0.027) (0.027,0.013)};
\draw[ch25arrow] (axis cs:0.843,-0.464) -- (axis cs:0.840,-0.457);
\addplot[ch25orbit] coordinates {(1.041,-0.426) (1.010,-0.359) (0.978,-0.296) (0.943,-0.233) (0.906,-0.172) (0.866,-0.112) (0.822,-0.055) (0.774,-0.002) (0.722,0.045) (0.664,0.087) (0.600,0.119) (0.532,0.140) (0.463,0.148) (0.393,0.146) (0.324,0.133) (0.256,0.114) (0.190,0.089) (0.126,0.061) (0.063,0.031) (0.027,0.013)};
\draw[ch25arrow] (axis cs:0.807,-0.039) -- (axis cs:0.802,-0.033);
\addplot[ch25orbit] coordinates {(1.041,0.532) (0.978,0.499) (0.916,0.466) (0.852,0.433) (0.790,0.400) (0.727,0.368) (0.665,0.336) (0.602,0.303) (0.539,0.271) (0.476,0.239) (0.413,0.207) (0.350,0.175) (0.287,0.144) (0.224,0.112) (0.162,0.081) (0.099,0.050) (0.036,0.018) (0.027,0.013)};
\draw[ch25arrow] (axis cs:0.695,0.351) -- (axis cs:0.688,0.347);
\addplot[ch25orbit] coordinates {(0.410,1.043) (0.398,0.970) (0.387,0.898) (0.374,0.827) (0.361,0.755) (0.347,0.683) (0.332,0.612) (0.316,0.542) (0.299,0.474) (0.280,0.406) (0.258,0.339) (0.234,0.272) (0.205,0.208) (0.171,0.146) (0.129,0.090) (0.076,0.044) (0.027,0.014)};
\draw[ch25arrow] (axis cs:0.342,0.661) -- (axis cs:0.341,0.653);
\addplot[ch25orbit] coordinates {(-0.407,1.042) (-0.398,0.967) (-0.390,0.897) (-0.380,0.825) (-0.371,0.754) (-0.360,0.683) (-0.349,0.610) (-0.337,0.540) (-0.325,0.470) (-0.311,0.399) (-0.296,0.329) (-0.279,0.259) (-0.259,0.191) (-0.236,0.124) (-0.207,0.060) (-0.167,0.002) (-0.105,-0.031) (-0.036,-0.017) (-0.027,-0.013)};
\draw[ch25arrow] (axis cs:-0.351,0.624) -- (axis cs:-0.350,0.616);
\fill[black] (axis cs:0,0) circle (1.3pt);
\end{axis}
\end{tikzpicture}}\hfill
\sdpanel{C}{\begin{tikzpicture}
\begin{axis}[sdplot, width=0.31\textwidth, height=0.31\textwidth, xmin=-1, xmax=1, ymin=-1, ymax=1, clip=true, xtick={-1,0,1}, ytick={-1,0,1}, xlabel={$y_1$}, ylabel={$y_2$}]
\addplot[ch25orbit] coordinates {(-1.040,-0.642) (-0.979,-0.677) (-0.916,-0.710) (-0.851,-0.739) (-0.784,-0.766) (-0.716,-0.790) (-0.647,-0.810) (-0.576,-0.828) (-0.505,-0.843) (-0.433,-0.855) (-0.361,-0.864) (-0.289,-0.870) (-0.218,-0.873) (-0.146,-0.874) (-0.076,-0.871) (-0.006,-0.866) (0.066,-0.857) (0.137,-0.846) (0.206,-0.832) (0.276,-0.814) (0.344,-0.793) (0.413,-0.768) (0.478,-0.741) (0.543,-0.710) (0.606,-0.674) (0.666,-0.635) (0.722,-0.593) (0.776,-0.546) (0.825,-0.496) (0.870,-0.441) (0.908,-0.383) (0.941,-0.319) (0.967,-0.252) (0.983,-0.183) (0.990,-0.112) (0.987,-0.041) (0.973,0.029) (0.951,0.096) (0.920,0.159) (0.881,0.220) (0.835,0.275) (0.785,0.324) (0.729,0.370) (0.670,0.410) (0.607,0.445) (0.543,0.475) (0.476,0.501) (0.407,0.522) (0.337,0.538) (0.266,0.550) (0.194,0.558) (0.123,0.560) (0.052,0.559) (-0.019,0.553) (-0.089,0.542) (-0.158,0.527) (-0.225,0.507) (-0.291,0.482) (-0.355,0.451) (-0.417,0.414) (-0.474,0.371) (-0.525,0.322) (-0.569,0.266) (-0.604,0.204) (-0.626,0.138) (-0.635,0.067) (-0.628,-0.003) (-0.606,-0.070) (-0.572,-0.132) (-0.526,-0.186) (-0.473,-0.233) (-0.413,-0.273) (-0.351,-0.304) (-0.284,-0.328) (-0.216,-0.345) (-0.146,-0.356) (-0.075,-0.359) (-0.004,-0.356) (0.065,-0.346) (0.134,-0.329) (0.200,-0.303) (0.262,-0.269) (0.318,-0.226) (0.364,-0.172) (0.396,-0.109) (0.407,-0.040) (0.395,0.030) (0.361,0.092) (0.312,0.142) (0.253,0.181) (0.188,0.208) (0.119,0.225) (0.049,0.230) (-0.021,0.226) (-0.090,0.209) (-0.155,0.181) (-0.211,0.138) (-0.250,0.080) (-0.260,0.011) (-0.235,-0.055) (-0.184,-0.104) (-0.120,-0.134) (-0.051,-0.147) (0.019,-0.144) (0.086,-0.123) (0.142,-0.081) (0.167,-0.015) (0.139,0.049) (0.079,0.085) (0.009,0.095) (-0.059,0.077) (-0.106,0.025) (-0.081,-0.038)};
\draw[ch25arrow] (axis cs:-0.136,-0.873) -- (axis cs:-0.128,-0.873);
\draw[ch25arrow] (axis cs:0.984,-0.166) -- (axis cs:0.985,-0.158);
\draw[ch25arrow] (axis cs:-0.193,0.517) -- (axis cs:-0.201,0.514);
\addplot[ch25orbit] coordinates {(1.040,0.642) (0.979,0.677) (0.916,0.710) (0.851,0.739) (0.784,0.766) (0.716,0.790) (0.647,0.810) (0.576,0.828) (0.505,0.843) (0.433,0.855) (0.361,0.864) (0.289,0.870) (0.218,0.873) (0.146,0.874) (0.076,0.871) (0.006,0.866) (-0.066,0.857) (-0.137,0.846) (-0.206,0.832) (-0.276,0.814) (-0.344,0.793) (-0.413,0.768) (-0.478,0.741) (-0.543,0.710) (-0.606,0.674) (-0.666,0.635) (-0.722,0.593) (-0.776,0.546) (-0.825,0.496) (-0.870,0.441) (-0.908,0.383) (-0.941,0.319) (-0.967,0.252) (-0.983,0.183) (-0.990,0.112) (-0.987,0.041) (-0.973,-0.029) (-0.951,-0.096) (-0.920,-0.159) (-0.881,-0.220) (-0.835,-0.275) (-0.785,-0.324) (-0.729,-0.370) (-0.670,-0.410) (-0.607,-0.445) (-0.543,-0.475) (-0.476,-0.501) (-0.407,-0.522) (-0.337,-0.538) (-0.266,-0.550) (-0.194,-0.558) (-0.123,-0.560) (-0.052,-0.559) (0.019,-0.553) (0.089,-0.542) (0.158,-0.527) (0.225,-0.507) (0.291,-0.482) (0.355,-0.451) (0.417,-0.414) (0.474,-0.371) (0.525,-0.322) (0.569,-0.266) (0.604,-0.204) (0.626,-0.138) (0.635,-0.067) (0.628,0.003) (0.606,0.070) (0.572,0.132) (0.526,0.186) (0.473,0.233) (0.413,0.273) (0.351,0.304) (0.284,0.328) (0.216,0.345) (0.146,0.356) (0.075,0.359) (0.004,0.356) (-0.065,0.346) (-0.134,0.329) (-0.200,0.303) (-0.262,0.269) (-0.318,0.226) (-0.364,0.172) (-0.396,0.109) (-0.407,0.040) (-0.395,-0.030) (-0.361,-0.092) (-0.312,-0.142) (-0.253,-0.181) (-0.188,-0.208) (-0.119,-0.225) (-0.049,-0.230) (0.021,-0.226) (0.090,-0.209) (0.155,-0.181) (0.211,-0.138) (0.250,-0.080) (0.260,-0.011) (0.235,0.055) (0.184,0.104) (0.120,0.134) (0.051,0.147) (-0.019,0.144) (-0.086,0.123) (-0.142,0.081) (-0.167,0.015) (-0.139,-0.049) (-0.079,-0.085) (-0.009,-0.095) (0.059,-0.077) (0.106,-0.025) (0.081,0.038)};
\draw[ch25arrow] (axis cs:0.136,0.873) -- (axis cs:0.128,0.873);
\draw[ch25arrow] (axis cs:-0.984,0.166) -- (axis cs:-0.985,0.158);
\draw[ch25arrow] (axis cs:0.193,-0.517) -- (axis cs:0.201,-0.514);
\fill[black] (axis cs:0,0) circle (1.3pt);
\end{axis}
\end{tikzpicture}}
\sdpanelrow
\sdpanel{D}{\begin{tikzpicture}
\begin{axis}[sdplot, width=0.31\textwidth, height=0.31\textwidth, xmin=-1, xmax=1, ymin=-1, ymax=1, clip=true, xtick={-1,0,1}, ytick={-1,0,1}, xlabel={$y_1$}, ylabel={$y_2$}]
\addplot[ch25orbit, domain=0:360, samples=41] ({0.300*cos(x)},{0.2121*sin(x)});
\draw[ch25arrow] (axis cs:0.010,0.2121) -- (axis cs:-0.000,0.2121);
\addplot[ch25orbit, domain=0:360, samples=41] ({0.550*cos(x)},{0.3889*sin(x)});
\draw[ch25arrow] (axis cs:0.010,0.3889) -- (axis cs:-0.000,0.3889);
\addplot[ch25orbit, domain=0:360, samples=41] ({0.800*cos(x)},{0.5657*sin(x)});
\draw[ch25arrow] (axis cs:0.010,0.5657) -- (axis cs:-0.000,0.5657);
\addplot[ch25orbit, domain=0:360, samples=41] ({0.980*cos(x)},{0.6930*sin(x)});
\draw[ch25arrow] (axis cs:0.010,0.6930) -- (axis cs:-0.000,0.6930);
\fill[black] (axis cs:0,0) circle (1.3pt);
\end{axis}
\end{tikzpicture}}\hfill
\sdpanel{E}{\begin{tikzpicture}
\begin{axis}[sdplot, width=0.31\textwidth, height=0.31\textwidth, xmin=-1, xmax=1, ymin=-1, ymax=1, clip=true, xtick={-1,0,1}, ytick={-1,0,1}, xlabel={$y_1$}, ylabel={$y_2$}]
\addplot[ch25orbit] coordinates {(0,0) (1.061,1.061)};
\draw[ch25arrow] (axis cs:0.557,0.557) -- (axis cs:0.550,0.550);
\addplot[ch25orbit] coordinates {(0,0) (-1.061,-1.061)};
\draw[ch25arrow] (axis cs:-0.557,-0.557) -- (axis cs:-0.550,-0.550);
\addplot[ch25orbit] coordinates {(1.043,0.397) (0.981,0.358) (0.918,0.320) (0.856,0.283) (0.795,0.247) (0.733,0.210) (0.672,0.175) (0.609,0.139) (0.547,0.105) (0.485,0.070) (0.422,0.037) (0.360,0.005) (0.296,-0.026) (0.232,-0.056) (0.168,-0.083) (0.102,-0.108) (0.034,-0.125) (-0.036,-0.120) (-0.034,-0.050) (-0.018,-0.024)};
\draw[ch25arrow] (axis cs:0.599,0.134) -- (axis cs:0.592,0.130);
\addplot[ch25orbit] coordinates {(1.042,0.055) (0.978,0.023) (0.914,-0.009) (0.850,-0.041) (0.785,-0.072) (0.722,-0.102) (0.656,-0.132) (0.592,-0.161) (0.527,-0.189) (0.461,-0.217) (0.393,-0.243) (0.327,-0.268) (0.260,-0.290) (0.192,-0.311) (0.124,-0.328) (0.054,-0.340) (-0.016,-0.343) (-0.084,-0.327) (-0.125,-0.270) (-0.118,-0.200) (-0.091,-0.135) (-0.055,-0.075) (-0.019,-0.023)};
\draw[ch25arrow] (axis cs:0.488,-0.205) -- (axis cs:0.480,-0.209);
\addplot[ch25orbit] coordinates {(1.041,-0.332) (0.975,-0.360) (0.907,-0.389) (0.838,-0.417) (0.773,-0.442) (0.707,-0.467) (0.641,-0.492) (0.573,-0.516) (0.505,-0.539) (0.437,-0.560) (0.368,-0.580) (0.300,-0.598) (0.232,-0.614) (0.163,-0.628) (0.092,-0.638) (0.021,-0.643) (-0.050,-0.641) (-0.119,-0.628) (-0.182,-0.597) (-0.224,-0.541) (-0.237,-0.472) (-0.228,-0.402) (-0.207,-0.335) (-0.179,-0.271) (-0.146,-0.209) (-0.110,-0.149) (-0.070,-0.091) (-0.029,-0.035) (-0.019,-0.023)};
\draw[ch25arrow] (axis cs:0.323,-0.592) -- (axis cs:0.316,-0.594);
\addplot[ch25orbit] coordinates {(-1.043,-0.397) (-0.981,-0.358) (-0.918,-0.320) (-0.856,-0.283) (-0.795,-0.247) (-0.733,-0.210) (-0.672,-0.175) (-0.609,-0.139) (-0.547,-0.105) (-0.485,-0.070) (-0.422,-0.037) (-0.360,-0.005) (-0.296,0.026) (-0.232,0.056) (-0.168,0.083) (-0.102,0.108) (-0.034,0.125) (0.036,0.120) (0.034,0.050) (0.018,0.024)};
\draw[ch25arrow] (axis cs:-0.599,-0.134) -- (axis cs:-0.592,-0.130);
\addplot[ch25orbit] coordinates {(-1.042,-0.055) (-0.978,-0.023) (-0.914,0.009) (-0.850,0.041) (-0.785,0.072) (-0.722,0.102) (-0.656,0.132) (-0.592,0.161) (-0.527,0.189) (-0.461,0.217) (-0.393,0.243) (-0.327,0.268) (-0.260,0.290) (-0.192,0.311) (-0.124,0.328) (-0.054,0.340) (0.016,0.343) (0.084,0.327) (0.125,0.270) (0.118,0.200) (0.091,0.135) (0.055,0.075) (0.019,0.023)};
\draw[ch25arrow] (axis cs:-0.488,0.205) -- (axis cs:-0.480,0.209);
\addplot[ch25orbit] coordinates {(-1.041,0.332) (-0.975,0.360) (-0.907,0.389) (-0.838,0.417) (-0.773,0.442) (-0.707,0.467) (-0.641,0.492) (-0.573,0.516) (-0.505,0.539) (-0.437,0.560) (-0.368,0.580) (-0.300,0.598) (-0.232,0.614) (-0.163,0.628) (-0.092,0.638) (-0.021,0.643) (0.050,0.641) (0.119,0.628) (0.182,0.597) (0.224,0.541) (0.237,0.472) (0.228,0.402) (0.207,0.335) (0.179,0.271) (0.146,0.209) (0.110,0.149) (0.070,0.091) (0.029,0.035) (0.019,0.023)};
\draw[ch25arrow] (axis cs:-0.323,0.592) -- (axis cs:-0.316,0.594);
\fill[black] (axis cs:0,0) circle (1.3pt);
\end{axis}
\end{tikzpicture}}\hfill
\sdpanel{F}{\begin{tikzpicture}
\begin{axis}[sdplot, width=0.31\textwidth, height=0.31\textwidth, xmin=-1, xmax=1, ymin=-1, ymax=1, clip=true, xtick={-1,0,1}, ytick={-1,0,1}, xlabel={$y_1$}, ylabel={$y_2$}]
\addplot[ch25orbit] coordinates {(0,0) (1.000,0.203)};
\draw[ch25arrow] (axis cs:0.560,0.114) -- (axis cs:0.550,0.111);
\addplot[ch25orbit] coordinates {(0,0) (1.000,0.883)};
\draw[ch25arrow] (axis cs:0.560,0.495) -- (axis cs:0.550,0.486);
\addplot[ch25orbit] coordinates {(0,0) (0.335,1.000)};
\draw[ch25arrow] (axis cs:0.188,0.560) -- (axis cs:0.184,0.550);
\addplot[ch25orbit] coordinates {(0,0) (-0.203,1.000)};
\draw[ch25arrow] (axis cs:-0.114,0.560) -- (axis cs:-0.111,0.550);
\addplot[ch25orbit] coordinates {(0,0) (-0.883,1.000)};
\draw[ch25arrow] (axis cs:-0.495,0.560) -- (axis cs:-0.486,0.550);
\addplot[ch25orbit] coordinates {(0,0) (-1.000,0.335)};
\draw[ch25arrow] (axis cs:-0.560,0.188) -- (axis cs:-0.550,0.184);
\addplot[ch25orbit] coordinates {(0,0) (-1.000,-0.203)};
\draw[ch25arrow] (axis cs:-0.560,-0.114) -- (axis cs:-0.550,-0.111);
\addplot[ch25orbit] coordinates {(0,0) (-1.000,-0.883)};
\draw[ch25arrow] (axis cs:-0.560,-0.495) -- (axis cs:-0.550,-0.486);
\addplot[ch25orbit] coordinates {(0,0) (-0.335,-1.000)};
\draw[ch25arrow] (axis cs:-0.188,-0.560) -- (axis cs:-0.184,-0.550);
\addplot[ch25orbit] coordinates {(0,0) (0.203,-1.000)};
\draw[ch25arrow] (axis cs:0.114,-0.560) -- (axis cs:0.111,-0.550);
\addplot[ch25orbit] coordinates {(0,0) (0.883,-1.000)};
\draw[ch25arrow] (axis cs:0.495,-0.560) -- (axis cs:0.486,-0.550);
\addplot[ch25orbit] coordinates {(0,0) (1.000,-0.335)};
\draw[ch25arrow] (axis cs:0.560,-0.188) -- (axis cs:0.550,-0.184);
\fill[black] (axis cs:0,0) circle (1.3pt);
\end{axis}
\end{tikzpicture}}
\caption{Phase portraits of $\mathbf y'=A\mathbf y$, with the equilibrium at the origin and arrows in the direction of increasing $t$. (A) Saddle, $A$ with the rows $(0,1)$ and $(2,1)$: eigenvalues $2$ and $-1$ with the eigenvectors $(1,2)$ and $(1,-1)$, drawn as straight lines. (B) Stable node, rows $(-1,0)$ and $(1,-3)$: eigenvalues $-1$ (slow direction $(2,1)$) and $-3$ (fast direction $\mathbf e_2$). (C) Stable spiral, rows $(-0.2,-2)$ and $(1,-0.2)$: eigenvalues $-0.2\pm i\sqrt2$. (D) Centre, rows $(0,-2)$ and $(1,0)$: eigenvalues $\pm i\sqrt2$, orbits the ellipses $y_1^2+2y_2^2=c$. (E) Degenerate node, rows $(-2,1)$ and $(-1,0)$: double eigenvalue $-1$ with the single eigenvector $(1,1)$. (F) Star, $A=-I$. For the unstable node, spiral, degenerate node and star, reverse the arrows.\label{fig:phase-portraits}}
\end{figure}

In a spiral or a centre the solutions turn about the origin, and the sense of rotation can be read off from a single entry of the matrix.

\begin{proposition}[Sense of rotation]
\label{prop:rotation-sense}
Let $A\in\mathbb R^{2\times2}$ have nonreal eigenvalues. Then $a_{21}\ne0$, and along every solution $\mathbf y\ne\mathbf 0$ of $\mathbf y'=A\mathbf y$ the quantity $y_1y_2'-y_2y_1'$ has the sign of $a_{21}$ at every time. Consequently every such solution turns counterclockwise about the origin if $a_{21}>0$ and clockwise if $a_{21}<0$.
\end{proposition}

\begin{proof}
If $a_{21}=0$, the matrix is triangular with the real eigenvalues $a_{11}$ and $a_{22}$ (\cref{prop:triangular-eigen}). Substituting $\mathbf y'=A\mathbf y$,
\[
y_1y_2'-y_2y_1'=y_1(a_{21}y_1+a_{22}y_2)-y_2(a_{11}y_1+a_{12}y_2)=a_{21}y_1^2+(a_{22}-a_{11})y_1y_2-a_{12}y_2^2 ,
\]
a quadratic form in $(y_1,y_2)$ whose discriminant is $(a_{22}-a_{11})^2+4a_{12}a_{21}=(a_{11}+a_{22})^2-4(a_{11}a_{22}-a_{12}a_{21})$, the discriminant of the characteristic polynomial, which is negative because the eigenvalues are not real. A quadratic form $\alpha y_1^2+\beta y_1y_2+\gamma y_2^2$ with $\beta^2<4\alpha\gamma$ is definite, with the sign of $\alpha$, by completing the square, so the quantity has the sign of $a_{21}$ wherever $\mathbf y\ne\mathbf 0$. Where $y_1\ne0$, the polar angle $\theta=\arctan(y_2/y_1)+\text{const}$ of the point $\mathbf y(t)$ has the derivative $\theta'=(y_1y_2'-y_2y_1')/(y_1^2+y_2^2)$ by the chain rule, and where $y_2\ne0$ the same follows from $\theta=-\arctan(y_1/y_2)+\text{const}$; so the angle increases if $a_{21}>0$ and decreases if $a_{21}<0$.
\end{proof}

At the point $\mathbf e_1$ the velocity is the first column $(a_{11},a_{21})$ of $A$, so the proposition says that the orbits cross the positive $y_1$-axis upwards when $a_{21}>0$; the argument shows that the rotation never reverses. In \cref{fig:phase-portraits}C and D, $a_{21}=1$ and the rotation is counterclockwise. The degenerate node in E has real eigenvalues, so the proposition does not apply, but the same computation gives $y_1y_2'-y_2y_1'=-(y_1-y_2)^2\le0$: its orbits, which turn by half a revolution in all, turn clockwise, and the only orbits that do not turn are those on the eigenvector line $y_1=y_2$, which the other orbits never meet.

When $\det A=0$ the origin is not the only equilibrium. If $A\ne0$ has the eigenvalue $0$ and a second eigenvalue $\tau=\operatorname{tr}A\ne0$, the kernel of $A$ is a line of equilibria, and every other solution, $\mathbf y(t)=c_1\mathbf v_1+c_2e^{\tau t}\mathbf v_2$, runs along a line parallel to the eigenvector $\mathbf v_2$, towards the line of equilibria if $\tau<0$ and away from it if $\tau>0$. If $A\ne0$ is nilpotent, it is similar to $J_2(0)$, whose solutions $(c_1+c_2t,c_2)$ move with constant velocity along lines parallel to the kernel, in opposite directions on its two sides. If $A=0$, nothing moves.

\subsection{The Trace--Determinant Plane}

The type of the equilibrium depends only on the eigenvalues, hence only on the trace and the determinant, and these two numbers draw a map of all linear systems in the plane.

\begin{theorem}[Trace and determinant]
\label{thm:trace-determinant}
Let $A\in\mathbb R^{2\times2}$, $\tau\coloneqq\operatorname{tr}A$, $\delta\coloneqq\det A$ and $\Delta\coloneqq\tau^2-4\delta$. The equilibrium $\mathbf 0$ of $\mathbf y'=A\mathbf y$ is
\begin{enumerate}
\item a saddle if $\delta<0$;
\item a node if $\delta>0$ and $\Delta>0$;
\item a star if $A=(\tau/2)I$, and a degenerate node otherwise, if $\delta>0$ and $\Delta=0$;
\item a spiral if $\delta>0$, $\Delta<0$ and $\tau\ne0$, and a centre if $\delta>0$ and $\tau=0$.
\end{enumerate}
In the cases (ii) to (iv), the solutions tend to $\mathbf 0$ as $t\to\infty$ if $\tau<0$ and as $t\to-\infty$ if $\tau>0$.
\end{theorem}

\begin{proof}
By \eqref{eq:eigen-2x2}, the eigenvalues are the roots $(\tau\pm\sqrt\Delta)/2$ of $\lambda^2-\tau\lambda+\delta$, with the sum $\tau$ and the product $\delta$ (\cref{thm:trace-det-eigen}). If $\delta<0$ they are real, since $\Delta>\tau^2\ge0$, and of opposite signs. If $\delta>0$ and $\Delta>0$ they are real, different and of the same sign, the sign of their sum $\tau$. If $\delta>0$ and $\Delta=0$ there is a double eigenvalue $\tau/2\ne0$. If $\Delta<0$ they are $\tau/2\pm i\sqrt{-\Delta}/2$, with the real part $\tau/2$; for $\tau=0$ the condition $\Delta<0$ is $\delta>0$. In each case \cref{thm:phase-classification} gives the type and the direction of time in which the solutions approach $\mathbf 0$.
\end{proof}

\Cref{fig:trace-determinant} shows the map. The parabola $\tau^2=4\delta$ separates nodes from spirals, the positive $\delta$-axis carries the centres, and the negative half-plane $\delta<0$ the saddles; \cref{tab:phase-types} summarises the classification. The types that fill open regions of the map, saddles, nodes and spirals, are robust: a small change of the entries of $A$ moves the point $(\tau,\delta)$ only slightly and does not change the type. Centres, stars and degenerate nodes lie on curves, and an arbitrarily small change can turn a centre into a spiral, inwards or outwards, or a degenerate node into a node or a spiral. This distinction becomes important in \cref{sec:nonlinear-systems}, where the linear system is only an approximation.

\begin{figure}
\sdfloat{%
\begin{tikzpicture}
\begin{axis}[sddiagram, width=0.6\textwidth, height=0.46\textwidth, xmin=-4.2, xmax=4.2, ymin=-2.2, ymax=4.6,
  xtick=\empty, ytick=\empty, xlabel={$\tau=\operatorname{tr}A$}, ylabel={$\delta=\det A$}]
\addplot[black, dashed, line width=0.7pt, domain=-4.1:4.1, samples=60] {x^2/4};
\addplot[black, line width=1.4pt, domain=0:4.6, samples=2] ({0},{x});
\node[font=\scriptsize, align=center] at (axis cs:-1.35,3.0) {stable\\spirals};
\node[font=\scriptsize, align=center] at (axis cs:1.35,3.0) {unstable\\spirals};
\node[font=\scriptsize, align=center] at (axis cs:-3.2,0.75) {stable\\nodes};
\node[font=\scriptsize, align=center] at (axis cs:3.2,0.75) {unstable\\nodes};
\node[font=\scriptsize] at (axis cs:1.5,-1.3) {saddles};
\node[font=\scriptsize, anchor=west] at (axis cs:0.1,4.3) {centres};
\end{axis}
\end{tikzpicture}}{The trace--determinant plane (\cref{thm:trace-determinant}). The parabola $\tau^2=4\delta$ (dashed) carries the stars and degenerate nodes and separates nodes from spirals; the positive $\delta$-axis (thick) carries the centres; on the $\tau$-axis, $\delta=0$, lie the matrices with a line of equilibria.\label{fig:trace-determinant}}
\end{figure}

\begin{table}
\sdfloat{%
\begin{sdtable}{colspec={llll}, row{1}={font=\small\bfseries}}
\toprule
type & eigenvalues & trace and determinant & solutions as $t\to\infty$ \\
\midrule
saddle & $\lambda_2<0<\lambda_1$ & $\delta<0$ & unbounded, except on a line \\
stable node & $\lambda_2<\lambda_1<0$ & $\tau<0$, $0<\delta<\tau^2/4$ & tend to $\mathbf 0$ \\
unstable node & $0<\lambda_2<\lambda_1$ & $\tau>0$, $0<\delta<\tau^2/4$ & unbounded \\
star, degenerate node & $\lambda_1=\lambda_2\ne0$ & $\delta=\tau^2/4>0$ & tend to $\mathbf 0$ iff $\tau<0$ \\
stable spiral & $a\pm ib$, $a<0$ & $\tau<0$, $\delta>\tau^2/4$ & tend to $\mathbf 0$ \\
unstable spiral & $a\pm ib$, $a>0$ & $\tau>0$, $\delta>\tau^2/4$ & unbounded \\
centre & $\pm ib$ & $\tau=0$, $\delta>0$ & periodic \\
\bottomrule
\end{sdtable}}{The equilibria of $\mathbf y'=A\mathbf y$ for real $2\times2$ matrices with $\det A\ne0$ (\cref{thm:phase-classification,thm:trace-determinant}); $\tau=\operatorname{tr}A$, $\delta=\det A$, $b>0$.\label{tab:phase-types}}
\end{table}

\subsection{The Damped Oscillator in the Phase Plane}

The oscillator $my''+cy'+ky=0$ of \cref{sec:oscillations}, with $m,k>0$ and $c\ge0$, becomes, with the velocity $v\coloneqq y'$ and in the notation $\omega_0=\sqrt{k/m}$, $\zeta=c/(2\sqrt{mk})$ of \eqref{eq:oscillator-normalised}, the system
\[
\begin{bmatrix}y\\v\end{bmatrix}'=\begin{bmatrix}0&1\\-\omega_0^2&-2\zeta\omega_0\end{bmatrix}\begin{bmatrix}y\\v\end{bmatrix},\qquad\tau=-2\zeta\omega_0,\quad\delta=\omega_0^2,\quad\Delta=4\omega_0^2(\zeta^2-1).
\]
Since $\delta>0$, the equilibrium is never a saddle, and \cref{thm:trace-determinant} translates \cref{prop:free-oscillations} into the language of the phase plane. Without damping, $\zeta=0$, it is a centre, and the orbits are the ellipses $my'^2+ky^2=2E$ of constant energy (\cref{prop:oscillator-energy}). An underdamped oscillator, $0<\zeta<1$, has $\Delta<0$ and $\tau<0$: a stable spiral, clockwise because $a_{21}=-\omega_0^2<0$ (\cref{prop:rotation-sense}), each turn being one full oscillation. A critically damped one, $\zeta=1$, has $\Delta=0$, and since the companion matrix has a single Jordan block for each eigenvalue (\cref{prop:nonderogatory}), the origin is a degenerate node, not a star. An overdamped one, $\zeta>1$, is a stable node, whose slow direction is the eigenvector $(1,\lambda_{\mathrm s})$ for the eigenvalue $\lambda_{\mathrm s}=\omega_0(-\zeta+\sqrt{\zeta^2-1})$ of smaller modulus: after the fast transient the body creeps back with the velocity $v\approx\lambda_{\mathrm s}y$, proportional to its displacement. \Cref{fig:damped-portraits} shows the three damped cases. In all of them the energy decreases along every solution, so the orbits cross the ellipses of constant energy inwards, as \cref{prop:oscillator-energy} predicts.

\begin{figure}
\sdpanel{A}{\begin{tikzpicture}
\begin{axis}[sdplot, width=0.31\textwidth, height=0.31\textwidth, xmin=-1, xmax=1, ymin=-1, ymax=1, clip=true, xtick={-1,0,1}, ytick={-1,0,1}, xlabel={$y$}, ylabel={$y'$}]
\addplot[ch25orbit] coordinates {(0.433,1.040) (0.488,0.993) (0.540,0.945) (0.590,0.894) (0.639,0.840) (0.684,0.785) (0.727,0.729) (0.767,0.669) (0.804,0.609) (0.839,0.546) (0.870,0.482) (0.898,0.416) (0.921,0.350) (0.942,0.282) (0.958,0.213) (0.970,0.143) (0.977,0.072) (0.980,0.000) (0.977,-0.071) (0.969,-0.142) (0.955,-0.212) (0.936,-0.280) (0.910,-0.345) (0.878,-0.407) (0.840,-0.468) (0.795,-0.523) (0.746,-0.573) (0.691,-0.617) (0.631,-0.656) (0.568,-0.687) (0.501,-0.711) (0.432,-0.728) (0.361,-0.738) (0.290,-0.741) (0.219,-0.737) (0.149,-0.727) (0.080,-0.710) (0.013,-0.688) (-0.052,-0.660) (-0.115,-0.626) (-0.175,-0.588) (-0.231,-0.545) (-0.283,-0.499) (-0.332,-0.448) (-0.376,-0.392) (-0.416,-0.334) (-0.449,-0.272) (-0.477,-0.207) (-0.498,-0.140) (-0.511,-0.071) (-0.516,-0.000) (-0.511,0.070) (-0.496,0.138) (-0.469,0.204) (-0.430,0.263) (-0.380,0.313) (-0.321,0.351) (-0.256,0.377) (-0.187,0.389) (-0.116,0.388) (-0.047,0.375) (0.019,0.351) (0.080,0.318) (0.136,0.275) (0.185,0.224) (0.225,0.166) (0.254,0.102) (0.270,0.034) (0.269,-0.037) (0.248,-0.104) (0.206,-0.160) (0.145,-0.195) (0.076,-0.205) (0.007,-0.192) (-0.055,-0.158) (-0.104,-0.109) (-0.136,-0.046) (-0.141,0.024) (-0.107,0.085) (-0.041,0.108) (0.026,0.086) (0.053,0.060)};
\draw[ch25arrow] (axis cs:0.799,0.618) -- (axis cs:0.803,0.611);
\draw[ch25arrow] (axis cs:0.844,-0.461) -- (axis cs:0.840,-0.467);
\draw[ch25arrow] (axis cs:-0.375,-0.394) -- (axis cs:-0.380,-0.387);
\addplot[ch25orbit] coordinates {(-0.433,-1.040) (-0.488,-0.993) (-0.540,-0.945) (-0.590,-0.894) (-0.639,-0.840) (-0.684,-0.785) (-0.727,-0.729) (-0.767,-0.669) (-0.804,-0.609) (-0.839,-0.546) (-0.870,-0.482) (-0.898,-0.416) (-0.921,-0.350) (-0.942,-0.282) (-0.958,-0.213) (-0.970,-0.143) (-0.977,-0.072) (-0.980,0.000) (-0.977,0.071) (-0.969,0.142) (-0.955,0.212) (-0.936,0.280) (-0.910,0.345) (-0.878,0.407) (-0.840,0.468) (-0.795,0.523) (-0.746,0.573) (-0.691,0.617) (-0.631,0.656) (-0.568,0.687) (-0.501,0.711) (-0.432,0.728) (-0.361,0.738) (-0.290,0.741) (-0.219,0.737) (-0.149,0.727) (-0.080,0.710) (-0.013,0.688) (0.052,0.660) (0.115,0.626) (0.175,0.588) (0.231,0.545) (0.283,0.499) (0.332,0.448) (0.376,0.392) (0.416,0.334) (0.449,0.272) (0.477,0.207) (0.498,0.140) (0.511,0.071) (0.516,0.000) (0.511,-0.070) (0.496,-0.138) (0.469,-0.204) (0.430,-0.263) (0.380,-0.313) (0.321,-0.351) (0.256,-0.377) (0.187,-0.389) (0.116,-0.388) (0.047,-0.375) (-0.019,-0.351) (-0.080,-0.318) (-0.136,-0.275) (-0.185,-0.224) (-0.225,-0.166) (-0.254,-0.102) (-0.270,-0.034) (-0.269,0.037) (-0.248,0.104) (-0.206,0.160) (-0.145,0.195) (-0.076,0.205) (-0.007,0.192) (0.055,0.158) (0.104,0.109) (0.136,0.046) (0.141,-0.024) (0.107,-0.085) (0.041,-0.108) (-0.026,-0.086) (-0.053,-0.060)};
\draw[ch25arrow] (axis cs:-0.799,-0.618) -- (axis cs:-0.803,-0.611);
\draw[ch25arrow] (axis cs:-0.844,0.461) -- (axis cs:-0.840,0.467);
\draw[ch25arrow] (axis cs:0.375,0.394) -- (axis cs:0.380,0.387);
\fill[black] (axis cs:0,0) circle (1.3pt);
\end{axis}
\end{tikzpicture}}\hfill
\sdpanel{B}{\begin{tikzpicture}
\begin{axis}[sdplot, width=0.31\textwidth, height=0.31\textwidth, xmin=-1, xmax=1, ymin=-1, ymax=1, clip=true, xtick={-1,0,1}, ytick={-1,0,1}, xlabel={$y$}, ylabel={$y'$}]
\addplot[ch25orbit] coordinates {(0,0) (1.061,-1.061)};
\draw[ch25arrow] (axis cs:0.557,-0.557) -- (axis cs:0.550,-0.550);
\addplot[ch25orbit] coordinates {(0,0) (-1.061,1.061)};
\draw[ch25arrow] (axis cs:-0.557,0.557) -- (axis cs:-0.550,0.550);
\addplot[ch25orbit] coordinates {(-0.246,1.041) (-0.210,0.977) (-0.174,0.912) (-0.139,0.848) (-0.105,0.785) (-0.072,0.722) (-0.038,0.659) (-0.006,0.595) (0.026,0.531) (0.057,0.467) (0.086,0.403) (0.115,0.337) (0.142,0.272) (0.166,0.207) (0.188,0.140) (0.206,0.072) (0.214,0.002) (0.195,-0.066) (0.125,-0.073) (0.062,-0.044) (0.024,-0.019)};
\draw[ch25arrow] (axis cs:0.014,0.555) -- (axis cs:0.017,0.548);
\addplot[ch25orbit] coordinates {(0.116,1.044) (0.146,0.980) (0.177,0.914) (0.207,0.848) (0.236,0.781) (0.264,0.715) (0.292,0.649) (0.319,0.581) (0.344,0.516) (0.368,0.450) (0.391,0.382) (0.413,0.314) (0.432,0.245) (0.449,0.176) (0.462,0.107) (0.470,0.037) (0.470,-0.034) (0.455,-0.103) (0.411,-0.157) (0.343,-0.174) (0.274,-0.161) (0.208,-0.136) (0.146,-0.103) (0.087,-0.066) (0.030,-0.025) (0.023,-0.019)};
\draw[ch25arrow] (axis cs:0.378,0.420) -- (axis cs:0.381,0.412);
\addplot[ch25orbit] coordinates {(0.486,1.041) (0.512,0.976) (0.538,0.908) (0.564,0.839) (0.590,0.770) (0.614,0.703) (0.636,0.636) (0.658,0.570) (0.679,0.500) (0.699,0.432) (0.717,0.363) (0.734,0.294) (0.748,0.223) (0.760,0.153) (0.768,0.083) (0.772,0.013) (0.770,-0.058) (0.758,-0.127) (0.733,-0.193) (0.688,-0.247) (0.624,-0.277) (0.554,-0.284) (0.485,-0.274) (0.418,-0.254) (0.353,-0.228) (0.289,-0.197) (0.228,-0.162) (0.169,-0.125) (0.111,-0.086) (0.054,-0.044) (0.023,-0.019)};
\draw[ch25arrow] (axis cs:0.742,0.253) -- (axis cs:0.744,0.245);
\addplot[ch25orbit] coordinates {(0.246,-1.041) (0.210,-0.977) (0.174,-0.912) (0.139,-0.848) (0.105,-0.785) (0.072,-0.722) (0.038,-0.659) (0.006,-0.595) (-0.026,-0.531) (-0.057,-0.467) (-0.086,-0.403) (-0.115,-0.337) (-0.142,-0.272) (-0.166,-0.207) (-0.188,-0.140) (-0.206,-0.072) (-0.214,-0.002) (-0.195,0.066) (-0.125,0.073) (-0.062,0.044) (-0.024,0.019)};
\draw[ch25arrow] (axis cs:-0.014,-0.555) -- (axis cs:-0.017,-0.548);
\addplot[ch25orbit] coordinates {(-0.116,-1.044) (-0.146,-0.980) (-0.177,-0.914) (-0.207,-0.848) (-0.236,-0.781) (-0.264,-0.715) (-0.292,-0.649) (-0.319,-0.581) (-0.344,-0.516) (-0.368,-0.450) (-0.391,-0.382) (-0.413,-0.314) (-0.432,-0.245) (-0.449,-0.176) (-0.462,-0.107) (-0.470,-0.037) (-0.470,0.034) (-0.455,0.103) (-0.411,0.157) (-0.343,0.174) (-0.274,0.161) (-0.208,0.136) (-0.146,0.103) (-0.087,0.066) (-0.030,0.025) (-0.023,0.019)};
\draw[ch25arrow] (axis cs:-0.378,-0.420) -- (axis cs:-0.381,-0.412);
\addplot[ch25orbit] coordinates {(-0.486,-1.041) (-0.512,-0.976) (-0.538,-0.908) (-0.564,-0.839) (-0.590,-0.770) (-0.614,-0.703) (-0.636,-0.636) (-0.658,-0.570) (-0.679,-0.500) (-0.699,-0.432) (-0.717,-0.363) (-0.734,-0.294) (-0.748,-0.223) (-0.760,-0.153) (-0.768,-0.083) (-0.772,-0.013) (-0.770,0.058) (-0.758,0.127) (-0.733,0.193) (-0.688,0.247) (-0.624,0.277) (-0.554,0.284) (-0.485,0.274) (-0.418,0.254) (-0.353,0.228) (-0.289,0.197) (-0.228,0.162) (-0.169,0.125) (-0.111,0.086) (-0.054,0.044) (-0.023,0.019)};
\draw[ch25arrow] (axis cs:-0.742,-0.253) -- (axis cs:-0.744,-0.245);
\fill[black] (axis cs:0,0) circle (1.3pt);
\end{axis}
\end{tikzpicture}}\hfill
\sdpanel{C}{\begin{tikzpicture}
\begin{axis}[sdplot, width=0.31\textwidth, height=0.31\textwidth, xmin=-1, xmax=1, ymin=-1, ymax=1, clip=true, xtick={-1,0,1}, ytick={-1,0,1}, xlabel={$y$}, ylabel={$y'$}]
\addplot[ch25orbit] coordinates {(0,0) (1.449,-0.388)};
\draw[ch25arrow] (axis cs:0.560,-0.150) -- (axis cs:0.550,-0.147);
\addplot[ch25orbit] coordinates {(0,0) (-1.449,0.388)};
\draw[ch25arrow] (axis cs:-0.560,0.150) -- (axis cs:-0.550,0.147);
\addplot[ch25orbit] coordinates {(0,0) (0.388,-1.449)};
\draw[ch25arrow] (axis cs:0.150,-0.560) -- (axis cs:0.147,-0.550);
\addplot[ch25orbit] coordinates {(0,0) (-0.388,1.449)};
\draw[ch25arrow] (axis cs:-0.150,0.560) -- (axis cs:-0.147,0.550);
\addplot[ch25orbit] coordinates {(0.798,1.042) (0.814,0.965) (0.829,0.893) (0.844,0.816) (0.858,0.745) (0.872,0.671) (0.886,0.596) (0.899,0.527) (0.911,0.452) (0.923,0.380) (0.933,0.311) (0.943,0.238) (0.951,0.168) (0.957,0.097) (0.961,0.025) (0.960,-0.046) (0.951,-0.116) (0.921,-0.179) (0.857,-0.209) (0.787,-0.205) (0.719,-0.191) (0.650,-0.174) (0.583,-0.156) (0.515,-0.138) (0.447,-0.120) (0.379,-0.102) (0.312,-0.083) (0.244,-0.065) (0.176,-0.047) (0.109,-0.029) (0.041,-0.011) (0.029,-0.008)};
\draw[ch25arrow] (axis cs:0.948,0.197) -- (axis cs:0.949,0.189);
\addplot[ch25orbit] coordinates {(0.550,1.044) (0.566,0.970) (0.581,0.901) (0.596,0.828) (0.612,0.753) (0.627,0.683) (0.641,0.612) (0.655,0.541) (0.668,0.472) (0.681,0.399) (0.693,0.327) (0.704,0.256) (0.714,0.186) (0.722,0.115) (0.727,0.044) (0.728,-0.027) (0.718,-0.096) (0.676,-0.152) (0.606,-0.157) (0.537,-0.143) (0.469,-0.126) (0.401,-0.108) (0.334,-0.089) (0.266,-0.071) (0.198,-0.053) (0.131,-0.035) (0.063,-0.017) (0.029,-0.008)};
\draw[ch25arrow] (axis cs:0.696,0.309) -- (axis cs:0.697,0.301);
\addplot[ch25orbit] coordinates {(0.053,1.043) (0.071,0.970) (0.088,0.901) (0.105,0.830) (0.123,0.757) (0.140,0.685) (0.157,0.614) (0.173,0.544) (0.189,0.473) (0.205,0.403) (0.221,0.333) (0.235,0.262) (0.249,0.193) (0.261,0.124) (0.271,0.054) (0.274,-0.016) (0.218,-0.057) (0.150,-0.040) (0.082,-0.022) (0.029,-0.008)};
\draw[ch25arrow] (axis cs:0.177,0.527) -- (axis cs:0.179,0.519);
\addplot[ch25orbit] coordinates {(-0.751,1.046) (-0.729,0.973) (-0.707,0.900) (-0.684,0.829) (-0.662,0.759) (-0.640,0.691) (-0.618,0.623) (-0.595,0.557) (-0.572,0.490) (-0.548,0.422) (-0.522,0.355) (-0.495,0.290) (-0.465,0.226) (-0.429,0.166) (-0.381,0.115) (-0.317,0.086) (-0.250,0.067) (-0.182,0.049) (-0.114,0.031) (-0.047,0.012) (-0.029,0.008)};
\draw[ch25arrow] (axis cs:-0.586,0.528) -- (axis cs:-0.583,0.521);
\addplot[ch25orbit] coordinates {(-0.797,-1.049) (-0.813,-0.971) (-0.828,-0.899) (-0.842,-0.830) (-0.856,-0.757) (-0.870,-0.683) (-0.883,-0.613) (-0.896,-0.543) (-0.908,-0.473) (-0.919,-0.403) (-0.930,-0.332) (-0.941,-0.260) (-0.949,-0.190) (-0.956,-0.118) (-0.961,-0.046) (-0.961,0.026) (-0.955,0.096) (-0.933,0.163) (-0.878,0.206) (-0.808,0.208) (-0.738,0.196) (-0.670,0.179) (-0.602,0.161) (-0.534,0.143) (-0.467,0.125) (-0.399,0.107) (-0.331,0.089) (-0.263,0.071) (-0.196,0.052) (-0.128,0.034) (-0.061,0.016) (-0.029,0.008)};
\draw[ch25arrow] (axis cs:-0.948,-0.201) -- (axis cs:-0.949,-0.193);
\addplot[ch25orbit] coordinates {(-0.551,-1.042) (-0.567,-0.969) (-0.582,-0.900) (-0.598,-0.827) (-0.613,-0.752) (-0.628,-0.682) (-0.642,-0.611) (-0.656,-0.541) (-0.669,-0.471) (-0.682,-0.399) (-0.694,-0.326) (-0.705,-0.255) (-0.715,-0.185) (-0.723,-0.114) (-0.728,-0.044) (-0.729,0.027) (-0.719,0.096) (-0.677,0.152) (-0.607,0.157) (-0.538,0.143) (-0.470,0.126) (-0.402,0.108) (-0.334,0.090) (-0.267,0.071) (-0.199,0.053) (-0.131,0.035) (-0.064,0.017) (-0.029,0.008)};
\draw[ch25arrow] (axis cs:-0.697,-0.308) -- (axis cs:-0.698,-0.300);
\addplot[ch25orbit] coordinates {(-0.055,-1.043) (-0.073,-0.970) (-0.090,-0.901) (-0.107,-0.830) (-0.124,-0.757) (-0.142,-0.685) (-0.159,-0.614) (-0.175,-0.544) (-0.191,-0.473) (-0.207,-0.403) (-0.222,-0.332) (-0.237,-0.262) (-0.251,-0.193) (-0.263,-0.124) (-0.273,-0.054) (-0.276,0.016) (-0.219,0.058) (-0.151,0.040) (-0.084,0.022) (-0.029,0.008)};
\draw[ch25arrow] (axis cs:-0.179,-0.526) -- (axis cs:-0.181,-0.518);
\addplot[ch25orbit] coordinates {(0.748,-1.040) (0.726,-0.968) (0.704,-0.896) (0.681,-0.825) (0.659,-0.755) (0.638,-0.688) (0.615,-0.620) (0.592,-0.551) (0.568,-0.483) (0.544,-0.416) (0.519,-0.350) (0.492,-0.285) (0.462,-0.222) (0.425,-0.161) (0.375,-0.112) (0.311,-0.084) (0.243,-0.065) (0.176,-0.047) (0.108,-0.029) (0.040,-0.011) (0.029,-0.008)};
\draw[ch25arrow] (axis cs:0.583,-0.526) -- (axis cs:0.581,-0.518);
\fill[black] (axis cs:0,0) circle (1.3pt);
\end{axis}
\end{tikzpicture}}
\caption{The damped oscillator $y''+2\zeta y'+y=0$ in the phase plane of position $y$ and velocity $y'$. (A) Underdamped, $\zeta=0.2$: a stable spiral, turning clockwise. (B) Critically damped, $\zeta=1$: a degenerate node with the single eigenvector line $y'=-y$. (C) Overdamped, $\zeta=2$: a stable node with the slow direction $y'=(-2+\sqrt3)y$ and the fast direction $y'=(-2-\sqrt3)y$. The time courses of the same motions, started at $(1,0)$, are shown in \cref{fig:damping}.\label{fig:damped-portraits}}
\end{figure}

\section{Stability of Linear Systems}
\label{sec:linear-stability}

\subsection{Stability of an Equilibrium}

\Cref{def:stability} defined stability for an equilibrium of a single autonomous equation. The definition carries over to systems with the Euclidean norm in place of the absolute value.

\begin{definition}[Stability of an equilibrium]
\label{def:stability-systems}
Let $U\subseteq\mathbb R^n$ be open, $\mathbf f\colon U\to\mathbb R^n$ continuously differentiable, and $\mathbf f(\mathbf y^*)=\mathbf 0$. The equilibrium $\mathbf y^*$ of $\mathbf y'=\mathbf f(\mathbf y)$ is \term{stable} if for every $\varepsilon>0$ there is $\delta>0$ such that for $\|\mathbf y_0-\mathbf y^*\|<\delta$ the solution with $\mathbf y(0)=\mathbf y_0$ exists for all $t\ge0$ and satisfies $\|\mathbf y(t)-\mathbf y^*\|<\varepsilon$ for all $t\ge0$. It is \term{asymptotically stable} if it is stable and there is $\delta>0$ such that $\mathbf y(t)\to\mathbf y^*$ as $t\to\infty$ whenever $\|\mathbf y_0-\mathbf y^*\|<\delta$. It is \term{unstable} if it is not stable. The same definitions apply to complex systems $\mathbf y'=A\mathbf y$ with $A\in\mathbb C^{n\times n}$ and the equilibrium $\mathbf 0$.
\end{definition}

By \cref{lem:orbits}\,(i), starting at the time $0$ is no restriction for an autonomous system. Stability says that solutions starting close to the equilibrium stay close; asymptotic stability, that they moreover return to it. A centre is stable but not asymptotically stable: the solutions circle around it at a constant distance, in the sense of the ellipses through their initial points. For linear systems, local and global behaviour coincide, because solutions can be scaled, and stability becomes a property of the matrix $e^{tA}$.

\begin{proposition}[Stability and the exponential]
\label{prop:stability-exponential}
Let $A\in\mathbb F^{n\times n}$. The equilibrium $\mathbf 0$ of $\mathbf y'=A\mathbf y$ is stable if and only if $\|e^{tA}\|_2$ is bounded for $t\ge0$, and asymptotically stable if and only if $e^{tA}\to0$ as $t\to\infty$; in the second case every solution tends to $\mathbf 0$.
\end{proposition}

\begin{proof}
If $\|e^{tA}\|_2\le M$ for $t\ge0$, then $\|\mathbf y(t)\|=\|e^{tA}\mathbf y_0\|\le M\|\mathbf y_0\|$, and $\delta\coloneqq\varepsilon/M$ works. Conversely, if $\mathbf 0$ is stable, choose $\delta$ for $\varepsilon=1$; for a unit vector $\mathbf x$, the solution starting at $\delta\mathbf x/2$ satisfies $\|e^{tA}\delta\mathbf x/2\|<1$, so $\|e^{tA}\|_2\le2/\delta$. If $e^{tA}\to0$, then $\|e^{tA}\|_2\le\|e^{tA}\|_{\mathrm F}\to0$, so the norm is bounded, $\mathbf 0$ is stable, and every solution $e^{tA}\mathbf y_0$ tends to $\mathbf 0$. Conversely, if the solutions starting in the ball of radius $\delta$ tend to $\mathbf 0$, then so do the columns $e^{tA}\mathbf e_j=(2/\delta)\,e^{tA}(\delta\mathbf e_j/2)$, and $e^{tA}\to0$.
\end{proof}

\begin{theorem}[Stability of linear systems]
\label{thm:linear-stability}
Let $A\in\mathbb C^{n\times n}$, and let $\alpha(A)\coloneqq\max\{\operatorname{Re}\lambda:\lambda\text{ an eigenvalue of }A\}$ be its spectral abscissa.
\begin{enumerate}
\item $e^{tA}\to0$ as $t\to\infty$, and $\mathbf 0$ is asymptotically stable, if and only if $\alpha(A)<0$, that is, every eigenvalue has negative real part.
\item $e^{tA}$ is bounded for $t\ge0$, and $\mathbf 0$ is stable, if and only if $\alpha(A)\le0$ and every eigenvalue with real part $0$ is semisimple.
\item If an eigenvalue has positive real part, or an eigenvalue with real part $0$ is not semisimple, then $\mathbf y'=A\mathbf y$ has unbounded solutions, and $\mathbf 0$ is unstable.
\end{enumerate}
\end{theorem}

\begin{proof}
Let $A=SJS^{-1}$ with a Jordan matrix $J$ (\cref{thm:jordan-form}). Then $e^{tA}=Se^{tJ}S^{-1}$, so $e^{tA}$ tends to zero or is bounded if and only if $e^{tJ}=S^{-1}e^{tA}S$ is, and $e^{tJ}$ is block diagonal with the blocks of \cref{prop:exp-jordan}, whose entries are $0$ or $e^{\mu t}t^j/j!$, of modulus $e^{(\operatorname{Re}\mu)t}t^j/j!$, for the eigenvalues $\mu$ and $j$ less than the size of the block.

If $\alpha(A)<0$, every entry tends to zero, since $t^je^{-ct}\to0$ for $c>0$ (\cref{prop:exp-growth}). If $\alpha(A)\le0$ and every eigenvalue on the imaginary axis is semisimple, its blocks have size $1$ and are the numbers $e^{\mu t}$ of modulus $1$, while the other entries tend to zero; so $e^{tJ}$ is bounded. This proves the implications from right to left in (i) and (ii).

For the converses and (iii), let $\lambda$ be an eigenvalue with an eigenvector $\mathbf v$. The solution $e^{tA}\mathbf v=e^{\lambda t}\mathbf v$ has the norm $e^{(\operatorname{Re}\lambda)t}\|\mathbf v\|$, which does not tend to zero if $\operatorname{Re}\lambda\ge0$ and is unbounded if $\operatorname{Re}\lambda>0$. If $\operatorname{Re}\lambda=0$ and $\lambda$ is not semisimple, it has a Jordan chain $\mathbf v_1,\mathbf v_2$ of length $2$, and by \cref{prop:chain-solutions} the function $e^{\lambda t}(\mathbf v_2+t\mathbf v_1)$ is a solution, of norm at least $t\|\mathbf v_1\|-\|\mathbf v_2\|$, which is unbounded. A bounded $e^{tA}$ admits no unbounded solution, and the stability statements follow from \cref{prop:stability-exponential}.
\end{proof}

The theorem is the continuous counterpart of \cref{thm:powers-size} for the powers of a matrix, with the left half-plane $\operatorname{Re}\lambda<0$ in the place of the unit disc $|\lambda|<1$. The two are connected by the exponential: the matrix $e^{hA}$, which advances the solutions by the time $h$, has the eigenvalues $e^{h\lambda}$ (\cref{prop:exp-basic}), of modulus $e^{h\operatorname{Re}\lambda}$, which is less than $1$ exactly when $\operatorname{Re}\lambda<0$; \cref{ch:difference-equations} studies the discrete systems $\mathbf x_{k+1}=A\mathbf x_k$ from this point of view. For systems of two equations, \cref{tab:phase-types} lists the stability of each type: stable nodes, spirals and degenerate nodes are asymptotically stable, centres are stable but not asymptotically stable, and saddles and the unstable types are unstable. The two masses of \cref{sec:ode-systems}, written as a system of four equations, have the eigenvalues $\pm i\omega_1$ and $\pm i\omega_2$ (\cref{sec:normal-modes}), all semisimple, so their equilibrium is stable but not asymptotically stable: without friction, the bodies vibrate for ever. The system $y_1'=y_2$, $y_2'=0$, the equation $y''=0$ of a free particle, has the matrix $J_2(0)$, whose double eigenvalue $0$ is not semisimple; its solutions $y_1=c_1+c_2t$ are unbounded, and the particle drifts away from any initial position with its initial velocity.

\subsection{Rates of Decay and Transient Growth}

When every eigenvalue has negative real part, the solutions decay exponentially, at a rate set by the spectral abscissa, but the theorem says nothing about how large they may become first.

\begin{proposition}[Exponential decay]
\label{prop:decay-estimate}
Let $A\in\mathbb C^{n\times n}$ and $\sigma\in\mathbb R$.
\begin{enumerate}
\item If $\alpha(A)<-\sigma$, there is a constant $C\ge1$ with $\|e^{tA}\|_2\le Ce^{-\sigma t}$ for all $t\ge0$. If every eigenvalue with real part $\alpha(A)$ is semisimple, this holds also for $\sigma=-\alpha(A)$.
\item Conversely, if $\|e^{tA}\|_2\le Ce^{-\sigma t}$ for all $t\ge0$ with some constant $C$, then $\alpha(A)\le-\sigma$.
\end{enumerate}
\end{proposition}

\begin{proof}
(i) With $A=SJS^{-1}$ as before, the entries of $e^{\sigma t}e^{tJ}$ have the moduli $e^{(\operatorname{Re}\mu+\sigma)t}t^j/j!$ or $0$. If $\operatorname{Re}\mu+\sigma<0$, such a function is continuous on $[0,\infty)$ and tends to zero, so it is bounded; under the additional hypothesis, the blocks with $\operatorname{Re}\mu+\sigma=0$ have size $1$ and their entries have modulus $1$. So $\|e^{\sigma t}e^{tJ}\|_2\le\|e^{\sigma t}e^{tJ}\|_{\mathrm F}\le C'$ for $t\ge0$, and $\|e^{tA}\|_2\le\|S\|_2\|e^{tJ}\|_2\|S^{-1}\|_2\le\kappa_2(S)C'e^{-\sigma t}$. Since $e^{0A}=I$, any such constant is at least $1$. (ii) By \cref{prop:exp-norm-bounds}, $e^{t\alpha(A)}\le\|e^{tA}\|_2\le Ce^{-\sigma t}$, so $e^{t(\alpha(A)+\sigma)}\le C$ for all $t\ge0$, which forces $\alpha(A)+\sigma\le0$.
\end{proof}

The constant $C$ cannot be dispensed with, and it can be large. For
\[
A_s=\begin{bmatrix}-1&s\\0&-1\end{bmatrix},\qquad e^{tA_s}=e^{-t}\begin{bmatrix}1&st\\0&1\end{bmatrix}
\]
by \cref{prop:exp-2x2}\,(ii), both eigenvalues are $-1$, and the entry $ste^{-t}$ reaches the value $s/e$ at $t=1$, by \cref{prop:exp-growth} and elementary calculus. The singular values of a matrix with the rows $(1,c)$ and $(0,1)$ are $(\sqrt{c^2+4}\pm c)/2$, the square roots of the eigenvalues of its Gram matrix, which has the trace $2+c^2$ and the determinant $1$ (\cref{prop:singular-eigen}); so
\[
\|e^{tA_s}\|_2=e^{-t}\,\frac{st+\sqrt{s^2t^2+4}}2 .
\]
For $s=100$ this rises from $1$ to its maximum $36.79$, attained at $t\approx1.00$, before it decays; for $s=10$ the maximum is $3.716$ at $t\approx0.98$. A disturbance in the second component is first amplified nearly $37$-fold by a system whose eigenvalues both promise decay at the rate $1$. The example of \cref{fig:transient}, with the eigenvalues $-1$ and $-2$, showed the same phenomenon for a diagonalisable matrix. \Cref{prop:exp-norm-bounds} identifies the matrices for which it cannot occur.

\begin{corollary}[Monotone decay]
\label{cor:monotone-decay}
Let $A\in\mathbb C^{n\times n}$. The norm $\|\mathbf y(t)\|$ of every solution of $\mathbf y'=A\mathbf y$ is nonincreasing in $t$ if and only if the Hermitian matrix $A+A^{\mathsf H}$ is negative semidefinite. If $A$ is normal, $\|e^{tA}\|_2=e^{t\alpha(A)}$, and the constant $C=1$ works in \cref{prop:decay-estimate}\,(i) for $\sigma=-\alpha(A)$.
\end{corollary}

\begin{proof}
The condition says that the numerical abscissa $\omega(A)$, the largest eigenvalue of $(A+A^{\mathsf H})/2$, is at most $0$. If so, $\|e^{tA}\mathbf y_0\|\le e^{t\omega(A)}\|\mathbf y_0\|\le\|\mathbf y_0\|$ by \cref{prop:exp-norm-bounds}, and since $\mathbf y(t+s)=e^{tA}\mathbf y(s)$, the norm does not increase. If $\omega(A)>0$, let $\mathbf y_0$ be a unit eigenvector of $(A+A^{\mathsf H})/2$ for its largest eigenvalue $\omega(A)$. As in the proof of \cref{prop:exp-norm-bounds}, the function $\varphi(t)\coloneqq\|e^{tA}\mathbf y_0\|^2$ has the derivative $\varphi'(0)=\mathbf y_0^{\mathsf H}(A^{\mathsf H}+A)\mathbf y_0=2\omega(A)>0$, so $\varphi(t)>\varphi(0)=1$ for small $t>0$, and the norm of this solution increases. For a normal $A$ the statement is \cref{prop:exp-norm-bounds} once more, since then $\omega(A)=\alpha(A)$.
\end{proof}

For a symmetric matrix, in particular, the solutions of a stable system decay monotonically in norm, as \cref{sec:eigenvalue-method} found. For nonnormal matrices the spectral abscissa governs the long-time behaviour and the numerical abscissa the short-time behaviour, and between the two lies the transient, which can be large; Trefethen and Embree~\cite{trefethen2005} describe how such transients arise in applications, for example in the flow of fluids, where a linearised system that is stable by its eigenvalues amplifies small disturbances by large factors.

\subsection{Criteria from the Coefficients}

The eigenvalues are the roots of the characteristic polynomial, and the question whether they all lie in the left half-plane can be decided from the coefficients, without computing the roots.

\begin{theorem}[Stability of polynomials of degree $2$ and $3$]
\label{thm:hurwitz-small}
Let the polynomials below have real coefficients.
\begin{enumerate}
\item All roots of $\lambda^2+a_1\lambda+a_0$ have negative real parts if and only if $a_1>0$ and $a_0>0$.
\item All roots of $\lambda^3+a_2\lambda^2+a_1\lambda+a_0$ have negative real parts if and only if $a_2>0$, $a_0>0$ and $a_2a_1>a_0$.
\end{enumerate}
\end{theorem}

\begin{proof}
(i) If the roots $r_1$, $r_2$ are real, then $a_1=-(r_1+r_2)$ and $a_0=r_1r_2$, and both roots are negative if and only if their sum is negative and their product positive. If they are $a\pm ib$ with $b\ne0$, then $a_1=-2a$ and $a_0=a^2+b^2>0$, so the real part is negative if and only if $a_1>0$.

(ii) Let $p(\lambda)\coloneqq\lambda^3+a_2\lambda^2+a_1\lambda+a_0$. A polynomial of odd degree with real coefficients has a real root (\cref{cor:ivt}), say $-r$, and then $p(\lambda)=(\lambda+r)(\lambda^2+s\lambda+q)$ with real $s$ and $q$, by \cref{cor:factor}. Comparing coefficients gives $a_2=r+s$, $a_1=q+rs$ and $a_0=rq$, and hence the identity
\[
a_2a_1-a_0=(r+s)(q+rs)-rq=s\,(q+r^2+rs)=s\,(q+ra_2).
\]
Suppose that all roots have negative real parts. Then $r>0$, and by (i) $s>0$ and $q>0$; so $a_2=r+s>0$, $a_0=rq>0$ and $a_2a_1-a_0=s(q+ra_2)>0$. Conversely, suppose that $a_2>0$, $a_0>0$ and $a_2a_1>a_0$; then also $a_1>0$. For $\lambda\ge0$ all terms of $p(\lambda)$ are nonnegative and $a_0>0$, so the real root $-r$ is negative: $r>0$. Then $q=a_0/r>0$, so $q+ra_2>0$, and the identity gives $s=(a_2a_1-a_0)/(q+ra_2)>0$. By (i), the roots of $\lambda^2+s\lambda+q$ have negative real parts, and so do all roots of $p$.
\end{proof}

By \cref{thm:charpoly-coefficients}, the characteristic polynomial of a $3\times3$ matrix is $\lambda^3-\operatorname{tr}A\,\lambda^2+E_2(A)\,\lambda-\det A$, where $E_2(A)$ is the sum of the three principal minors of order $2$. So $\mathbf y'=A\mathbf y$ with a real $3\times3$ matrix is asymptotically stable if and only if
\[
\operatorname{tr}A<0,\qquad\det A<0,\qquad\operatorname{tr}A\cdot E_2(A)<\det A ;
\]
for $2\times2$ matrices, (i) gives the condition $\operatorname{tr}A<0$, $\det A>0$ of \cref{tab:phase-types}. The polynomial $\lambda^3+\lambda^2+\lambda+c$ has all its roots in the left half-plane exactly for $0<c<1$; at $c=1$ it factors as $(\lambda+1)(\lambda^2+1)$, and two roots sit on the imaginary axis, where the condition $a_2a_1>a_0$ turns into an equality. For polynomials of every degree, the answer was given by Routh and by Hurwitz~\cite{routh1877,hurwitz1895}.

\begin{theorem}[Routh--Hurwitz criterion]
\label{thm:routh-hurwitz}
Let $p(\lambda)=a_0\lambda^n+a_1\lambda^{n-1}+\cdots+a_n$ with real coefficients and $a_0>0$, let $a_k\coloneqq0$ for $k<0$ and $k>n$, and let $H$ be the $n\times n$ \term{Hurwitz matrix} with the entries $h_{ij}=a_{2j-i}$,
\[
H=\begin{bmatrix}a_1&a_3&a_5&\cdots\\a_0&a_2&a_4&\cdots\\0&a_1&a_3&\cdots\\0&a_0&a_2&\cdots\\\vdots&&&\ddots\end{bmatrix}.
\]
All roots of $p$ have negative real parts if and only if the leading principal minors $\Delta_1,\ldots,\Delta_n$ of $H$ are all positive.
\end{theorem}

\begin{proof}[sketch]
For $n=3$ and $a_0=1$, the minors are $\Delta_1=a_1$, $\Delta_2=a_1a_2-a_3$ and $\Delta_3=a_3\Delta_2$, and the theorem is \cref{thm:hurwitz-small}\,(ii) in the other numbering of the coefficients. In general, write $p(\lambda)=a_0\prod_j(\lambda-\lambda_j)$ and let $\omega$ run along the real axis from $-\infty$ to $\infty$. The point $i\omega$ then moves up the imaginary axis, and the argument of each factor $i\omega-\lambda_j$ increases by $\pi$ if $\operatorname{Re}\lambda_j<0$ and decreases by $\pi$ if $\operatorname{Re}\lambda_j>0$. So all roots lie in the left half-plane if and only if $p$ has no root on the imaginary axis and the argument of $p(i\omega)$ increases by $n\pi$. Writing $p(i\omega)=U(\omega)+iV(\omega)$ with real polynomials $U$ and $V$, this happens exactly when the roots of $U$ and $V$ are real and interlace in a particular way, and Routh's algorithm, a Euclidean algorithm for the pair $U$, $V$ in the spirit of Sturm sequences, tests the interlacing through the signs of the numbers it produces. Hurwitz expressed these numbers as the minors $\Delta_k$. Gantmacher~\cite[Chap.~XV]{gantmacher1959vol2} gives the complete proof.
\end{proof}

\subsection{Lyapunov's Equation}

\Cref{thm:linear-stability} decides stability through the eigenvalues. Lyapunov found a second way, which needs no eigenvalues and which extends to nonlinear systems: a function that decreases along every solution, the higher-dimensional analogue of the energy of a damped oscillator. For linear systems the natural candidates are the quadratic functions $V(\mathbf y)=\mathbf y^{\mathsf H}P\mathbf y$ with a positive definite $P$, whose level sets are ellipsoids. Along a solution of $\mathbf y'=A\mathbf y$,
\[
\frac d{dt}\,\mathbf y^{\mathsf H}P\mathbf y=\mathbf y'^{\mathsf H}P\mathbf y+\mathbf y^{\mathsf H}P\mathbf y'=\mathbf y^{\mathsf H}\bigl(A^{\mathsf H}P+PA\bigr)\mathbf y ,
\]
so $V$ decreases along every solution if $A^{\mathsf H}P+PA$ is negative definite. The question is whether such a $P$ exists.

\begin{theorem}[Lyapunov's theorem]
\label{thm:lyapunov-equation}
Let $A\in\mathbb C^{n\times n}$.
\begin{enumerate}
\item If $\alpha(A)<0$, then for every Hermitian $Q\in\mathbb C^{n\times n}$ the \term{Lyapunov equation}
\begin{equation}
A^{\mathsf H}P+PA=-Q
\label{eq:lyapunov}
\end{equation}
has exactly one solution $P$. It is Hermitian, it is real if $A$ and $Q$ are real, and
\[
P=\int_0^\infty e^{tA^{\mathsf H}}Qe^{tA}\,dt\coloneqq\lim_{T\to\infty}\int_0^Te^{tA^{\mathsf H}}Qe^{tA}\,dt .
\]
If $Q$ is positive definite, so is $P$.
\item Conversely, if \eqref{eq:lyapunov} holds with Hermitian positive definite $P$ and $Q$, then $\alpha(A)<0$, and every solution of $\mathbf y'=A\mathbf y$ satisfies
\[
\|\mathbf y(t)\|\le\sqrt{\kappa_2(P)}\,e^{-\gamma t}\|\mathbf y(0)\|\qquad(t\ge0),\qquad\gamma\coloneqq\frac{\lambda_1(Q)}{2\lambda_n(P)} .
\]
\end{enumerate}
\end{theorem}

\begin{proof}
(i) The equation is a Sylvester equation \eqref{eq:sylvester} with the matrices $A^{\mathsf H}$ and $A$ in the places of $A$ and $B$. The eigenvalues of $A^{\mathsf H}$ are the conjugates $\bar\lambda_i$ of those of $A$, and $\bar\lambda_i=-\lambda_j$ would mean $\operatorname{Re}\lambda_i+\operatorname{Re}\lambda_j=0$, which is impossible when all real parts are negative. So $A^{\mathsf H}$ and $-A$ have no common eigenvalue, and by \cref{thm:sylvester-equation} the linear map $\mathcal L(X)\coloneqq A^{\mathsf H}X+XA$ on $\mathbb C^{n\times n}$ is invertible, and \eqref{eq:lyapunov} has exactly one solution $P$. Taking the conjugate transpose of \eqref{eq:lyapunov} gives $P^{\mathsf H}A+A^{\mathsf H}P^{\mathsf H}=-Q^{\mathsf H}=-Q$, so $P^{\mathsf H}$ is a solution as well, and $P^{\mathsf H}=P$. If $A$ and $Q$ are real, conjugating \eqref{eq:lyapunov} shows that $\bar P$ is a solution as well, so $\bar P=P$ and $P$ is real.

For the integral, let $F(t)\coloneqq e^{tA^{\mathsf H}}Qe^{tA}$ and $P_T\coloneqq\int_0^TF$. Since $(e^{tA})^{\mathsf H}=e^{tA^{\mathsf H}}$ (\cref{prop:exp-basic}), the product rule and \cref{thm:exp-derivative} give $F'=A^{\mathsf H}F+FA$, so, integrating entry by entry,
\[
\mathcal L(P_T)=A^{\mathsf H}P_T+P_TA=\int_0^TF'=F(T)-Q .
\]
Subtracting $\mathcal L(P)=-Q$ gives $\mathcal L(P_T-P)=F(T)$, that is, $P_T-P=\mathcal L^{-1}(F(T))$. By \cref{thm:linear-stability}, $F(T)\to0$ as $T\to\infty$, and since the entries of $\mathcal L^{-1}(X)$ are fixed linear combinations of the entries of $X$, $P_T\to P$ entry by entry. If $Q$ is positive definite and $\mathbf x\ne\mathbf 0$, the function $g(t)\coloneqq\mathbf x^{\mathsf H}F(t)\mathbf x=(e^{tA}\mathbf x)^{\mathsf H}Q(e^{tA}\mathbf x)$ is continuous, nonnegative, and positive at $t=0$, so $\mathbf x^{\mathsf H}P_T\mathbf x=\int_0^Tg$ is nondecreasing in $T$ and positive for $T>0$, and its limit $\mathbf x^{\mathsf H}P\mathbf x$ is positive.

(ii) Let $\mathbf y$ be a solution and $V(t)\coloneqq\mathbf y(t)^{\mathsf H}P\mathbf y(t)$. By \cref{thm:rayleigh}, $\lambda_1(P)\|\mathbf y\|^2\le V\le\lambda_n(P)\|\mathbf y\|^2$, and by the computation before the theorem
\[
V'=-\mathbf y^{\mathsf H}Q\mathbf y\le-\lambda_1(Q)\|\mathbf y\|^2\le-\frac{\lambda_1(Q)}{\lambda_n(P)}\,V=-2\gamma V .
\]
So $(e^{2\gamma t}V)'\le0$, and $V(t)\le e^{-2\gamma t}V(0)$ (\cref{cor:monotone-derivative}). Hence $\lambda_1(P)\|\mathbf y(t)\|^2\le e^{-2\gamma t}\lambda_n(P)\|\mathbf y(0)\|^2$, which is the estimate, since $\kappa_2(P)=\lambda_n(P)/\lambda_1(P)$ for a positive definite $P$ (\cref{prop:singular-eigen}). All solutions tend to $\mathbf 0$, so $\alpha(A)<0$ by \cref{thm:linear-stability}.
\end{proof}

A function $V$ that is positive away from the equilibrium and decreases along the solutions is a \term{Lyapunov function}. The theorem says that a linear system is asymptotically stable if and only if it has a quadratic Lyapunov function, and that one is found by solving a linear system of equations for the entries of $P$, the Lyapunov equation, which is a Sylvester equation and can be solved by the method of Bartels and Stewart that \cref{sec:schur} described; the routine \code{linalg.solve\_sylvester} solves it, and Volume~VIII treats its algorithm. Its level sets are ellipsoids that every orbit crosses inwards. For a real symmetric $A$, the choice $P=I$ gives $A^{\mathsf T}P+PA=2A$, so $\|\mathbf y\|^2$ itself is a Lyapunov function when $A$ is negative definite, in accordance with \cref{cor:monotone-decay}.

For the matrix of \cref{fig:transient}, with the rows $(-1,10)$ and $(0,-2)$, and $Q=I$, the equation \eqref{eq:lyapunov} has the entries $-2p_{11}=-1$, $10p_{11}-3p_{12}=0$ and $20p_{12}-4p_{22}=-1$, so
\[
P=\begin{bmatrix}1/2&5/3\\5/3&103/12\end{bmatrix},\qquad\det P=\frac{109}{72},
\]
which is positive definite, since $p_{11}>0$ and $\det P>0$ (\cref{thm:pd-criteria}). Its eigenvalues are $0.1698$ and $8.913$, so $\kappa_2(P)=52.48$, and part (ii) guarantees $\|\mathbf y(t)\|\le7.24\,e^{-0.0561t}\|\mathbf y(0)\|$. The bound is valid but far from sharp: the largest value of $\|e^{tA}\|_2$ is $2.563$, and the solutions decay at the rate $1$. A different $Q$ gives a different ellipsoid and a different bound, and every choice certifies stability.

\section{Nonhomogeneous Systems}
\label{sec:nonhomogeneous-systems}

\subsection{Variation of Constants and Duhamel's Formula}

By \cref{thm:system-solution-space}\,(iii), the inhomogeneous system $\mathbf y'=A(t)\mathbf y+\mathbf b(t)$ is solved once one particular solution is known. As for a single equation (\cref{sec:variation-parameters}), it is found by letting the constants of the general homogeneous solution vary: with a fundamental matrix $Y$, try $\mathbf y=Y(t)\mathbf c(t)$. Then $\mathbf y'=AY\mathbf c+Y\mathbf c'$, and the system holds if and only if $Y\mathbf c'=\mathbf b$, that is, $\mathbf c'=Y^{-1}\mathbf b$, an equation that is solved by an integration.

\begin{theorem}[Variation of constants]
\label{thm:duhamel}
Let $A\colon I\to\mathbb F^{n\times n}$ and $\mathbf b\colon I\to\mathbb F^n$ be continuous, $t_0\in I$ and $\mathbf y_0\in\mathbb F^n$, and let $\Phi$ be the transition matrix of $\mathbf y'=A(t)\mathbf y$. The solution of $\mathbf y'=A(t)\mathbf y+\mathbf b(t)$ with $\mathbf y(t_0)=\mathbf y_0$ is
\begin{equation}
\mathbf y(t)=\Phi(t,t_0)\,\mathbf y_0+\int_{t_0}^t\Phi(t,s)\,\mathbf b(s)\,ds .
\label{eq:variation-constants}
\end{equation}
In particular, for a constant matrix $A$ we have \term{Duhamel's formula}
\begin{equation}
\mathbf y(t)=e^{(t-t_0)A}\mathbf y_0+\int_{t_0}^te^{(t-s)A}\,\mathbf b(s)\,ds .
\label{eq:duhamel-system}
\end{equation}
\end{theorem}

\begin{proof}
Let $Y$ be a fundamental matrix. Its inverse $Y(s)^{-1}=\operatorname{adj}Y(s)/\det Y(s)$ (\cref{thm:adjugate}) has continuous entries, since the determinant has no zero (\cref{prop:fundamental-matrices}), so $\mathbf z(t)\coloneqq\int_{t_0}^tY(s)^{-1}\mathbf b(s)\,ds$ is continuously differentiable with $\mathbf z'=Y^{-1}\mathbf b$, by the fundamental theorem of calculus applied to each entry. Since $\Phi(t,s)=Y(t)Y(s)^{-1}$ by \cref{prop:transition-matrix} and $Y(t)$ does not depend on the variable of integration, the right side of \eqref{eq:variation-constants} is
\[
\mathbf y(t)=Y(t)\bigl(Y(t_0)^{-1}\mathbf y_0+\mathbf z(t)\bigr).
\]
By the product rule, $\mathbf y'=Y'(Y(t_0)^{-1}\mathbf y_0+\mathbf z)+Y\mathbf z'=A\mathbf y+\mathbf b$, and $\mathbf y(t_0)=\mathbf y_0$. By \cref{thm:linear-existence}, this is the only solution. For a constant $A$, $\Phi(t,s)=e^{(t-s)A}$ by \cref{prop:transition-matrix}\,(iv).
\end{proof}

The formula has the same reading as for a single equation (\cref{sec:linear-first-order,sec:variation-parameters}). The first term is the free motion from the initial state. In the integral, the input $\mathbf b(s)\,ds$ received during the short interval at the time $s$ is a small change of state, which the homogeneous system then carries forward to the time $t$ by the transition matrix $\Phi(t,s)$; the solution is the sum of all these propagated contributions, by superposition. Since $e^{(t-s)A}=e^{tA}e^{-sA}$, Duhamel's formula can also be written as $e^{(t-t_0)A}\mathbf y_0+e^{tA}\int_{t_0}^te^{-sA}\mathbf b(s)\,ds$, the form in which it was announced in \cref{ch:canonical-forms}, and for $n=1$ and $A=-a$ it is the formula \eqref{eq:duhamel} for $y'+ay=q(t)$.

The theorem contains the variation of parameters of \cref{ch:linear-higher-order}. For the companion system \eqref{eq:first-order-system} of $Ly=g$, with the input $g(s)\mathbf e_k$, the $k$th column $\Phi(t,s)\mathbf e_k$ of the transition matrix is the solution with the value $\mathbf e_k$ at the time $s$, and by \cref{prop:companion-system} its first component is the solution of $Ly=0$ with the initial values $y(s)=\cdots=y^{(k-2)}(s)=0$ and $y^{(k-1)}(s)=1$, which is the Cauchy function $K(t,s)$ of \cref{thm:variation-general}. So the first component of \eqref{eq:variation-constants} with $\mathbf y_0=\mathbf 0$ is $\int_{t_0}^tK(t,s)g(s)\,ds$, the formula \eqref{eq:cauchy-function}, and for constant coefficients it is the convolution with the impulse response of \cref{cor:convolution}.

\subsection{Exponential Inputs}

For inputs that are exponentials times constant vectors, and therefore also for sines and cosines, a particular solution can be found without integration, as by undetermined coefficients in \cref{sec:undetermined}.

\begin{theorem}[Exponential inputs]
\label{thm:exponential-inputs}
Let $A\in\mathbb C^{n\times n}$, $\mu\in\mathbb C$ and $\mathbf c\in\mathbb C^n$, and consider $\mathbf y'=A\mathbf y+e^{\mu t}\mathbf c$.
\begin{enumerate}
\item If $\mu$ is not an eigenvalue of $A$, then $\mathbf y_{\mathrm p}(t)=e^{\mu t}(\mu I-A)^{-1}\mathbf c$ is a solution, and it is the only solution of the form $e^{\mu t}\mathbf v$ with a constant vector $\mathbf v$.
\item If $\mu$ is an eigenvalue of $A$ of index $m$, there is a solution $e^{\mu t}\mathbf q(t)$ with a vector polynomial $\mathbf q$ of degree at most $m$.
\end{enumerate}
\end{theorem}

\begin{proof}
(i) For $\mathbf y=e^{\mu t}\mathbf v$, the system reads $\mu e^{\mu t}\mathbf v=e^{\mu t}(A\mathbf v+\mathbf c)$, that is, $(\mu I-A)\mathbf v=\mathbf c$, which has exactly one solution $\mathbf v$ since $\mu I-A$ is invertible. (ii) Let $N\coloneqq A-\mu I$, let $G\coloneqq G_\mu(A)$ and let $R$ be the sum of the other generalised eigenspaces, so that $\mathbb C^n=G\oplus R$ and $\mathbf c=\mathbf c_G+\mathbf c_R$ with $\mathbf c_G\in G$ and $\mathbf c_R\in R$ (\cref{thm:generalised-eigenspaces}). The space $R$ is invariant under $N$, and $N$ is injective on it, since its kernel $E_\mu$ lies in $G$ and $G\cap R=\{\mathbf 0\}$; so $N$ maps $R$ onto $R$ (\cref{cor:injective-surjective}), and there is $\mathbf v_R\in R$ with $-N\mathbf v_R=\mathbf c_R$. Then $e^{\mu t}\mathbf v_R$ solves $\mathbf y'=A\mathbf y+e^{\mu t}\mathbf c_R$, by the computation in (i). On $G$, $N^m=0$ by \cref{thm:jordan-minimal}; let
\[
\mathbf q_G(t)\coloneqq\sum_{j=0}^{m-1}\frac{t^{j+1}}{(j+1)!}N^j\mathbf c_G .
\]
A function $e^{\mu t}\mathbf q(t)$ solves $\mathbf y'=A\mathbf y+e^{\mu t}\mathbf c_G$ if and only if $\mathbf q'=N\mathbf q+\mathbf c_G$, and for $\mathbf q_G$
\[
\mathbf q_G'-N\mathbf q_G=\sum_{j=0}^{m-1}\frac{t^j}{j!}N^j\mathbf c_G-\sum_{j=1}^{m}\frac{t^j}{j!}N^j\mathbf c_G=\mathbf c_G ,
\]
since $N^m\mathbf c_G=\mathbf 0$. By superposition, $e^{\mu t}(\mathbf q_G(t)+\mathbf v_R)$ is a solution, and its polynomial has degree at most $m$.
\end{proof}

A polynomial factor can be needed only when the input has a component in the generalised eigenspace of $\mu$, that is, when the input pushes the system at one of its own exponential rates; if $\mu$ is semisimple, it is needed whenever that component is not zero, since then $N=0$ on $G_\mu(A)$ and $(\mu I-A)\mathbf v=\mathbf c$ has no solution. This is \term{resonance}, as in \cref{thm:undetermined}, where the multiplicity of $\mu$ as a root of the characteristic polynomial played the role of the index. For the harmonic oscillator $y''+y=\cos t$, written as $\mathbf y'=A\mathbf y+\operatorname{Re}(e^{it}\mathbf e_2)$ with the rows $(0,1)$ and $(-1,0)$ of $A$, the number $\mu=i$ is a simple eigenvalue with the eigenvector $(1,i)$, and $-i$ has the eigenvector $(1,-i)$. Then $\mathbf e_2=\mathbf c_G+\mathbf c_R$ with $\mathbf c_G=(-i/2,1/2)$ and $\mathbf c_R=(i/2,1/2)$. On $R$, $N=A-iI$ acts as multiplication by $-2i$, so $\mathbf v_R=\mathbf c_R/(2i)=(1/4,-i/4)$, and $m=1$ gives $\mathbf q_G=t\mathbf c_G$. The first component of the solution $e^{it}(t\mathbf c_G+\mathbf v_R)$ is $e^{it}(1/4-it/2)$, and its real part
\[
\operatorname{Re}\Bigl((\cos t+i\sin t)\Bigl(\frac14-\frac{it}2\Bigr)\Bigr)=\frac{\cos t}4+\frac t2\sin t
\]
solves $y''+y=\cos t$. It is the resonant solution $t\sin(t)/2$ of \cref{sec:variation-parameters} plus the homogeneous solution $\cos(t)/4$.

\subsection{Bounded Inputs and Steady States}

A stable system responds to a bounded input with a bounded output, and to an input that settles down with an output that settles down.

\begin{theorem}[Bounded inputs]
\label{thm:bounded-input}
Let $A\in\mathbb C^{n\times n}$ with $\alpha(A)<0$, and let $C\ge1$ and $\sigma>0$ be constants with $\|e^{tA}\|_2\le Ce^{-\sigma t}$ for $t\ge0$ (\cref{prop:decay-estimate}). Let $\mathbf b$ be continuous on $[t_0,\infty)$ with $\|\mathbf b(t)\|\le\beta$. Then every solution of $\mathbf y'=A\mathbf y+\mathbf b(t)$ satisfies
\[
\|\mathbf y(t)\|\le Ce^{-\sigma(t-t_0)}\|\mathbf y(t_0)\|+\frac{C\beta}\sigma\qquad(t\ge t_0).
\]
If moreover $\mathbf b(t)\to\mathbf b_\infty$ as $t\to\infty$, then $\mathbf y(t)\to-A^{-1}\mathbf b_\infty$.
\end{theorem}

\begin{proof}
By Duhamel's formula and \cref{lem:integral-norm},
\[
\begin{aligned}
\|\mathbf y(t)\|&\le\|e^{(t-t_0)A}\|_2\|\mathbf y(t_0)\|+\int_{t_0}^t\|e^{(t-s)A}\|_2\|\mathbf b(s)\|\,ds\\
&\le Ce^{-\sigma(t-t_0)}\|\mathbf y(t_0)\|+C\beta\int_{t_0}^te^{-\sigma(t-s)}\,ds ,
\end{aligned}
\]
and the last integral is $(1-e^{-\sigma(t-t_0)})/\sigma\le1/\sigma$. For the second statement, $A$ is invertible, since $0$ is not an eigenvalue, and $\mathbf w\coloneqq\mathbf y+A^{-1}\mathbf b_\infty$ satisfies $\mathbf w'=A\mathbf w+(\mathbf b(t)-\mathbf b_\infty)$. Let $\varepsilon>0$ and choose $T\ge t_0$ with $\|\mathbf b(s)-\mathbf b_\infty\|\le\varepsilon$ for $s\ge T$, and let $\beta'$ be the maximum of $\|\mathbf b-\mathbf b_\infty\|$ on $[t_0,T]$. Splitting the integral in Duhamel's formula for $\mathbf w$ at $T$ gives, for $t\ge T$,
\[
\|\mathbf w(t)\|\le Ce^{-\sigma(t-t_0)}\|\mathbf w(t_0)\|+C\beta'(T-t_0)e^{-\sigma(t-T)}+\frac{C\varepsilon}\sigma ,
\]
whose first two terms tend to zero as $t\to\infty$; so $\|\mathbf w(t)\|<2C\varepsilon/\sigma$ for large $t$, and $\mathbf w(t)\to\mathbf 0$.
\end{proof}

The limit $-A^{-1}\mathbf b_\infty$ is the \term{steady state} of the constant input $\mathbf b_\infty$, the equilibrium of $\mathbf y'=A\mathbf y+\mathbf b_\infty$; the theorem says that every solution approaches it, whatever its initial value. If the fresh water flowing into the first tank of \cref{sec:ode-systems} carries salt at the concentration $\gamma$, the system becomes $\mathbf y'=A\mathbf y+(\gamma,0)$ with the matrix \eqref{eq:tank-matrix}, whose inverse has the rows $(-1,-1/2)$ and $(-1,-1)$, and the steady state is $-A^{-1}(\gamma,0)=(\gamma,\gamma)$: in the end both tanks, of unit volume, have the concentration of the inflow, as they must. The hypothesis $\alpha(A)<0$ cannot be dropped. For the free particle $\mathbf y'=J_2(0)\mathbf y+\mathbf b$ with a constant $\mathbf b=(b_1,b_2)$, $b_2\ne0$, the solutions grow like $t^2$, and for the undamped oscillator the bounded input $\cos t$ produced the unbounded solution $t\sin(t)/2$.

\subsection{Periodic Inputs and the Frequency Response}

\begin{theorem}[Periodic inputs]
\label{thm:periodic-response}
Let $A\in\mathbb C^{n\times n}$, $\mathbf c\in\mathbb C^n$, and let $\omega\in\mathbb R$ be such that $i\omega$ is not an eigenvalue of $A$. Then
\[
\mathbf y_{\mathrm p}(t)=e^{i\omega t}H(i\omega)\,\mathbf c,\qquad H(z)\coloneqq(zI-A)^{-1},
\]
solves $\mathbf y'=A\mathbf y+e^{i\omega t}\mathbf c$. If $\alpha(A)<0$, every solution $\mathbf y$ satisfies $\|\mathbf y(t)-\mathbf y_{\mathrm p}(t)\|\le Ce^{-\sigma(t-t_0)}\|\mathbf y(t_0)-\mathbf y_{\mathrm p}(t_0)\|$ with the constants of \cref{prop:decay-estimate}, and for $\omega\ne0$, $\mathbf y_{\mathrm p}$ is the only solution that is periodic with the period $2\pi/|\omega|$. If $A$ and $\mathbf c$ are real, $\operatorname{Re}\mathbf y_{\mathrm p}$ solves $\mathbf y'=A\mathbf y+\cos(\omega t)\,\mathbf c$.
\end{theorem}

\begin{proof}
The first statement is \cref{thm:exponential-inputs}\,(i) with $\mu=i\omega$. The difference $\mathbf y-\mathbf y_{\mathrm p}$ solves the homogeneous system, so it is $e^{(t-t_0)A}(\mathbf y(t_0)-\mathbf y_{\mathrm p}(t_0))$. If $\mathbf z$ is a solution with the period $T=2\pi/|\omega|$, then $\mathbf w\coloneqq\mathbf z-\mathbf y_{\mathrm p}$ has the period $T$ and tends to $\mathbf 0$, so $\mathbf w(t)=\mathbf w(t+kT)\to\mathbf 0$ as $k\to\infty$ for every $t$, and $\mathbf w=\mathbf 0$. The last statement is \cref{prop:real-solutions}\,(i) applied to the inhomogeneous system, whose input is the real part of $e^{i\omega t}\mathbf c$.
\end{proof}

The matrix $H(i\omega)$ is the \term{frequency response} of the system: an input that oscillates with the frequency $\omega$ produces, once the transient has died out, an output that oscillates with the same frequency, the vector of amplitudes and phases being $H(i\omega)\mathbf c$. The function $H(z)=(zI-A)^{-1}$ of a complex variable is the \term{resolvent} of $A$. For the companion system of $p(D)y=g$ with constant coefficients, $p(\lambda)=\lambda^k+a_{k-1}\lambda^{k-1}+\cdots+a_0$, the first component of $H(i\omega)\mathbf e_k$ is $1/p(i\omega)$: the vector $\mathbf v\coloneqq(1,i\omega,\ldots,(i\omega)^{k-1})/p(i\omega)$ satisfies $(i\omega I-C)\mathbf v=\mathbf e_k$ for the matrix $C$ of \eqref{eq:companion-system}, since the first $k-1$ rows read $i\omega v_j-v_{j+1}=0$ and the last row reads $i\omega v_k+\sum_ja_{j-1}v_j=p(i\omega)/p(i\omega)=1$. So the theorem generalises the amplitude response $R(\omega)=1/|p(i\omega)|$ of \cref{thm:frequency-response}. The eigenvalues of $H(i\omega)$ are the numbers $1/(i\omega-\lambda_j)$, so $\|H(i\omega)\|_2\ge1/\min_j|i\omega-\lambda_j|$ by \cref{prop:spectral-radius-norm}, with equality for normal matrices (\cref{prop:singular-eigen}): the response is large at frequencies close to the imaginary parts of eigenvalues that lie near the imaginary axis, which is resonance in the language of systems. For nonnormal matrices it can be much larger than this bound, for the same reason as the transient growth of \cref{sec:linear-stability}, a subject that Trefethen and Embree develop through pseudospectra~\cite{trefethen2005}. For $\omega=0$ the input is constant and the response $H(0)\mathbf c=-A^{-1}\mathbf c$ is the steady state of \cref{thm:bounded-input}.

\section{Second-Order Systems and Normal Modes}
\label{sec:normal-modes}

\subsection{Undamped Vibrations}

Mechanical systems near a position of stable equilibrium, such as the masses and springs of \cref{sec:ode-systems}, a building swaying in the wind, or the atoms of a molecule, are described by equations of the form
\begin{equation}
M\mathbf y''+K\mathbf y=\mathbf 0 ,
\label{eq:vibrations}
\end{equation}
where $\mathbf y\in\mathbb R^n$ collects the displacements from equilibrium, the \term{mass matrix} $M$ and the \term{stiffness matrix} $K$ are real, symmetric and positive definite, and the kinetic and potential energies are the quadratic forms $\langle\mathbf y',M\mathbf y'\rangle/2$ and $\langle\mathbf y,K\mathbf y\rangle/2$ (\cref{sec:quadratic-forms}). With $\mathbf v\coloneqq\mathbf y'$, the system is equivalent to the first-order system of order $2n$ with the matrix
\[
B=\begin{bmatrix}0&I_n\\-M^{-1}K&0\end{bmatrix},
\]
so it has exactly one solution for given initial positions $\mathbf y(0)$ and velocities $\mathbf y'(0)$, by \cref{thm:linear-system}. Rather than computing $e^{tB}$, we use \cref{thm:generalised-eigen}: there is a real invertible $S$ with $S^{\mathsf T}MS=I$ and $S^{\mathsf T}KS=\Lambda=\operatorname{diag}(\lambda_1,\ldots,\lambda_n)$, and its columns $\mathbf s_j$ satisfy $K\mathbf s_j=\lambda_jM\mathbf s_j$ and $\mathbf s_i^{\mathsf T}M\mathbf s_j=\delta_{ij}$.

\begin{theorem}[Normal modes]
\label{thm:normal-modes}
Let $M,K\in\mathbb R^{n\times n}$ be symmetric and positive definite, and let $S$, $\Lambda$ and $\mathbf s_j$ be as above. Then $\lambda_j>0$, and with $\omega_j\coloneqq\sqrt{\lambda_j}$ the real solutions of \eqref{eq:vibrations} are exactly the functions
\begin{equation}
\mathbf y(t)=\sum_{j=1}^n\bigl(a_j\cos\omega_jt+b_j\sin\omega_jt\bigr)\,\mathbf s_j ,
\label{eq:mode-expansion}
\end{equation}
with real constants $a_j$, $b_j$; the solution with the initial values $\mathbf y(0)$ and $\mathbf y'(0)$ has $a_j=\mathbf s_j^{\mathsf T}M\mathbf y(0)$ and $b_j=\mathbf s_j^{\mathsf T}M\mathbf y'(0)/\omega_j$. The energy
\[
E\coloneqq\frac12\langle\mathbf y',M\mathbf y'\rangle+\frac12\langle\mathbf y,K\mathbf y\rangle=\sum_{j=1}^n\frac12\bigl(z_j'^2+\omega_j^2z_j^2\bigr),\qquad\mathbf z\coloneqq S^{-1}\mathbf y ,
\]
is constant along every solution, and so is each term of the sum.
\end{theorem}

\begin{proof}
$\lambda_j=\mathbf s_j^{\mathsf T}K\mathbf s_j>0$, since $K$ is positive definite and $\mathbf s_j\ne\mathbf 0$. From $S^{\mathsf T}MS=I$ we get $S^{-1}=S^{\mathsf T}M$. Let $\mathbf y$ be a solution and $\mathbf z\coloneqq S^{-1}\mathbf y=S^{\mathsf T}M\mathbf y$. Then
\[
\mathbf z''=S^{\mathsf T}M\mathbf y''=-S^{\mathsf T}K\mathbf y=-S^{\mathsf T}KS\,\mathbf z=-\Lambda\mathbf z ,
\]
that is, $z_j''+\omega_j^2z_j=0$ for each $j$, and by \cref{lem:harmonic-unique}, in the form of \cref{sec:linear-operators}, $z_j=a_j\cos\omega_jt+b_j\sin\omega_jt$ with $a_j=z_j(0)$ and $\omega_jb_j=z_j'(0)$. So $\mathbf y=S\mathbf z$ has the form \eqref{eq:mode-expansion}, with the stated constants. Conversely, every function \eqref{eq:mode-expansion} is a solution, since $M\mathbf y''=-\sum_j\omega_j^2(\cdots)M\mathbf s_j=-\sum_j(\cdots)K\mathbf s_j=-K\mathbf y$. For the energy, $\langle\mathbf y',M\mathbf y'\rangle=\langle S\mathbf z',MS\mathbf z'\rangle=\|\mathbf z'\|^2$ and $\langle\mathbf y,K\mathbf y\rangle=\langle\mathbf z,\Lambda\mathbf z\rangle=\sum_j\omega_j^2z_j^2$, and each $z_j'^2+\omega_j^2z_j^2$ is constant by \cref{prop:oscillator-energy} with $m=1$ and $k=\omega_j^2$.
\end{proof}

The vectors $\mathbf s_j$ are the \term{normal modes} and the numbers $\omega_j$ the \term{natural frequencies} of the system. If the system is released from rest in the shape of a mode, $\mathbf y(0)=\mathbf s_j$ and $\mathbf y'(0)=\mathbf 0$, it moves as $\mathbf y(t)=\cos(\omega_jt)\,\mathbf s_j$: all coordinates oscillate in step, with one frequency, and the shape stays the same. Every motion is a superposition of such motions, with amplitudes fixed by the initial values, and each mode keeps its own energy for ever: there is no exchange of energy between modes. The theorem keeps the promise made after \cref{thm:generalised-eigen}. In the language of \cref{sec:linear-stability}, the matrix $B$ has the eigenvalues $\pm i\omega_j$, since an eigenvector $(\mathbf u,\lambda\mathbf u)$ for $\lambda$ requires $(\lambda^2M+K)\mathbf u=\mathbf 0$, so that $\lambda^2=-\lambda_j$ for some $j$; they are all semisimple, by \cref{thm:linear-stability}\,(iii), since the solutions are bounded; the equilibrium is stable but not asymptotically stable. A motion that combines modes with incommensurable frequencies, frequencies whose ratios are irrational, never repeats exactly.

\subsection{Two Masses and Beats}

For the two masses of \cref{fig:systems-examples}B, $M=mI_2$ and $K=kT_2$, the generalised eigenvalue problem $K\mathbf s=\lambda M\mathbf s$ is the eigenvalue problem $T_2\mathbf s=(\lambda m/k)\mathbf s$, with the eigenvalues $1$ and $3$ of $T_2$. So
\[
\omega_1=\sqrt{\frac km},\quad\mathbf s_1=\frac1{\sqrt{2m}}\begin{bmatrix}1\\1\end{bmatrix},\qquad\omega_2=\sqrt{\frac{3k}m},\quad\mathbf s_2=\frac1{\sqrt{2m}}\begin{bmatrix}1\\-1\end{bmatrix}.
\]
In the first mode the masses move together, the middle spring is never stretched, and each mass oscillates on its outer spring alone, with the frequency $\sqrt{k/m}$. In the second they move against each other; the middle spring is stretched by twice the displacement, the restoring force on each mass is $3k$ times its displacement, and the frequency is $\sqrt{3k/m}$.

The modes explain a striking experiment. Let the outer springs have the constant $k$ and the middle spring a much smaller constant $\kappa$, so that the masses are only weakly coupled. Then $K$ has the rows $(k+\kappa,-\kappa)$ and $(-\kappa,k+\kappa)$, the modes are the same, and the frequencies are $\omega_1=\sqrt{k/m}$ and $\omega_2=\sqrt{(k+2\kappa)/m}$, close to each other. Release the first mass from the displacement $1$ and the second from rest at its equilibrium: $\mathbf y(0)=\mathbf e_1$, $\mathbf y'(0)=\mathbf 0$. By \eqref{eq:mode-expansion} and the identities for sums of cosines,
\begin{gather*}
y_1(t)=\frac{\cos\omega_1t+\cos\omega_2t}2=\cos\frac{(\omega_2-\omega_1)t}2\,\cos\frac{(\omega_1+\omega_2)t}2,\\
y_2(t)=\frac{\cos\omega_1t-\cos\omega_2t}2=\sin\frac{(\omega_2-\omega_1)t}2\,\sin\frac{(\omega_1+\omega_2)t}2 .
\end{gather*}
Both masses oscillate rapidly, with nearly the frequency $\omega_1$, inside slowly varying envelopes. After the time $\pi/(\omega_2-\omega_1)$ the first mass has come to rest and the second swings with the full amplitude, and then the energy flows back (\cref{fig:modes}A). These are the beats of \cref{sec:oscillations}, now as an exchange of energy between two oscillators. The energy of each mode is constant; what moves back and forth is the energy of each mass, which is not a mode.

\subsection{Chains and the Vibrating String}

For $n$ equal masses joined by $n+1$ equal springs, $M=mI$ and $K=kT_n$, and \cref{thm:second-difference-eigen} gives the modes and frequencies explicitly: the $j$th mode has the components $\sin(ij\pi/(n+1))$, $i=1,\ldots,n$, and
\[
\omega_j=2\sqrt{\frac km}\,\sin\frac{j\pi}{2(n+1)}\qquad(j=1,\ldots,n).
\]
The slowest mode is a single arch, in which all masses move in the same direction; the $j$th mode changes sign $j-1$ times along the chain (\cref{fig:modes}B). The frequencies increase with the number of sign changes, from about $\pi\sqrt{k/m}/(n+1)$ for the first mode to nearly $2\sqrt{k/m}$ for the last.

The chain is a model of a vibrating string. Let a string of length $L$, with the mass $\rho$ per unit length, be stretched by the tension $F$, and replace it by $n$ beads of mass $\rho h$ at the spacing $h=L/(n+1)$. For small transverse displacements, the string between two neighbouring beads pulls each of them towards the other with a transverse force of about $F$ times the slope $(y_{i+1}-y_i)/h$, so the beads behave like masses joined by springs of the constant $F/h$. The frequencies are then
\[
\omega_j=2\sqrt{\frac F{\rho h^2}}\,\sin\frac{j\pi h}{2L}\to\frac{j\pi}L\sqrt{\frac F\rho}\qquad(h\to0),
\]
since $\sin x/x\to1$, and the modes approach the functions $\sin(j\pi x/L)$. The frequencies of the string are the integer multiples of the fundamental frequency $(\pi/L)\sqrt{F/\rho}$, the harmonics on which music is built, and the modes are standing waves. The limit itself is described by a partial differential equation, the wave equation, which this volume does not treat; as for the matrices $T_n$ in \cref{sec:structured-eigen}, the low modes of the chain approximate those of the string well, and the high ones do not.

\begin{figure}
\sdpanel{A}{\begin{tikzpicture}
\begin{axis}[sdplot, width=0.56\textwidth, height=0.36\textwidth, xmin=0, xmax=70, ymin=-1.1, ymax=1.1,
  xtick={0,20,40,60}, ytick={-1,0,1}, xlabel={$t$}]
\addplot[black, line width=0.5pt, domain=0:70, samples=280] {(cos(deg(x))+cos(deg(1.0954451*x)))/2};
\addplot[black!45, line width=0.5pt, domain=0:70, samples=280] {(cos(deg(x))-cos(deg(1.0954451*x)))/2};
\end{axis}
\end{tikzpicture}}\hfill
\sdpanel{B}{\begin{tikzpicture}
\begin{axis}[sdplot, width=0.4\textwidth, height=0.36\textwidth, xmin=0, xmax=7, ymin=-1.15, ymax=1.15,
  xtick={0,1,2,3,4,5,6,7}, ytick={-1,0,1}, xlabel={mass $i$}]
\addplot[black, line width=0.6pt, mark=*, mark size=1.3pt, samples at={0,...,7}] {sin(deg(x*pi/7))};
\addplot[black, dashed, line width=0.6pt, mark=square*, mark size=1.2pt, mark options={solid}, samples at={0,...,7}] {sin(deg(2*x*pi/7))};
\addplot[black!45, line width=0.6pt, mark=triangle*, mark size=1.5pt, samples at={0,...,7}] {sin(deg(3*x*pi/7))};
\end{axis}
\end{tikzpicture}}
\caption{Normal modes. (A) Two weakly coupled masses with $m=k=1$ and $\kappa=0.1$, so that $\omega_1=1$ and $\omega_2=\sqrt{1.2}=1.0954$, released with $\mathbf y(0)=\mathbf e_1$ and $\mathbf y'(0)=\mathbf 0$: the displacements $y_1$ (black) and $y_2$ (grey). The energy passes from the first mass to the second and back; the first mass is at rest near $t=\pi/(\omega_2-\omega_1)=32.9$. (B) The first three modes of a chain of $n=6$ masses between two walls (positions $0$ and $7$): $\sin(ij\pi/7)$ for $j=1$ (solid), $2$ (dashed) and $3$ (grey), with the frequencies $2\sin(j\pi/14)=0.445$, $0.868$ and $1.247$ for $k=m=1$.\label{fig:modes}}
\end{figure}

\subsection{Damping}

Friction adds a term proportional to the velocities: $M\mathbf y''+C\mathbf y'+K\mathbf y=\mathbf 0$ with a real symmetric \term{damping matrix} $C$.

\begin{proposition}[Damped vibrations]
\label{prop:damped-vibrations}
Let $M$, $C$ and $K\in\mathbb R^{n\times n}$ be symmetric and positive definite. Along every solution of $M\mathbf y''+C\mathbf y'+K\mathbf y=\mathbf 0$ the energy $E$ of \cref{thm:normal-modes} satisfies $E'=-\langle\mathbf y',C\mathbf y'\rangle\le0$, and every eigenvalue of the matrix of the equivalent first-order system has negative real part. Consequently every solution tends to $\mathbf 0$, together with its derivative, exponentially.
\end{proposition}

\begin{proof}
Since $M$ and $K$ are symmetric, $E'=\langle\mathbf y',M\mathbf y''\rangle+\langle\mathbf y',K\mathbf y\rangle=\langle\mathbf y',M\mathbf y''+K\mathbf y\rangle=-\langle\mathbf y',C\mathbf y'\rangle$. The first-order system for $(\mathbf y,\mathbf v)$, $\mathbf v=\mathbf y'$, has the matrix with the block rows $(0,I)$ and $(-M^{-1}K,-M^{-1}C)$. If $(\mathbf u,\mathbf w)\ne\mathbf 0$ is an eigenvector for $\lambda\in\mathbb C$, then $\mathbf w=\lambda\mathbf u$ and $-M^{-1}K\mathbf u-M^{-1}C\mathbf w=\lambda\mathbf w$, so $\mathbf u\ne\mathbf 0$ and $(\lambda^2M+\lambda C+K)\mathbf u=\mathbf 0$. Multiplying by $\mathbf u^{\mathsf H}$ gives $\lambda^2\mu+\lambda\gamma+\kappa=0$ with $\mu\coloneqq\mathbf u^{\mathsf H}M\mathbf u$, $\gamma\coloneqq\mathbf u^{\mathsf H}C\mathbf u$ and $\kappa\coloneqq\mathbf u^{\mathsf H}K\mathbf u$, which are positive numbers, since the matrices are positive definite (\cref{def:hermitian-definite}). By \cref{thm:hurwitz-small}\,(i), applied to $\lambda^2+(\gamma/\mu)\lambda+\kappa/\mu$, the root $\lambda$ has negative real part. The last statement follows from \cref{prop:decay-estimate}.
\end{proof}

When $C=aM+bK$ with constants $a,b\ge0$, the case of \term{proportional damping}, the transformation of \cref{thm:normal-modes} decouples the damped system as well: $z_j''+(a+b\lambda_j)z_j'+\lambda_jz_j=0$, a damped oscillator for each mode, with the damping ratio $(a+b\lambda_j)/(2\omega_j)$ of \eqref{eq:oscillator-normalised}. For general $C$ the modes are coupled by the damping, and the system must be treated as a first-order system of order $2n$, as in \cref{sec:eigenvalue-method}.

\subsection{Estimating the Lowest Frequency}

The lowest natural frequency is often the most important one, since it describes the slowest and largest motions of a structure. It can be estimated without solving an eigenvalue problem.

\begin{proposition}[Rayleigh's principle]
\label{prop:rayleigh-frequency}
Under the hypotheses of \cref{thm:normal-modes}, with $\lambda_1\le\cdots\le\lambda_n$,
\[
\omega_1^2=\min_{\mathbf x\ne\mathbf 0}\frac{\langle\mathbf x,K\mathbf x\rangle}{\langle\mathbf x,M\mathbf x\rangle},\qquad\omega_n^2=\max_{\mathbf x\ne\mathbf 0}\frac{\langle\mathbf x,K\mathbf x\rangle}{\langle\mathbf x,M\mathbf x\rangle},
\]
and the minimum is attained exactly at the nonzero solutions $\mathbf x$ of $K\mathbf x=\lambda_1M\mathbf x$.
\end{proposition}

\begin{proof}
Let $M=R^{\mathsf T}R$ be the Cholesky factorisation (\cref{thm:cholesky}) and $\mathbf x=R^{-1}\mathbf u$. The quotient becomes $\langle\mathbf u,G\mathbf u\rangle/\|\mathbf u\|^2$ with the symmetric matrix $G\coloneqq R^{-\mathsf T}KR^{-1}$, whose eigenvalues are the $\lambda_j$, by the proof of \cref{thm:generalised-eigen}. By \cref{thm:rayleigh}, its minimum and maximum over $\mathbf u\ne\mathbf 0$ are $\lambda_1$ and $\lambda_n$, the minimum being attained exactly at the eigenvectors of $G$ for $\lambda_1$, and $\mathbf u\mapsto R^{-1}\mathbf u$ maps these onto the solutions of $K\mathbf x=\lambda_1M\mathbf x$.
\end{proof}

The quotient has a mechanical meaning. If the system moved as $\cos(\omega t)\,\mathbf x$ with a fixed shape $\mathbf x$, its largest potential energy would be $\langle\mathbf x,K\mathbf x\rangle/2$ and its largest kinetic energy $\omega^2\langle\mathbf x,M\mathbf x\rangle/2$, and equating the two gives $\omega^2$ as the quotient. Every guessed shape therefore gives an upper bound for $\omega_1^2$, and, by \cref{prop:rayleigh-quadratic}, an error of order $\varepsilon$ in the shape causes an error of order only $\varepsilon^2$ in the estimate. For the chain of $n=10$ masses with $k=m=1$, the exact value is $\omega_1^2=4\sin^2(\pi/22)=0.081014$. The parabolic shape $x_i=i(11-i)$, which resembles the single arch of the first mode, gives the quotient $5/61=0.081967$, so that the estimated frequency is too high by only $0.59$ per cent. The constant shape $x_i=1$, a poor guess that stretches only the two end springs, gives $2/10$, which is about $2.5$ times too large.

\section{Nonlinear Systems Near an Equilibrium}
\label{sec:nonlinear-systems}

\subsection{Linearisation}

Most systems in nature are nonlinear, and few nonlinear systems can be solved by formulas. Near an equilibrium, however, a smooth system is approximately linear, and the theory of the preceding sections describes it there. Let $U\subseteq\mathbb R^n$ be open, let $\mathbf f\colon U\to\mathbb R^n$ be continuously differentiable, and consider the autonomous system $\mathbf y'=\mathbf f(\mathbf y)$. By the results of \cref{sec:existence-uniqueness} in their form for systems, recalled in \cref{sec:ode-systems}, every initial value problem has a unique maximal solution, and \cref{lem:orbits} describes the orbits. If $\mathbf y^*$ is an equilibrium, the derivative $A\coloneqq D\mathbf f(\mathbf y^*)$, the Jacobian matrix of $\mathbf f$ at $\mathbf y^*$ (\cref{def:differentiable-rn}), gives
\[
\mathbf f(\mathbf y)=A(\mathbf y-\mathbf y^*)+\mathbf r(\mathbf y),\qquad\frac{\|\mathbf r(\mathbf y)\|}{\|\mathbf y-\mathbf y^*\|}\to0\quad(\mathbf y\to\mathbf y^*),
\]
and the \term{linearisation} of the system at $\mathbf y^*$ is the linear system $\mathbf u'=A\mathbf u$ for the deviation $\mathbf u=\mathbf y-\mathbf y^*$. For the predator--prey system \eqref{eq:lotka-volterra}, \cref{sec:ode-systems} found it by dropping the product of the deviations, which is the remainder $\mathbf r$. The question is which properties of the linearisation the nonlinear system inherits. Asymptotic stability is inherited, with a complete proof through the quadratic Lyapunov functions of \cref{sec:linear-stability}.

\begin{theorem}[Linearised stability for systems]
\label{thm:linearised-stability-systems}
Let $\mathbf f\colon U\to\mathbb R^n$ be continuously differentiable on the open set $U$, let $\mathbf f(\mathbf y^*)=\mathbf 0$, and suppose that every eigenvalue of $A=D\mathbf f(\mathbf y^*)$ has negative real part. Then $\mathbf y^*$ is asymptotically stable. More precisely, there are constants $\delta>0$, $C\ge1$ and $\gamma>0$ such that every solution with $\|\mathbf y(0)-\mathbf y^*\|<\delta$ exists for all $t\ge0$ and satisfies $\|\mathbf y(t)-\mathbf y^*\|\le Ce^{-\gamma t}\|\mathbf y(0)-\mathbf y^*\|$.
\end{theorem}

\begin{proof}
By \cref{thm:lyapunov-equation}, there is a real symmetric positive definite $P$ with $A^{\mathsf T}P+PA=-I$. Let $V(\mathbf u)\coloneqq\langle\mathbf u,P\mathbf u\rangle$, so that $\lambda_1(P)\|\mathbf u\|^2\le V(\mathbf u)\le\lambda_n(P)\|\mathbf u\|^2$. Choose $\rho>0$ such that the closed ball $\bar B(\mathbf y^*,\rho)$ lies in $U$ and $\|\mathbf r(\mathbf y)\|\le\|\mathbf y-\mathbf y^*\|/(4\lambda_n(P))$ on it, which is possible by the property of $\mathbf r$. For a solution $\mathbf y$ with values in this ball, $\mathbf u\coloneqq\mathbf y-\mathbf y^*$ satisfies $\mathbf u'=A\mathbf u+\mathbf r(\mathbf y)$, and
\[
\begin{aligned}
\frac d{dt}V(\mathbf u)&=\langle\mathbf u,(A^{\mathsf T}P+PA)\mathbf u\rangle+2\langle\mathbf u,P\mathbf r(\mathbf y)\rangle\\
&\le-\|\mathbf u\|^2+2\lambda_n(P)\|\mathbf u\|\,\|\mathbf r(\mathbf y)\|\le-\frac12\|\mathbf u\|^2\le-\frac{V(\mathbf u)}{2\lambda_n(P)},
\end{aligned}
\]
by the Cauchy--Schwarz inequality and $\|P\|_2=\lambda_n(P)$.

Let $c\coloneqq\lambda_1(P)\rho^2$, and let $\mathbf y$ be the maximal solution with $V(\mathbf y(0)-\mathbf y^*)<c$; this holds if $\|\mathbf y(0)-\mathbf y^*\|<\delta\coloneqq\rho\sqrt{\lambda_1(P)/\lambda_n(P)}$. Points with $V(\mathbf u)<c$ satisfy $\|\mathbf u\|<\rho$. Let $T$ be the supremum of the times $\tau<\omega_+$ such that $V(\mathbf u(t))<c$ for all $t\in[0,\tau]$. On $[0,T)$ the solution lies in the ball, so $V(\mathbf u(t))$ does not increase and stays at most $V(\mathbf u(0))<c$. If $T<\omega_+$, continuity would give $V(\mathbf u(T))\le V(\mathbf u(0))<c$ and then $V(\mathbf u(t))<c$ for $t$ slightly beyond $T$, contradicting the definition of $T$; so $T=\omega_+$. The solution therefore stays in the compact ball $\bar B(\mathbf y^*,\rho)\subseteq U$, and by \cref{thm:maximal-solution} in its form for systems, $\omega_+=\infty$. Finally, the differential inequality gives $V(\mathbf u(t))\le e^{-t/(2\lambda_n(P))}V(\mathbf u(0))$ as in the proof of \cref{thm:lyapunov-equation}, and hence the estimate with $C=\sqrt{\kappa_2(P)}$ and $\gamma=1/(4\lambda_n(P))$.
\end{proof}

The theorem is the system version of \cref{thm:linearised-stability}\,(i), where $f'(y^*)<0$ played the role of the eigenvalues. The second half of that theorem, instability when $f'(y^*)>0$, also carries over.

\begin{theorem}[Linearised instability]
\label{thm:linearised-instability}
Under the hypotheses of \cref{thm:linearised-stability-systems}, except that some eigenvalue of $A=D\mathbf f(\mathbf y^*)$ has positive real part, the equilibrium $\mathbf y^*$ is unstable.
\end{theorem}

\begin{proof}[sketch]
Split $\mathbb R^n$ into the sum $E_+$ of the real generalised eigenspaces of the eigenvalues with positive real part and the sum $E_-$ of the others, both invariant under $A$. Choose $\sigma>0$ smaller than the real parts of the eigenvalues in $E_+$. Then the restriction of $\sigma I-A$ to $E_+$ and the restriction of $A-\sigma I$ to $E_-$ have only eigenvalues with negative real part, and Lyapunov equations for them, with the identity on the right side, provide positive definite quadratic forms $V_+$ on $E_+$ and $V_-$ on $E_-$. The function $W\coloneqq V_+(\mathbf u_+)-V_-(\mathbf u_-)$ of the components of $\mathbf u=\mathbf y-\mathbf y^*$ then satisfies $W'=2\sigma W+\|\mathbf u_+\|^2+\|\mathbf u_-\|^2$ along the solutions of the linearisation, so that $W'\ge c\,(W+\|\mathbf u\|^2)$ with a constant $c>0$ wherever $W\ge0$, and, as in the proof of \cref{thm:linearised-stability-systems}, the remainder $\mathbf r$ does not spoil an inequality of this kind in a small ball around $\mathbf y^*$. A solution that starts in the cone $W>0$, however close to $\mathbf y^*$, has $W$ growing at least like $e^{ct}$ while it stays in the ball, and since $W$ is bounded on the ball it must leave the ball. The details are in Perko~\cite[Chap.~2]{perko2001} and Hirsch, Smale and Devaney~\cite{hirsch2013}.
\end{proof}

When the linearisation has an eigenvalue with real part zero and none with positive real part, the linearisation does not decide stability. The systems
\[
\mathbf y'=\begin{bmatrix}0&-1\\1&0\end{bmatrix}\mathbf y\pm\|\mathbf y\|^2\mathbf y
\]
have the same linearisation at $\mathbf 0$, a centre. But $\rho\coloneqq\|\mathbf y\|^2$ satisfies $\rho'=2\langle\mathbf y,\mathbf y'\rangle=\pm2\|\mathbf y\|^4=\pm2\rho^2$, since the rotation part is perpendicular to $\mathbf y$. With the minus sign, $\rho(t)=\rho_0/(1+2\rho_0t)\to0$: the origin is asymptotically stable, and the solutions spiral in, slowly. With the plus sign, $\rho(t)=\rho_0/(1-2\rho_0t)$ blows up at $t=1/(2\rho_0)$: the origin is unstable. The phenomenon is that of \cref{thm:linearised-stability} for $f'(y^*)=0$: a small nonlinear term decides what the linear part leaves open, just as a small change of the matrix turns a centre into a spiral (\cref{sec:phase-plane}).

At an equilibrium without eigenvalues on the imaginary axis, a \term{hyperbolic} equilibrium, the linearisation describes not only the stability but the whole phase portrait near the equilibrium, up to a continuous deformation.

\begin{theorem}[Hartman--Grobman]
\label{thm:hartman-grobman}
Let $\mathbf f\colon U\to\mathbb R^n$ be continuously differentiable, $\mathbf f(\mathbf y^*)=\mathbf 0$, and suppose that no eigenvalue of $A=D\mathbf f(\mathbf y^*)$ has real part zero. Then there are a neighbourhood $N$ of $\mathbf y^*$, a neighbourhood $N_0$ of $\mathbf 0$ and a continuous bijection $\mathbf h\colon N\to N_0$ with a continuous inverse and $\mathbf h(\mathbf y^*)=\mathbf 0$, such that $\mathbf h$ maps the pieces in $N$ of the orbits of $\mathbf y'=\mathbf f(\mathbf y)$ onto the pieces in $N_0$ of the orbits of $\mathbf u'=A\mathbf u$, preserving the direction of time.
\end{theorem}

\begin{proof}[sketch]
One looks for $\mathbf h$ in the form $\mathbf h(\mathbf y)=\mathbf y-\mathbf y^*+\mathbf g(\mathbf y)$ with a small bounded continuous $\mathbf g$, such that $\mathbf h$ carries the solutions of the nonlinear system at the time $1$ to those of the linear one, $\mathbf h\circ\boldsymbol\varphi_1=e^A\circ\mathbf h$, where $\boldsymbol\varphi_1$ maps an initial value to the value of the solution at the time $1$. After $\mathbf f$ has been modified outside a small ball so that it differs from its linear part by a function with a small Lipschitz constant, this equation becomes a fixed point equation for $\mathbf g$ in the space of bounded continuous functions, and the splitting of $\mathbb R^n$ into the directions in which $e^A$ contracts and those in which it expands, which exists because no eigenvalue lies on the imaginary axis, makes it a contraction. Its fixed point gives $\mathbf h$, and a similar argument for the continuous time completes the proof. The theorem is due to Grobman and to Hartman~\cite{grobman1959,hartman1960}; complete proofs are in Hartman~\cite[Chap.~IX]{hartman2002} and Perko~\cite[Chap.~2]{perko2001}.
\end{proof}

By the theorem, a saddle of the linearisation is a saddle of the nonlinear system, with two orbits that come into the equilibrium and two that leave it. These orbits, together with the equilibrium, form the \term{stable} and the \term{unstable manifold} of the equilibrium, and the stable manifold theorem, which accompanies the Hartman--Grobman theorem, adds that they are smooth curves tangent at the equilibrium to the stable and unstable lines of the linearisation~\cite[Chap.~2]{perko2001}; the incoming orbits, the \term{separatrices}, divide the plane near the saddle into regions in which the solutions behave differently. Nodes and spirals of the linearisation are nodes and spirals of the nonlinear system up to deformation. Centres, stars and degenerate nodes, the nonhyperbolic and borderline cases, are not covered, in agreement with the example above.

\subsection{The Pendulum}

A pendulum of length $l$ under gravity $g$ satisfies $\theta''=-(g/l)\sin\theta$ for the angle $\theta$ from the downward vertical (\cref{sec:exact}). Measuring time in units of $\sqrt{l/g}$ turns the equation into $\theta''+\sin\theta=0$, and with the angular velocity $v\coloneqq\theta'$ into the system
\begin{equation}
\theta'=v,\qquad v'=-\sin\theta .
\label{eq:pendulum}
\end{equation}
Its equilibria are the points $(k\pi,0)$, $k\in\mathbb Z$, the pendulum hanging down for even $k$ and standing upright for odd $k$. The Jacobian matrix is
\[
D\mathbf f(\theta,v)=\begin{bmatrix}0&1\\-\cos\theta&0\end{bmatrix},
\]
which at $(\pi,0)$ has the rows $(0,1)$ and $(1,0)$ and the eigenvalues $\pm1$: the upright position is a saddle, and unstable by \cref{thm:linearised-instability}. At $(0,0)$ the linearisation is the harmonic oscillator, a centre, and the linearisation does not decide. The energy decides. As \cref{sec:exact} showed, $E(\theta,v)\coloneqq v^2/2-\cos\theta$ is constant along every solution, so every orbit lies on a level curve of $E$.

\begin{proposition}[The pendulum]
\label{prop:pendulum}
For the system \eqref{eq:pendulum}:
\begin{enumerate}
\item The equilibrium $(0,0)$ is stable but not asymptotically stable.
\item For $0<\theta_0<\pi$, the solution with $\theta(0)=\theta_0$ and $v(0)=0$ is periodic, with the period
\[
T(\theta_0)=4\int_0^{\theta_0}\frac{d\theta}{\sqrt{2(\cos\theta-\cos\theta_0)}}=4K\Bigl(\sin\frac{\theta_0}2\Bigr),
\]
where
\[
K(k)\coloneqq\int_0^{\pi/2}\frac{d\varphi}{\sqrt{1-k^2\sin^2\varphi}}\qquad(0\le k<1).
\]
\item The period increases strictly with the amplitude $\theta_0$, tends to $2\pi$ as $\theta_0\to0$ and to infinity as $\theta_0\to\pi$.
\end{enumerate}
\end{proposition}

\begin{proof}
(i) The function $E+1=v^2/2+2\sin^2(\theta/2)$ is positive except at the points $(2k\pi,0)$, where it is $0$. Let $0<\varepsilon<\pi$, let $m>0$ be the minimum of $E+1$ on the circle of radius $\varepsilon$ about the origin, and choose $\delta\in(0,\varepsilon)$ with $E+1<m$ on the disc of radius $\delta$, which is possible by continuity. A solution that starts in this disc keeps the value $E+1<m$, so it never meets the circle and, by continuity, stays inside it; it exists for all time, by \cref{thm:maximal-solution}, since it stays in a compact set. So the origin is stable. It is not asymptotically stable: a solution starting at a point $\ne(0,0)$ close to it has $E\ne-1$, and if it tended to $(0,0)$, the continuous function $E$ would tend to $-1$.

(ii) Let $(\theta,v)$ be the solution and $T_0\coloneqq\int_0^{\theta_0}d\eta/\sqrt{2(\cos\eta-\cos\theta_0)}$. The integral converges: near $\eta=\theta_0$, $\cos\eta-\cos\theta_0\ge(\sin\theta_0)(\theta_0-\eta)/2$ by the mean value theorem and the continuity of $\sin$, and the integrand is at most a constant times $(\theta_0-\eta)^{-1/2}$ (\cref{prop:p-integrals}). By conservation of energy, $v^2=2(\cos\theta-\cos\theta_0)$ along the solution, so $\cos\theta(t)\ge\cos\theta_0$ for all $t$. The angles with this property form the disjoint intervals $[2k\pi-\theta_0,2k\pi+\theta_0]$, $k\in\mathbb Z$, and the continuous function $\theta$ starts in the interval with $k=0$, so it stays there (\cref{cor:ivt}): $|\theta(t)|\le\theta_0$ for all $t$. Since $v'(0)=-\sin\theta_0<0$, the velocity is negative for small $t>0$; let $\tau\in(0,\infty]$ be the supremum of the times $s$ with $v<0$ on $(0,s)$. On $(0,\tau)$, $\theta$ decreases strictly, $v=-\sqrt{2(\cos\theta-\cos\theta_0)}$, and the substitution $\eta=\theta(s)$, made on $[s_0,t]$ with $0<s_0<t<\tau$ and followed by the limit $s_0\to0$, gives
\[
t=\int_0^t\frac{-\theta'(s)}{\sqrt{2(\cos\theta(s)-\cos\theta_0)}}\,ds=\int_{\theta(t)}^{\theta_0}\frac{d\eta}{\sqrt{2(\cos\eta-\cos\theta_0)}}\le2T_0 ,
\]
so $\tau\le2T_0$ is finite and $v(\tau)=0$. Then $\cos\theta(\tau)=\cos\theta_0$ with $\theta(\tau)<\theta_0$ and $|\theta(\tau)|\le\theta_0$, so $\theta(\tau)=-\theta_0$, and the formula gives $\tau=2T_0$. Now $(\tilde\theta,\tilde v)(t)\coloneqq(-\theta(t-\tau),-v(t-\tau))$ is also a solution, since the system is unchanged by $(\theta,v)\mapsto(-\theta,-v)$ and by shifts of time, and it has the value $(-\theta_0,0)=(\theta(\tau),v(\tau))$ at $t=\tau$. By uniqueness it equals $(\theta,v)$, and at $t=2\tau$ this gives $(\theta,v)(2\tau)=(\theta_0,0)$. By \cref{lem:orbits}\,(iii) the solution is periodic with the period $2\tau=4T_0$. For the second expression, write $k\coloneqq\sin(\theta_0/2)$, use $\cos\theta=1-2\sin^2(\theta/2)$, so that $\cos\theta-\cos\theta_0=2(k^2-\sin^2(\theta/2))$, and substitute $\sin(\theta/2)=k\sin\varphi$, $0\le\varphi<\pi/2$. Then $\cos(\theta/2)\,d\theta/2=k\cos\varphi\,d\varphi$, $\sqrt{2(\cos\theta-\cos\theta_0)}=2k\cos\varphi$, and $\cos(\theta/2)=\sqrt{1-k^2\sin^2\varphi}$, so the integrand becomes $d\varphi/\sqrt{1-k^2\sin^2\varphi}$.

(iii) As $\theta_0$ increases in $(0,\pi)$, $k=\sin(\theta_0/2)$ increases in $(0,1)$, and so does the integrand of $K(k)$ at every $\varphi\in(0,\pi/2]$; so $T$ increases strictly. Since $1-k^2\le1-k^2\sin^2\varphi\le1$, we have $\pi/2\le K(k)\le(\pi/2)/\sqrt{1-k^2}$, which tends to $\pi/2$ as $k\to0$, so $T\to2\pi$. For the other limit, $1-k^2\sin^2\varphi=(1-k^2)+k^2\cos^2\varphi\le\varepsilon^2+\cos^2\varphi$ with $\varepsilon^2\coloneqq1-k^2$, and with $\psi=\pi/2-\varphi$ and $\sin\psi\le\psi$,
\[
K(k)\ge\int_0^{\pi/2}\frac{d\psi}{\sqrt{\varepsilon^2+\sin^2\psi}}\ge\int_0^{\pi/2}\frac{d\psi}{\sqrt{\varepsilon^2+\psi^2}}=\operatorname{arsinh}\frac\pi{2\varepsilon},
\]
by \cref{prop:inverse-hyperbolic}, which tends to infinity as $k\to1$.
\end{proof}

The integral $K(k)$ is the \term{complete elliptic integral of the first kind}; it has no expression by elementary functions, and Volume~X treats its computation. \Cref{sec:systems-numerics} compares the period with a numerical solution of the pendulum. By the binomial series (\cref{thm:binomial-series}) and $\int_0^{\pi/2}\sin^2\varphi\,d\varphi=\pi/4$, $K(k)=(\pi/2)(1+k^2/4+\cdots)$, so for small amplitudes $T\approx2\pi(1+\theta_0^2/16)$, and the period of the small oscillations, $2\pi$ in our units or $2\pi\sqrt{l/g}$ in the original ones, is that of the linearisation.

The phase portrait in \cref{fig:pendulum}A shows the whole picture. The level curves of $E$ with $-1<E<1$ are closed curves around the equilibria $(2k\pi,0)$, on which the pendulum swings back and forth. The level $E=1$ consists of the saddles $((2k+1)\pi,0)$ and the curves $v=\pm2\cos(\theta/2)$ that join them: on these \term{separatrices} the pendulum creeps towards the upright position, which it approaches but, by \cref{lem:orbits}\,(ii), never reaches. For $E>1$ the pendulum has enough energy to pass over the top, and it rotates for ever in one direction, $v$ never vanishing. The plane of $(\theta,v)$ repeats with the period $2\pi$ in $\theta$, since $\theta$ and $\theta+2\pi$ describe the same position.

With friction, $\theta''+\gamma\theta'+\sin\theta=0$ with $\gamma>0$, the energy decreases, $E'=-\gamma v^2$, and the linearisation at $(0,0)$ has the trace $-\gamma$ and the determinant $1$: the hanging position is asymptotically stable, a spiral for $\gamma<2$, by \cref{thm:linearised-stability-systems,thm:hartman-grobman}. A pendulum started with a large velocity rotates a few times, loses energy at each turn, and is finally captured by one of the hanging positions, around which it spirals in (\cref{fig:pendulum}B). Which of them captures it depends on its initial state in a complicated way, and the boundaries between the regions of capture are the stable manifolds of the upright positions.

\begin{figure}
\sdpanel{A}{\begin{tikzpicture}
\begin{axis}[sdplot, width=0.92\textwidth, height=0.36\textwidth, xmin=-7, xmax=7, ymin=-3, ymax=3, clip=true, xtick={-6.2832,-3.1416,0,3.1416,6.2832}, ytick={-2,0,2}, xlabel={$\theta$}, ylabel={$v=\theta'$}, xticklabels={$-2\pi$,$-\pi$,$0$,$\pi$,$2\pi$}]
\addplot[ch25orbit, domain=0:360, samples=33] ({0.0000+2*rad(asin(0.3500*sin(x)))},{0.7000*cos(x)});
\draw[ch25arrow] (axis cs:-0.020,0.700) -- (axis cs:0.000,0.700);
\draw[ch25arrow] (axis cs:0.020,-0.700) -- (axis cs:0.000,-0.700);
\addplot[ch25orbit, domain=0:360, samples=33] ({0.0000+2*rad(asin(0.6000*sin(x)))},{1.2000*cos(x)});
\draw[ch25arrow] (axis cs:-0.020,1.200) -- (axis cs:0.000,1.200);
\draw[ch25arrow] (axis cs:0.020,-1.200) -- (axis cs:0.000,-1.200);
\addplot[ch25orbit, domain=0:360, samples=33] ({0.0000+2*rad(asin(0.8500*sin(x)))},{1.7000*cos(x)});
\draw[ch25arrow] (axis cs:-0.020,1.700) -- (axis cs:0.000,1.700);
\draw[ch25arrow] (axis cs:0.020,-1.700) -- (axis cs:0.000,-1.700);
\addplot[ch25orbit, domain=0:360, samples=33] ({-6.2832+2*rad(asin(0.6000*sin(x)))},{1.2000*cos(x)});
\draw[ch25arrow] (axis cs:-6.303,1.200) -- (axis cs:-6.283,1.200);
\draw[ch25arrow] (axis cs:-6.263,-1.200) -- (axis cs:-6.283,-1.200);
\addplot[ch25orbit, domain=0:360, samples=33] ({6.2832+2*rad(asin(0.6000*sin(x)))},{1.2000*cos(x)});
\draw[ch25arrow] (axis cs:6.263,1.200) -- (axis cs:6.283,1.200);
\draw[ch25arrow] (axis cs:6.303,-1.200) -- (axis cs:6.283,-1.200);
\addplot[black, line width=0.8pt, domain=-7:7, samples=57] {2*abs(cos(deg(x/2)))};
\addplot[black, line width=0.8pt, domain=-7:7, samples=57] {-2*abs(cos(deg(x/2)))};
\draw[ch25arrow] (axis cs:-1.595,1.397) -- (axis cs:-1.571,1.414);
\draw[ch25arrow] (axis cs:-1.547,-1.431) -- (axis cs:-1.571,-1.414);
\draw[ch25arrow] (axis cs:1.547,1.431) -- (axis cs:1.571,1.414);
\draw[ch25arrow] (axis cs:1.595,-1.397) -- (axis cs:1.571,-1.414);
\draw[ch25arrow] (axis cs:-4.736,1.431) -- (axis cs:-4.712,1.414);
\draw[ch25arrow] (axis cs:-4.688,-1.397) -- (axis cs:-4.712,-1.414);
\draw[ch25arrow] (axis cs:4.688,1.397) -- (axis cs:4.712,1.414);
\draw[ch25arrow] (axis cs:4.736,-1.431) -- (axis cs:4.712,-1.414);
\addplot[ch25orbit, domain=-7:7, samples=45] {sqrt(2*(2+cos(deg(x))))};
\addplot[ch25orbit, domain=-7:7, samples=45] {-sqrt(2*(2+cos(deg(x))))};
\draw[ch25arrow] (axis cs:-4.527,1.906) -- (axis cs:-4.500,1.892);
\draw[ch25arrow] (axis cs:-4.473,-1.878) -- (axis cs:-4.500,-1.892);
\draw[ch25arrow] (axis cs:-0.030,2.449) -- (axis cs:0.000,2.449);
\draw[ch25arrow] (axis cs:0.030,-2.449) -- (axis cs:0.000,-2.449);
\draw[ch25arrow] (axis cs:4.473,1.878) -- (axis cs:4.500,1.892);
\draw[ch25arrow] (axis cs:4.527,-1.906) -- (axis cs:4.500,-1.892);
\fill[black] (axis cs:-6.2832,0) circle (1.4pt);
\fill[black] (axis cs:0.0000,0) circle (1.4pt);
\fill[black] (axis cs:6.2832,0) circle (1.4pt);
\draw[black, line width=0.5pt, fill=white] (axis cs:-3.1416,0) circle (1.4pt);
\draw[black, line width=0.5pt, fill=white] (axis cs:3.1416,0) circle (1.4pt);
\end{axis}
\end{tikzpicture}}
\sdpanelrow
\sdpanel{B}{\begin{tikzpicture}
\begin{axis}[sdplot, width=0.92\textwidth, height=0.36\textwidth, xmin=-7, xmax=7, ymin=-3, ymax=3, clip=true, xtick={-6.2832,-3.1416,0,3.1416,6.2832}, ytick={-2,0,2}, xlabel={$\theta$}, ylabel={$v=\theta'$}, xticklabels={$-2\pi$,$-\pi$,$0$,$\pi$,$2\pi$}]
\addplot[ch25orbit] coordinates {(-6.900,2.000) (-6.670,1.998) (-6.451,1.973) (-6.246,1.929) (-6.066,1.871) (-5.900,1.803) (-5.750,1.728) (-5.615,1.649) (-5.486,1.564) (-5.372,1.479) (-5.265,1.390) (-5.164,1.299) (-5.070,1.207) (-4.983,1.114) (-4.903,1.021) (-4.830,0.928) (-4.759,0.831) (-4.697,0.735) (-4.642,0.642) (-4.591,0.544) (-4.549,0.450) (-4.513,0.352) (-4.484,0.252) (-4.465,0.155) (-4.454,0.056) (-4.454,-0.041) (-4.463,-0.140) (-4.483,-0.237) (-4.515,-0.337) (-4.557,-0.435) (-4.610,-0.531) (-4.674,-0.624) (-4.748,-0.716) (-4.835,-0.809) (-4.933,-0.897) (-5.046,-0.984) (-5.175,-1.068) (-5.319,-1.143) (-5.484,-1.210) (-5.676,-1.262) (-5.893,-1.289) (-6.118,-1.281) (-6.327,-1.239) (-6.502,-1.176) (-6.650,-1.101) (-6.777,-1.019) (-6.885,-0.935) (-6.983,-0.844) (-7.067,-0.752) (-7.141,-0.656) (-7.202,-0.564)};
\draw[ch25arrow] (axis cs:-5.911,1.808) -- (axis cs:-5.865,1.786);
\draw[ch25arrow] (axis cs:-4.458,-0.091) -- (axis cs:-4.461,-0.115);
\addplot[ch25orbit] coordinates {(-6.900,2.950) (-6.679,2.931) (-6.461,2.898) (-6.259,2.853) (-6.062,2.795) (-5.882,2.730) (-5.720,2.662) (-5.563,2.586) (-5.410,2.505) (-5.274,2.427) (-5.143,2.346) (-5.016,2.263) (-4.894,2.179) (-4.777,2.095) (-4.664,2.012) (-4.546,1.923) (-4.433,1.836) (-4.317,1.746) (-4.206,1.659) (-4.093,1.572) (-3.978,1.485) (-3.863,1.400) (-3.741,1.314) (-3.614,1.232) (-3.477,1.152) (-3.326,1.080) (-3.148,1.020) (-2.933,0.988) (-2.710,1.002) (-2.511,1.051) (-2.338,1.113) (-2.178,1.181) (-2.019,1.253) (-1.865,1.326) (-1.708,1.399) (-1.550,1.469) (-1.385,1.537) (-1.212,1.600) (-1.025,1.657) (-0.823,1.702) (-0.609,1.730) (-0.383,1.735) (-0.159,1.714) (0.045,1.670) (0.225,1.610) (0.383,1.540) (0.526,1.463) (0.654,1.381) (0.768,1.297) (0.874,1.208) (0.973,1.116) (1.059,1.026) (1.137,0.935) (1.208,0.844) (1.276,0.747) (1.335,0.650) (1.386,0.554) (1.429,0.459) (1.466,0.361) (1.494,0.264) (1.515,0.164) (1.526,0.066) (1.527,-0.034) (1.519,-0.132) (1.500,-0.232) (1.470,-0.329) (1.429,-0.428) (1.376,-0.523) (1.311,-0.619) (1.234,-0.710) (1.144,-0.798) (1.034,-0.886) (0.909,-0.966) (0.758,-1.040) (0.576,-1.099) (0.364,-1.132) (0.137,-1.125) (-0.067,-1.080) (-0.235,-1.014) (-0.372,-0.938) (-0.489,-0.854) (-0.590,-0.765) (-0.676,-0.672) (-0.748,-0.579) (-0.809,-0.483) (-0.859,-0.385) (-0.898,-0.287) (-0.925,-0.189) (-0.941,-0.093) (-0.946,0.006) (-0.939,0.105) (-0.919,0.205) (-0.885,0.302) (-0.836,0.396) (-0.767,0.490) (-0.676,0.580) (-0.554,0.662) (-0.390,0.728) (-0.174,0.755) (0.038,0.721) (0.200,0.652) (0.323,0.570) (0.417,0.481) (0.491,0.388) (0.546,0.294) (0.585,0.199) (0.609,0.103) (0.618,0.005) (0.611,-0.091) (0.585,-0.189) (0.537,-0.284) (0.461,-0.375) (0.342,-0.457) (0.146,-0.505) (-0.064,-0.470) (-0.201,-0.393) (-0.294,-0.305) (-0.356,-0.212) (-0.395,-0.116) (-0.410,-0.020) (-0.403,0.078) (-0.368,0.173) (-0.294,0.264) (-0.141,0.335) (0.072,0.303) (0.184,0.219) (0.247,0.125) (0.274,0.029) (0.266,-0.067) (0.214,-0.162) (0.118,-0.220)};
\draw[ch25arrow] (axis cs:-5.165,2.359) -- (axis cs:-5.124,2.333);
\draw[ch25arrow] (axis cs:-0.959,1.671) -- (axis cs:-0.905,1.683);
\draw[ch25arrow] (axis cs:-0.860,-0.379) -- (axis cs:-0.872,-0.351);
\addplot[ch25orbit] coordinates {(6.900,-2.000) (6.670,-1.998) (6.451,-1.973) (6.246,-1.929) (6.066,-1.871) (5.900,-1.803) (5.750,-1.728) (5.615,-1.649) (5.486,-1.564) (5.372,-1.479) (5.265,-1.390) (5.164,-1.299) (5.070,-1.207) (4.983,-1.114) (4.903,-1.021) (4.830,-0.928) (4.759,-0.831) (4.697,-0.735) (4.642,-0.642) (4.591,-0.544) (4.549,-0.450) (4.513,-0.352) (4.484,-0.252) (4.465,-0.155) (4.454,-0.056) (4.454,0.041) (4.463,0.140) (4.483,0.237) (4.515,0.337) (4.557,0.435) (4.610,0.531) (4.674,0.624) (4.748,0.716) (4.835,0.809) (4.933,0.897) (5.046,0.984) (5.175,1.068) (5.319,1.143) (5.484,1.210) (5.676,1.262) (5.893,1.289) (6.118,1.281) (6.327,1.239) (6.502,1.176) (6.650,1.101) (6.777,1.019) (6.885,0.935) (6.983,0.844) (7.067,0.752) (7.141,0.656) (7.202,0.564)};
\draw[ch25arrow] (axis cs:5.911,-1.808) -- (axis cs:5.865,-1.786);
\draw[ch25arrow] (axis cs:4.458,0.091) -- (axis cs:4.461,0.115);
\addplot[ch25orbit] coordinates {(6.900,-2.950) (6.679,-2.931) (6.461,-2.898) (6.259,-2.853) (6.062,-2.795) (5.882,-2.730) (5.720,-2.662) (5.563,-2.586) (5.410,-2.505) (5.274,-2.427) (5.143,-2.346) (5.016,-2.263) (4.894,-2.179) (4.777,-2.095) (4.664,-2.012) (4.546,-1.923) (4.433,-1.836) (4.317,-1.746) (4.206,-1.659) (4.093,-1.572) (3.978,-1.485) (3.863,-1.400) (3.741,-1.314) (3.614,-1.232) (3.477,-1.152) (3.326,-1.080) (3.148,-1.020) (2.933,-0.988) (2.710,-1.002) (2.511,-1.051) (2.338,-1.113) (2.178,-1.181) (2.019,-1.253) (1.865,-1.326) (1.708,-1.399) (1.550,-1.469) (1.385,-1.537) (1.212,-1.600) (1.025,-1.657) (0.823,-1.702) (0.609,-1.730) (0.383,-1.735) (0.159,-1.714) (-0.045,-1.670) (-0.225,-1.610) (-0.383,-1.540) (-0.526,-1.463) (-0.654,-1.381) (-0.768,-1.297) (-0.874,-1.208) (-0.973,-1.116) (-1.059,-1.026) (-1.137,-0.935) (-1.208,-0.844) (-1.276,-0.747) (-1.335,-0.650) (-1.386,-0.554) (-1.429,-0.459) (-1.466,-0.361) (-1.494,-0.264) (-1.515,-0.164) (-1.526,-0.066) (-1.527,0.034) (-1.519,0.132) (-1.500,0.232) (-1.470,0.329) (-1.429,0.428) (-1.376,0.523) (-1.311,0.619) (-1.234,0.710) (-1.144,0.798) (-1.034,0.886) (-0.909,0.966) (-0.758,1.040) (-0.576,1.099) (-0.364,1.132) (-0.137,1.125) (0.067,1.080) (0.235,1.014) (0.372,0.938) (0.489,0.854) (0.590,0.765) (0.676,0.672) (0.748,0.579) (0.809,0.483) (0.859,0.385) (0.898,0.287) (0.925,0.189) (0.941,0.093) (0.946,-0.006) (0.939,-0.105) (0.919,-0.205) (0.885,-0.302) (0.836,-0.396) (0.767,-0.490) (0.676,-0.580) (0.554,-0.662) (0.390,-0.728) (0.174,-0.755) (-0.038,-0.721) (-0.200,-0.652) (-0.323,-0.570) (-0.417,-0.481) (-0.491,-0.388) (-0.546,-0.294) (-0.585,-0.199) (-0.609,-0.103) (-0.618,-0.005) (-0.611,0.091) (-0.585,0.189) (-0.537,0.284) (-0.461,0.375) (-0.342,0.457) (-0.146,0.505) (0.064,0.470) (0.201,0.393) (0.294,0.305) (0.356,0.212) (0.395,0.116) (0.410,0.020) (0.403,-0.078) (0.368,-0.173) (0.294,-0.264) (0.141,-0.335) (-0.072,-0.303) (-0.184,-0.219) (-0.247,-0.125) (-0.274,-0.029) (-0.266,0.067) (-0.214,0.162) (-0.118,0.220)};
\draw[ch25arrow] (axis cs:5.165,-2.359) -- (axis cs:5.124,-2.333);
\draw[ch25arrow] (axis cs:0.959,-1.671) -- (axis cs:0.905,-1.683);
\draw[ch25arrow] (axis cs:0.860,0.379) -- (axis cs:0.872,0.351);
\addplot[ch25orbit] coordinates {(-1.000,2.950) (-0.764,2.953) (-0.528,2.942) (-0.308,2.917) (-0.091,2.878) (0.109,2.827) (0.291,2.769) (0.469,2.700) (0.629,2.629) (0.784,2.551) (0.922,2.475) (1.056,2.396) (1.186,2.314) (1.311,2.230) (1.431,2.146) (1.547,2.063) (1.668,1.972) (1.783,1.884) (1.894,1.799) (2.008,1.710) (2.116,1.626) (2.227,1.540) (2.339,1.456) (2.459,1.368) (2.579,1.286) (2.709,1.203) (2.849,1.126) (3.007,1.056) (3.196,1.001) (3.419,0.983) (3.637,1.011) (3.829,1.066) (3.999,1.131) (4.156,1.200) (4.311,1.273) (4.461,1.344) (4.620,1.417) (4.780,1.487) (4.947,1.554) (5.121,1.615) (5.310,1.668) (5.513,1.709) (5.737,1.732) (5.962,1.731) (6.177,1.703) (6.379,1.653) (6.557,1.589) (6.713,1.515) (6.853,1.435) (6.978,1.350) (7.089,1.264) (7.193,1.174) (7.205,1.163)};
\draw[ch25arrow] (axis cs:-0.156,2.892) -- (axis cs:-0.102,2.880);
\draw[ch25arrow] (axis cs:2.186,1.572) -- (axis cs:2.224,1.543);
\addplot[ch25orbit] coordinates {(1.000,-2.950) (0.764,-2.953) (0.528,-2.942) (0.308,-2.917) (0.091,-2.878) (-0.109,-2.827) (-0.291,-2.769) (-0.469,-2.700) (-0.629,-2.629) (-0.784,-2.551) (-0.922,-2.475) (-1.056,-2.396) (-1.186,-2.314) (-1.311,-2.230) (-1.431,-2.146) (-1.547,-2.063) (-1.668,-1.972) (-1.783,-1.884) (-1.894,-1.799) (-2.008,-1.710) (-2.116,-1.626) (-2.227,-1.540) (-2.339,-1.456) (-2.459,-1.368) (-2.579,-1.286) (-2.709,-1.203) (-2.849,-1.126) (-3.007,-1.056) (-3.196,-1.001) (-3.419,-0.983) (-3.637,-1.011) (-3.829,-1.066) (-3.999,-1.131) (-4.156,-1.200) (-4.311,-1.273) (-4.461,-1.344) (-4.620,-1.417) (-4.780,-1.487) (-4.947,-1.554) (-5.121,-1.615) (-5.310,-1.668) (-5.513,-1.709) (-5.737,-1.732) (-5.962,-1.731) (-6.177,-1.703) (-6.379,-1.653) (-6.557,-1.589) (-6.713,-1.515) (-6.853,-1.435) (-6.978,-1.350) (-7.089,-1.264) (-7.193,-1.174) (-7.205,-1.163)};
\draw[ch25arrow] (axis cs:0.156,-2.892) -- (axis cs:0.102,-2.880);
\draw[ch25arrow] (axis cs:-2.186,-1.572) -- (axis cs:-2.224,-1.543);
\fill[black] (axis cs:-6.2832,0) circle (1.4pt);
\fill[black] (axis cs:0.0000,0) circle (1.4pt);
\fill[black] (axis cs:6.2832,0) circle (1.4pt);
\draw[black, line width=0.5pt, fill=white] (axis cs:-3.1416,0) circle (1.4pt);
\draw[black, line width=0.5pt, fill=white] (axis cs:3.1416,0) circle (1.4pt);
\end{axis}
\end{tikzpicture}}
\caption{Phase portraits of the pendulum, with the angle $\theta$ and the angular velocity $v=\theta'$. (A) Without friction, $\theta''+\sin\theta=0$: the level curves of the energy $v^2/2-\cos\theta$, with closed orbits around the stable equilibria (filled dots), the separatrices $v=\pm2\cos(\theta/2)$ (thick) through the saddles (open dots), and rotations above and below them. (B) With friction, $\theta''+0.25\,\theta'+\sin\theta=0$: solutions computed numerically; rotations slow down until the pendulum is captured by a hanging position, around which it spirals in.\label{fig:pendulum}}
\end{figure}

\subsection{Predators and Prey}

The Lotka--Volterra system \eqref{eq:lotka-volterra} has two equilibria with $u,w\ge0$. At $(0,0)$ the Jacobian matrix is $\operatorname{diag}(a,-c)$, a saddle: without predators the prey grow exponentially, and without prey the predators die out. At $(u^*,w^*)=(c/d,a/b)$ the linearisation, found in \cref{sec:ode-systems}, has the eigenvalues $\pm i\sqrt{ac}$, a centre, which the linearisation does not decide. Once more a conserved quantity decides.

\begin{proposition}[Lotka--Volterra system]
\label{prop:lotka-volterra}
Let $a,b,c,d>0$, and let $\mathbf y=(u,w)$ be a solution of \eqref{eq:lotka-volterra} with $u(0)>0$ and $w(0)>0$.
\begin{enumerate}
\item The function $H(u,w)\coloneqq du-c\ln u+bw-a\ln w$ is constant along the solution, which exists for all $t$ and stays in the open quadrant.
\item If $\mathbf y(0)\ne(u^*,w^*)$, the solution is periodic, and its orbit is a closed curve around $(u^*,w^*)$, traversed counterclockwise.
\item Over a period $T$, the averages of the populations are the equilibrium values: $\int_0^Tu=Tu^*$ and $\int_0^Tw=Tw^*$.
\end{enumerate}
\end{proposition}

\begin{proof}
(i) The axes $u=0$ and $w=0$ are unions of orbits, since $u\equiv0$ solves the first equation and $w\equiv0$ the second, so by \cref{lem:orbits}\,(ii) the solution cannot reach them, and it stays in the open quadrant, where $H$ is defined. There
\[
\begin{aligned}
\frac d{dt}H(u,w)&=\Bigl(d-\frac cu\Bigr)u(a-bw)+\Bigl(b-\frac aw\Bigr)w(-c+du)\\
&=(du-c)(a-bw)+(bw-a)(du-c)=0 .
\end{aligned}
\]
The function $H$ is the sum of $h_1(u)\coloneqq du-c\ln u$ and $h_2(w)\coloneqq bw-a\ln w$, each of which tends to infinity at $0$ and at infinity (\cref{prop:ln-growth}); so the level set of $H$ through $\mathbf y(0)$ is closed and bounded and lies in the open quadrant, and the solution, which stays on it, exists for all $t$ by \cref{thm:maximal-solution}.

(ii) Let $L$ be this level set and $\xi\coloneqq u-u^*$, $\eta\coloneqq w-w^*$. The equations read $u'=-bu\eta$ and $w'=dw\xi$, so
\[
\xi\eta'-\eta\xi'=dw\xi^2+bu\eta^2>0\qquad\text{for }(\xi,\eta)\ne\mathbf 0 .
\]
The function $H$ has the second derivatives $h_1''=c/u^2>0$ and $h_2''=a/w^2>0$, so its restriction to every line is strictly convex, and its gradient vanishes at $(u^*,w^*)$. So the equilibrium is the strict minimum of $H$, it does not lie on $L$, and on each ray from it $H$ increases strictly with the distance, so that the ray meets $L$ in at most one point. On the compact set $L$, the continuous functions $u$, $w$ and $\xi^2+\eta^2$ have positive minima, and $\xi^2+\eta^2$ is bounded; hence the polar angle about $(u^*,w^*)$ has the derivative $(dw\xi^2+bu\eta^2)/(\xi^2+\eta^2)$, which is bounded below by a positive constant, as in the proof of \cref{prop:rotation-sense}. The angle therefore increases by $2\pi$ in a finite time $T$, and then the solution is again on the ray through $\mathbf y(0)$ and on $L$, so $\mathbf y(T)=\mathbf y(0)$. By \cref{lem:orbits}\,(iii), the solution is periodic.

(iii) $(\ln w)'=w'/w=-c+du$; integrating over a period, the left side gives $\ln w(T)-\ln w(0)=0$, so $cT=d\int_0^Tu$. In the same way $(\ln u)'=a-bw$ gives $aT=b\int_0^Tw$.
\end{proof}

The populations therefore oscillate for ever, with the predators lagging behind the prey: when prey are plentiful the predators multiply, which reduces the prey, after which the predators decline and the prey recover (\cref{fig:lv-competition}A). The amplitude of the oscillation depends on the initial state, as for a centre, and a disturbance moves the system to another closed orbit rather than back to the old one. Part (iii) has a surprising consequence. If both species are harvested at a rate proportional to their size, $a$ decreases and $c$ increases by the same amount, so the average prey population $c/d$ increases and the average predator population $a/b$ decreases. The model is too simple to describe real populations in detail, but its oscillations and its average behaviour are a first explanation of the cycles observed in nature; Hirsch, Smale and Devaney~\cite{hirsch2013} and Boyce, DiPrima and Meade~\cite[Chap.~9]{boyce2017} discuss it and its refinements.

\subsection{Competing Species}

Two species that compete for the same resource, rather than eat each other, give a different picture. In the system
\begin{equation}
u'=u(3-u-2w),\qquad w'=w(2-u-w),
\label{eq:competition}
\end{equation}
each species alone grows logistically (\cref{sec:separable}) towards its carrying capacity, $3$ for $u$ and $2$ for $w$, and each is hindered by the other~\cite[Sect.~6.4]{strogatz2015}. The equilibria are $(0,0)$, $(3,0)$, $(0,2)$ and the coexistence point $(1,1)$, where both brackets vanish. The Jacobian matrix
\[
D\mathbf f(u,w)=\begin{bmatrix}3-2u-2w&-2u\\-w&2-u-2w\end{bmatrix}
\]
is $\operatorname{diag}(3,2)$ at $(0,0)$, an unstable node; it has the eigenvalues $-3$ and $-1$ at $(3,0)$, and $-1$ and $-2$ at $(0,2)$, where it is triangular, so these are stable nodes; and at $(1,1)$ it has the rows $(-1,-2)$ and $(-1,-1)$, with $\det=-1<0$, a saddle with the eigenvalues $-1\pm\sqrt2$. By \cref{thm:linearised-stability-systems,thm:linearised-instability}, the two equilibria at which one species has died out are asymptotically stable, and the coexistence point is unstable. By \cref{thm:hartman-grobman}, two orbits come into the saddle. In the computed solutions of \cref{fig:lv-competition}B, traced backwards in time, one of them comes from the neighbourhood of the origin and the other from far away, and together with the saddle they form a curve that divides the quadrant; the computed solutions that start above it end at $(0,2)$, those below it at $(3,0)$. One species dies out, and which one depends on the initial populations. This is an instance of the principle of competitive exclusion.

\begin{figure}
\sdpanel{A}{\begin{tikzpicture}
\begin{axis}[sdplot, width=0.47\textwidth, height=0.47\textwidth, xmin=0, xmax=7, ymin=0, ymax=7, clip=true, xtick={0,2,4,6}, ytick={0,2,4,6}, xlabel={$u$ (prey)}, ylabel={$w$ (predators)}]
\addplot[ch25orbit] coordinates {(3.032,2.000) (3.025,2.101) (3.005,2.200) (2.973,2.295) (2.932,2.386) (2.883,2.472) (2.828,2.553) (2.767,2.630) (2.701,2.701) (2.630,2.767) (2.553,2.828) (2.472,2.883) (2.386,2.932) (2.295,2.973) (2.200,3.005) (2.101,3.025) (2.000,3.032) (1.899,3.025) (1.801,3.001) (1.708,2.962) (1.624,2.908) (1.549,2.844) (1.484,2.772) (1.429,2.696) (1.383,2.617) (1.345,2.538) (1.314,2.459) (1.288,2.380) (1.268,2.303) (1.253,2.227) (1.242,2.151) (1.236,2.075) (1.234,2.000) (1.236,1.925) (1.243,1.849) (1.256,1.774) (1.276,1.700) (1.303,1.627) (1.337,1.557) (1.381,1.492) (1.432,1.432) (1.492,1.381) (1.557,1.337) (1.627,1.303) (1.700,1.276) (1.774,1.256) (1.849,1.243) (1.925,1.236) (2.000,1.234) (2.075,1.236) (2.151,1.242) (2.227,1.253) (2.303,1.268) (2.380,1.288) (2.459,1.314) (2.538,1.345) (2.617,1.383) (2.696,1.429) (2.772,1.484) (2.844,1.549) (2.908,1.624) (2.962,1.708) (3.001,1.801) (3.025,1.899) (3.032,2.000)};
\draw[ch25arrow] (axis cs:3.032,1.980) -- (axis cs:3.032,2.000);
\draw[ch25arrow] (axis cs:1.234,2.020) -- (axis cs:1.234,2.000);
\draw[ch25arrow] (axis cs:2.020,3.032) -- (axis cs:2.000,3.032);
\draw[ch25arrow] (axis cs:1.980,1.234) -- (axis cs:2.000,1.234);
\addplot[ch25orbit] coordinates {(3.972,2.000) (3.955,2.193) (3.907,2.379) (3.838,2.557) (3.752,2.726) (3.653,2.884) (3.545,3.032) (3.428,3.172) (3.304,3.304) (3.172,3.428) (3.032,3.545) (2.884,3.653) (2.726,3.752) (2.557,3.838) (2.379,3.907) (2.193,3.955) (2.000,3.972) (1.808,3.953) (1.624,3.890) (1.458,3.786) (1.316,3.651) (1.201,3.495) (1.109,3.333) (1.038,3.173) (0.982,3.018) (0.938,2.871) (0.904,2.732) (0.878,2.600) (0.857,2.473) (0.842,2.351) (0.831,2.233) (0.824,2.116) (0.822,2.000) (0.825,1.884) (0.833,1.768) (0.847,1.650) (0.870,1.532) (0.904,1.414) (0.951,1.299) (1.014,1.191) (1.095,1.095) (1.191,1.014) (1.299,0.951) (1.414,0.904) (1.532,0.870) (1.650,0.847) (1.768,0.833) (1.884,0.825) (2.000,0.822) (2.116,0.824) (2.233,0.831) (2.351,0.842) (2.473,0.857) (2.600,0.878) (2.732,0.904) (2.871,0.938) (3.018,0.982) (3.173,1.038) (3.333,1.109) (3.495,1.201) (3.651,1.316) (3.786,1.458) (3.890,1.624) (3.953,1.808) (3.972,2.000)};
\draw[ch25arrow] (axis cs:3.972,1.980) -- (axis cs:3.972,2.000);
\draw[ch25arrow] (axis cs:0.822,2.020) -- (axis cs:0.822,2.000);
\draw[ch25arrow] (axis cs:2.020,3.972) -- (axis cs:2.000,3.972);
\draw[ch25arrow] (axis cs:1.980,0.822) -- (axis cs:2.000,0.822);
\addplot[ch25orbit] coordinates {(5.218,2.000) (5.182,2.313) (5.089,2.615) (4.960,2.898) (4.808,3.163) (4.640,3.411) (4.460,3.644) (4.271,3.863) (4.072,4.072) (3.863,4.271) (3.644,4.460) (3.411,4.640) (3.163,4.808) (2.898,4.960) (2.615,5.089) (2.313,5.182) (2.000,5.218) (1.687,5.173) (1.397,5.030) (1.153,4.793) (0.965,4.499) (0.830,4.189) (0.734,3.894) (0.666,3.625) (0.617,3.383) (0.581,3.165) (0.554,2.966) (0.533,2.784) (0.517,2.614) (0.506,2.453) (0.498,2.299) (0.493,2.148) (0.491,2.000) (0.493,1.852) (0.499,1.701) (0.511,1.548) (0.531,1.392) (0.563,1.232) (0.612,1.072) (0.685,0.921) (0.789,0.789) (0.921,0.685) (1.072,0.612) (1.232,0.563) (1.392,0.531) (1.548,0.511) (1.701,0.499) (1.852,0.493) (2.000,0.491) (2.148,0.493) (2.299,0.498) (2.453,0.506) (2.614,0.517) (2.784,0.533) (2.966,0.554) (3.165,0.581) (3.383,0.617) (3.625,0.666) (3.894,0.734) (4.189,0.830) (4.499,0.965) (4.793,1.153) (5.030,1.397) (5.173,1.687) (5.218,2.000)};
\draw[ch25arrow] (axis cs:5.218,1.980) -- (axis cs:5.218,2.000);
\draw[ch25arrow] (axis cs:0.491,2.020) -- (axis cs:0.491,2.000);
\draw[ch25arrow] (axis cs:2.020,5.218) -- (axis cs:2.000,5.218);
\draw[ch25arrow] (axis cs:1.980,0.491) -- (axis cs:2.000,0.491);
\addplot[ch25orbit] coordinates {(6.583,2.000) (6.520,2.445) (6.369,2.869) (6.168,3.264) (5.938,3.631) (5.691,3.973) (5.433,4.294) (5.165,4.597) (4.887,4.887) (4.597,5.165) (4.294,5.433) (3.973,5.691) (3.631,5.938) (3.264,6.168) (2.869,6.369) (2.445,6.520) (2.000,6.583) (1.557,6.499) (1.161,6.220) (0.857,5.767) (0.657,5.243) (0.532,4.747) (0.453,4.315) (0.402,3.948) (0.366,3.634) (0.341,3.361) (0.323,3.120) (0.309,2.904) (0.299,2.705) (0.292,2.518) (0.286,2.341) (0.283,2.169) (0.282,2.000) (0.283,1.831) (0.287,1.659) (0.296,1.483) (0.310,1.300) (0.334,1.110) (0.376,0.915) (0.446,0.725) (0.562,0.562) (0.725,0.446) (0.915,0.376) (1.110,0.334) (1.300,0.310) (1.483,0.296) (1.659,0.287) (1.831,0.283) (2.000,0.282) (2.169,0.283) (2.341,0.286) (2.518,0.292) (2.705,0.299) (2.904,0.309) (3.120,0.323) (3.361,0.341) (3.634,0.366) (3.948,0.402) (4.315,0.453) (4.747,0.532) (5.243,0.657) (5.767,0.857) (6.220,1.161) (6.499,1.557) (6.583,2.000)};
\draw[ch25arrow] (axis cs:6.583,1.980) -- (axis cs:6.583,2.000);
\draw[ch25arrow] (axis cs:0.282,2.020) -- (axis cs:0.282,2.000);
\draw[ch25arrow] (axis cs:2.020,6.583) -- (axis cs:2.000,6.583);
\draw[ch25arrow] (axis cs:1.980,0.282) -- (axis cs:2.000,0.282);
\fill[black] (axis cs:2,2) circle (1.4pt);
\draw[black, line width=0.5pt, fill=white] (axis cs:0,0) circle (1.4pt);
\end{axis}
\end{tikzpicture}}\hfill
\sdpanel{B}{\begin{tikzpicture}
\begin{axis}[sdplot, width=0.47\textwidth, height=0.47\textwidth, xmin=0, xmax=3.5, ymin=0, ymax=2.5, clip=true, xtick={0,1,2,3}, ytick={0,1,2}, xlabel={$u$}, ylabel={$w$}]
\addplot[black, line width=0.9pt, dashed] coordinates {(1.000,1.000) (1.050,1.035) (1.101,1.070) (1.152,1.105) (1.203,1.140) (1.256,1.176) (1.311,1.212) (1.367,1.248) (1.424,1.285) (1.479,1.321) (1.536,1.357) (1.594,1.393) (1.650,1.428) (1.713,1.467) (1.773,1.503) (1.840,1.544) (1.900,1.580) (1.965,1.619) (2.020,1.651) (2.078,1.685) (2.142,1.722) (2.211,1.762) (2.287,1.806) (2.371,1.853) (2.431,1.887) (2.496,1.923) (2.565,1.962) (2.639,2.003) (2.719,2.046) (2.805,2.093) (2.898,2.144) (2.999,2.198) (3.109,2.256) (3.168,2.287) (3.230,2.319) (3.294,2.353) (3.362,2.388) (3.432,2.425) (3.507,2.463) (3.585,2.503)};
\draw[ch25arrow] (axis cs:2.279,1.799) -- (axis cs:2.267,1.794);
\addplot[black, line width=0.9pt, dashed] coordinates {(1.000,1.000) (0.950,0.965) (0.900,0.929) (0.851,0.892) (0.802,0.856) (0.753,0.819) (0.704,0.781) (0.656,0.744) (0.608,0.704) (0.561,0.666) (0.514,0.626) (0.468,0.587) (0.422,0.547) (0.378,0.506) (0.335,0.466) (0.292,0.424) (0.250,0.381) (0.211,0.338) (0.173,0.295) (0.136,0.250) (0.102,0.205) (0.070,0.160) (0.043,0.114) (0.019,0.066) (0.006,0.029)};
\draw[ch25arrow] (axis cs:0.436,0.561) -- (axis cs:0.446,0.568);
\addplot[ch25orbit] coordinates {(1.000,1.000) (0.950,1.035) (0.901,1.071) (0.853,1.108) (0.805,1.144) (0.758,1.181) (0.711,1.219) (0.665,1.257) (0.619,1.295) (0.574,1.334) (0.530,1.373) (0.487,1.413) (0.444,1.453) (0.402,1.494) (0.362,1.535) (0.322,1.576) (0.283,1.618) (0.244,1.661) (0.208,1.704) (0.172,1.747) (0.137,1.791) (0.104,1.835) (0.073,1.880) (0.043,1.925) (0.015,1.971) (0.014,1.974)};
\draw[ch25arrow] (axis cs:0.469,1.432) -- (axis cs:0.459,1.439);
\addplot[ch25orbit] coordinates {(1.000,1.000) (1.050,0.965) (1.101,0.930) (1.152,0.896) (1.204,0.862) (1.256,0.828) (1.308,0.795) (1.362,0.762) (1.415,0.730) (1.469,0.698) (1.524,0.666) (1.579,0.634) (1.635,0.603) (1.691,0.572) (1.747,0.542) (1.804,0.512) (1.861,0.483) (1.919,0.453) (1.977,0.424) (2.036,0.395) (2.094,0.368) (2.153,0.340) (2.212,0.313) (2.272,0.286) (2.333,0.259) (2.394,0.232) (2.455,0.206) (2.516,0.181) (2.577,0.156) (2.639,0.132) (2.701,0.107) (2.764,0.084) (2.826,0.061) (2.889,0.038) (2.952,0.016) (2.972,0.010)};
\draw[ch25arrow] (axis cs:1.939,0.441) -- (axis cs:1.952,0.437);
\addplot[ch25orbit] coordinates {(0.050,0.500) (0.058,0.557) (0.066,0.614) (0.073,0.670) (0.081,0.728) (0.089,0.784) (0.097,0.841) (0.104,0.899) (0.111,0.956) (0.118,1.013) (0.124,1.068) (0.130,1.126) (0.135,1.182) (0.139,1.239) (0.143,1.296) (0.145,1.352) (0.146,1.410) (0.146,1.465) (0.144,1.521) (0.140,1.576) (0.134,1.632) (0.124,1.688) (0.112,1.743) (0.096,1.797) (0.075,1.850) (0.051,1.902) (0.024,1.954) (0.014,1.974)};
\draw[ch25arrow] (axis cs:0.117,1.011) -- (axis cs:0.119,1.021);
\addplot[ch25orbit] coordinates {(0.500,0.050) (0.581,0.055) (0.660,0.059) (0.739,0.063) (0.817,0.066) (0.900,0.070) (0.979,0.073) (1.061,0.076) (1.146,0.079) (1.224,0.081) (1.304,0.083) (1.384,0.085) (1.464,0.087) (1.543,0.088) (1.621,0.090) (1.698,0.090) (1.781,0.091) (1.861,0.091) (1.944,0.091) (2.024,0.091) (2.104,0.090) (2.185,0.089) (2.265,0.088) (2.342,0.085) (2.419,0.083) (2.498,0.079) (2.576,0.074) (2.654,0.068) (2.730,0.060) (2.807,0.050) (2.881,0.035) (2.953,0.015) (2.972,0.009)};
\draw[ch25arrow] (axis cs:1.353,0.085) -- (axis cs:1.367,0.085);
\addplot[ch25orbit] coordinates {(0.150,0.120) (0.214,0.152) (0.281,0.181) (0.350,0.209) (0.420,0.235) (0.491,0.260) (0.564,0.283) (0.638,0.305) (0.712,0.326) (0.787,0.345) (0.865,0.364) (0.943,0.381) (1.021,0.397) (1.097,0.411) (1.177,0.423) (1.253,0.434) (1.331,0.442) (1.408,0.448) (1.487,0.452) (1.566,0.453) (1.645,0.450) (1.721,0.444) (1.798,0.434) (1.873,0.419) (1.947,0.400) (2.019,0.378) (2.089,0.353) (2.157,0.327) (2.226,0.300) (2.292,0.273) (2.359,0.245) (2.426,0.217) (2.493,0.190) (2.560,0.163) (2.628,0.136) (2.696,0.109) (2.765,0.083) (2.834,0.058) (2.903,0.033) (2.972,0.010)};
\draw[ch25arrow] (axis cs:1.121,0.415) -- (axis cs:1.134,0.417);
\addplot[ch25orbit] coordinates {(3.450,0.700) (3.376,0.683) (3.287,0.661) (3.207,0.641) (3.134,0.623) (3.053,0.602) (2.980,0.583) (2.903,0.562) (2.826,0.540) (2.751,0.518) (2.679,0.496) (2.609,0.472) (2.538,0.445) (2.472,0.416) (2.412,0.382) (2.370,0.335) (2.381,0.280) (2.430,0.237) (2.489,0.202) (2.554,0.171) (2.620,0.141) (2.687,0.114) (2.755,0.088) (2.823,0.062) (2.892,0.037) (2.961,0.013) (2.972,0.010)};
\draw[ch25arrow] (axis cs:2.832,0.541) -- (axis cs:2.819,0.538);
\addplot[ch25orbit] coordinates {(3.450,1.600) (3.353,1.562) (3.263,1.527) (3.180,1.494) (3.102,1.463) (3.029,1.434) (2.961,1.407) (2.867,1.368) (2.782,1.333) (2.704,1.300) (2.632,1.270) (2.545,1.233) (2.468,1.200) (2.398,1.169) (2.319,1.135) (2.250,1.103) (2.176,1.069) (2.112,1.039) (2.046,1.008) (1.982,0.976) (1.915,0.941) (1.849,0.905) (1.789,0.870) (1.731,0.833) (1.677,0.793) (1.633,0.747) (1.613,0.694) (1.634,0.640) (1.681,0.596) (1.737,0.558) (1.796,0.522) (1.857,0.488) (1.919,0.455) (1.982,0.423) (2.046,0.391) (2.111,0.360) (2.176,0.330) (2.242,0.299) (2.307,0.270) (2.374,0.241) (2.441,0.212) (2.509,0.184) (2.577,0.156) (2.645,0.129) (2.713,0.103) (2.782,0.077) (2.851,0.052) (2.920,0.027) (2.972,0.010)};
\draw[ch25arrow] (axis cs:2.293,1.121) -- (axis cs:2.280,1.117);
\addplot[ch25orbit] coordinates {(2.400,2.450) (2.321,2.396) (2.247,2.346) (2.179,2.299) (2.115,2.255) (2.056,2.213) (2.001,2.174) (1.924,2.120) (1.854,2.070) (1.790,2.024) (1.732,1.982) (1.661,1.930) (1.598,1.883) (1.540,1.841) (1.476,1.792) (1.418,1.749) (1.357,1.703) (1.303,1.662) (1.247,1.620) (1.192,1.579) (1.138,1.538) (1.082,1.496) (1.027,1.456) (0.972,1.416) (0.914,1.376) (0.855,1.340) (0.791,1.308) (0.717,1.290) (0.642,1.305) (0.582,1.339) (0.529,1.380) (0.480,1.422) (0.432,1.466) (0.386,1.511) (0.341,1.556) (0.298,1.602) (0.255,1.649) (0.215,1.695) (0.175,1.743) (0.137,1.791) (0.101,1.840) (0.067,1.889) (0.034,1.939) (0.014,1.974)};
\draw[ch25arrow] (axis cs:1.540,1.838) -- (axis cs:1.529,1.832);
\addplot[ch25orbit] coordinates {(1.200,2.450) (1.157,2.403) (1.105,2.346) (1.058,2.294) (1.015,2.247) (0.967,2.193) (0.924,2.145) (0.878,2.093) (0.830,2.040) (0.789,1.993) (0.747,1.947) (0.701,1.897) (0.655,1.848) (0.610,1.801) (0.562,1.755) (0.512,1.712) (0.457,1.673) (0.390,1.645) (0.313,1.649) (0.250,1.683) (0.201,1.725) (0.158,1.771) (0.118,1.819) (0.082,1.867) (0.049,1.917) (0.018,1.967) (0.014,1.974)};
\draw[ch25arrow] (axis cs:0.817,2.022) -- (axis cs:0.808,2.014);
\addplot[ch25orbit] coordinates {(0.350,2.450) (0.326,2.391) (0.303,2.334) (0.281,2.279) (0.258,2.223) (0.236,2.168) (0.213,2.115) (0.188,2.062) (0.162,2.010) (0.129,1.960) (0.070,1.924) (0.020,1.966) (0.015,1.974)};
\draw[ch25arrow] (axis cs:0.266,2.238) -- (axis cs:0.261,2.228);
\fill[black] (axis cs:3,0) circle (1.4pt);
\fill[black] (axis cs:0,2) circle (1.4pt);
\draw[black, line width=0.5pt, fill=white] (axis cs:0,0) circle (1.4pt);
\draw[black, line width=0.5pt, fill=white] (axis cs:1,1) circle (1.4pt);
\end{axis}
\end{tikzpicture}}
\caption{(A) Orbits of the Lotka--Volterra system $u'=u-0.5uw$, $w'=-w+0.5uw$, the level curves $H=H(2,2)+\eta$ for $\eta=0.1$, $0.3$, $0.65$ and $1.1$ of the conserved quantity $H=0.5u-\ln u+0.5w-\ln w$, around the equilibrium $(2,2)$; they are traversed counterclockwise. (B) The competition model \eqref{eq:competition}: stable nodes at $(3,0)$ and $(0,2)$ (filled), unstable equilibria at $(0,0)$ and $(1,1)$ (open), solutions computed numerically, and the stable manifold of the saddle $(1,1)$ (dashed), which separates the initial states that lead to $(3,0)$ from those that lead to $(0,2)$.\label{fig:lv-competition}}
\end{figure}

\subsection{Conservative and Gradient Systems}

The pendulum without friction and the competing species show two opposite kinds of behaviour: orbits that keep a conserved quantity and return for ever, and orbits that settle at equilibria. Two classes of systems show these behaviours in pure form, and they correspond to the two kinds of potentials met earlier: the conservative fields of \cref{sec:conservative} and the potentials $V$ with $y'=-V'(y)$ of \cref{sec:autonomous}.

\begin{proposition}[Gradient and Hamiltonian systems]
\label{prop:gradient-hamiltonian}
\leavevmode
\begin{enumerate}
\item Let $V\colon U\to\mathbb R$ be twice continuously differentiable on an open $U\subseteq\mathbb R^n$, and consider the \term{gradient system} $\mathbf y'=-\nabla V(\mathbf y)$. Along every solution, $V(\mathbf y(t))$ decreases strictly unless the solution is an equilibrium; the equilibria are the critical points of $V$; the linearisation at an equilibrium has only real eigenvalues; and a critical point at which the Hessian $\nabla^2V$ is positive definite is asymptotically stable.
\item Let $H\colon U\to\mathbb R$ be twice continuously differentiable on an open $U\subseteq\mathbb R^2$, and consider the \term{Hamiltonian system} $\theta'=\partial H/\partial v$, $v'=-\partial H/\partial\theta$. Then $H$ is constant along every solution, the linearisation at every equilibrium has the trace $0$, and an equilibrium is never asymptotically stable unless $H$ is constant on a neighbourhood of it.
\end{enumerate}
\end{proposition}

\begin{proof}
(i) By the chain rule, $(V\circ\mathbf y)'=\langle\nabla V(\mathbf y),\mathbf y'\rangle=-\|\nabla V(\mathbf y)\|^2$, which is negative unless $\nabla V(\mathbf y(t))=\mathbf 0$, and a solution that passes through an equilibrium is constant (\cref{lem:orbits}). The derivative of $-\nabla V$ is $-\nabla^2V$, a symmetric matrix by \cref{thm:schwarz}, whose eigenvalues are real (\cref{prop:hermitian-eigen}); if $\nabla^2V(\mathbf y^*)$ is positive definite, they are negative, and \cref{thm:linearised-stability-systems} applies. (ii) $(H\circ\mathbf y)'=H_\theta\theta'+H_vv'=H_\theta H_v-H_vH_\theta=0$. The Jacobian matrix has the rows $(H_{v\theta},H_{vv})$ and $(-H_{\theta\theta},-H_{\theta v})$, whose trace vanishes by \cref{thm:schwarz}. If an equilibrium $\mathbf y^*$ were asymptotically stable, every solution starting in a neighbourhood of it would tend to $\mathbf y^*$, and since $H$ is constant along it and continuous, $H(\mathbf y(0))=H(\mathbf y^*)$ for all initial values in that neighbourhood.
\end{proof}

A gradient system models motion dominated by friction, a ball rolling in very viscous oil, that slides downhill on the graph of $V$ without overshooting: it has no closed orbits and no spirals. A Hamiltonian system models motion without friction. The pendulum \eqref{eq:pendulum} is Hamiltonian with $H=E=v^2/2-\cos\theta$, and its linearisations, with the trace $0$, are centres or saddles, the eigenvalues being $\pm\lambda$ or $\pm i\beta$. The Lotka--Volterra system becomes Hamiltonian in the logarithmic variables $p=\ln u$ and $q=\ln w$, with the conserved quantity $H$ of \cref{prop:lotka-volterra}. Linear systems show the same dichotomy: a symmetric $A$ gives a gradient system (\cref{sec:eigenvalue-method}), while the vibrating system \eqref{eq:vibrations} is of the frictionless kind, with the energy as conserved quantity and the eigenvalues $\pm i\omega_j$ on the imaginary axis.

\subsection{Limit Cycles}

The closed orbits of the pendulum and of the Lotka--Volterra system come in families, one through each point near the centre. Nonlinear systems can also have an isolated closed orbit, a \term{limit cycle}, towards which the neighbouring orbits spiral, so that the system oscillates with an amplitude fixed by the system rather than by the initial state. The heartbeat and the oscillations of electronic circuits are examples. Van der Pol's equation $y''-\mu(1-y^2)y'+y=0$, with $\mu>0$, models an oscillator that is pushed by negative damping when $|y|<1$ and damped when $|y|>1$~\cite{vanderpol1926}. It has exactly one limit cycle, which attracts every other nonconstant solution; Strogatz~\cite[Sect.~7.4]{strogatz2015} derives this from a theorem of Li\'enard. In the plane, a solution that stays bounded and away from the equilibria must approach a closed orbit, by the following theorem. For a solution defined for $t\ge0$, its $\omega$-\term{limit set} consists of the points $\mathbf p$ such that $\mathbf y(t_k)\to\mathbf p$ for some sequence $t_k\to\infty$.

\begin{theorem}[Poincar\'e--Bendixson]
\label{thm:poincare-bendixson}
Let $\mathbf f\colon U\to\mathbb R^2$ be continuously differentiable on an open $U\subseteq\mathbb R^2$, and let $\mathbf y$ be a solution of $\mathbf y'=\mathbf f(\mathbf y)$ that stays, for all $t\ge0$, in a compact set $K\subseteq U$ containing no equilibrium. Then the $\omega$-limit set of $\mathbf y$ is the orbit of a periodic solution: either $\mathbf y$ itself is periodic, or its orbit spirals towards a closed orbit.
\end{theorem}

\begin{proof}[sketch]
The $\omega$-limit set $\Omega$ is nonempty, by the compactness of $K$, and it is closed, bounded and a union of orbits, by the continuous dependence of solutions on their initial values. Through a point $\mathbf p\in\Omega$, which is not an equilibrium, draw a short segment $S$ transversal to the vector field. The orbit of $\mathbf y$ crosses $S$ infinitely often, and since in the plane a curve formed by an arc of an orbit and a piece of $S$ separates the plane, by the Jordan curve theorem, the successive crossings move monotonically along $S$. Hence $\Omega$ meets $S$ in the single point $\mathbf p$. Applied to a point of the $\omega$-limit set of the orbit through $\mathbf p$, which lies in $\Omega$, the same argument shows that this orbit returns to $S$ at $\mathbf p$, so it is closed, and that $\Omega$ is this closed orbit. The theorem goes back to Poincar\'e and Bendixson~\cite{poincare1881,bendixson1901}; complete proofs are in Hartman~\cite[Chap.~VII]{hartman2002} and Perko~\cite[Chap.~3]{perko2001}.
\end{proof}

The theorem is used by finding a \term{trapping region}, a compact region without equilibria that solutions enter but cannot leave, such as an annulus whose boundary the vector field crosses inwards; the theorem then guarantees a closed orbit inside. The proof uses the topology of the plane, through the Jordan curve theorem, and the theorem fails in higher dimensions.

\begin{note}{Chaos}
In three dimensions, bounded solutions need not approach equilibria or closed orbits. Lorenz found in 1963, in a system of three equations derived from a model of convection in the atmosphere, $y_1'=\sigma(y_2-y_1)$, $y_2'=ry_1-y_2-y_1y_3$, $y_3'=y_1y_2-by_3$, with $\sigma=10$, $b=8/3$ and $r=28$, that the solutions wander for ever on a bounded set of complicated structure, a \emph{strange attractor}, and that solutions starting very close together separate exponentially fast, so that their long-term behaviour cannot be predicted from imprecise initial data~\cite{lorenz1963}. This sensitive dependence on initial conditions, together with the bounded wandering, is called \emph{chaos}. It does not contradict the uniqueness theorem, which guarantees that each initial value determines its solution, but it shows that differences of solutions can in fact grow exponentially, as the bound $e^{L|t-t_0|}$ of \cref{thm:uniqueness} permits. Strogatz~\cite[Chap.~9]{strogatz2015} gives an introduction.
\end{note}

\section{Solving Systems Numerically: A First Look}
\label{sec:systems-numerics}

\subsection{Euler's Method for Systems}

Euler's method of \cref{sec:euler} applies to a system $\mathbf y'=\mathbf f(t,\mathbf y)$ without change: $\mathbf y_{k+1}=\mathbf y_k+h\,\mathbf f(t_k,\mathbf y_k)$, with vectors in place of numbers. Its convergence theorem (\cref{thm:euler-convergence}) and the discrete Gr\"onwall inequality on which it rests carry over when absolute values are replaced by norms, and the global error is again of order $h$~\cite[Chap.~II]{hairer1993}. For a linear system $\mathbf y'=A\mathbf y$ the method gives $\mathbf y_k=(I+hA)^k\mathbf y_0$, and with $h=1/m$ the value after $m$ steps is $(I+A/m)^m\mathbf y_0$, which converges to $e^A\mathbf y_0$ by \cref{thm:exp-limit-matrix}. As for a single equation, convergence as $h\to0$ is not the whole story, and the linear systems show what else matters.

\subsection{Regions of Stability}

\begin{proposition}[Stability of Euler's methods]
\label{prop:euler-stability}
Let $A\in\mathbb C^{n\times n}$ and $h>0$.
\begin{enumerate}
\item The iterates $(I+hA)^k\mathbf y_0$ of Euler's method tend to $\mathbf 0$ for every $\mathbf y_0$ if and only if $|1+h\lambda|<1$ for every eigenvalue $\lambda$ of $A$.
\item If $1/h$ is not an eigenvalue of $A$, the backward Euler method $\mathbf y_{k+1}=\mathbf y_k+hA\mathbf y_{k+1}$ gives $\mathbf y_k=(I-hA)^{-k}\mathbf y_0$, and these iterates tend to $\mathbf 0$ for every $\mathbf y_0$ if and only if $|1-h\lambda|>1$ for every eigenvalue $\lambda$. In particular, if $\alpha(A)<0$ this holds for every $h>0$.
\end{enumerate}
\end{proposition}

\begin{proof}
(i) By \cref{thm:spectral-mapping}, the eigenvalues of $I+hA$ are the numbers $1+h\lambda$, and by \cref{thm:powers-size}\,(i) its powers tend to zero if and only if all of them have modulus less than $1$. (ii) The equation for $\mathbf y_{k+1}$ is $(I-hA)\mathbf y_{k+1}=\mathbf y_k$, and $I-hA$ is invertible because its eigenvalues $1-h\lambda$ are not zero. The eigenvalues of $(I-hA)^{-1}$ are the reciprocals $1/(1-h\lambda)$, since $B\mathbf v=\mu\mathbf v$ with $\mu\ne0$ is equivalent to $B^{-1}\mathbf v=\mathbf v/\mu$, and \cref{thm:powers-size}\,(i) applies again. If $\operatorname{Re}\lambda<0$, then $|1-h\lambda|\ge\operatorname{Re}(1-h\lambda)=1-h\operatorname{Re}\lambda>1$.
\end{proof}

The set of complex numbers $z=h\lambda$ for which a method applied to the test equation $y'=\lambda y$ produces decaying iterates is its \term{region of stability}. For Euler's method it is the open disc $|1+z|<1$ of radius $1$ about $-1$, and for the backward Euler method the exterior $|1-z|>1$ of the closed disc of radius $1$ about $1$ (\cref{fig:stability-regions}). The exact solutions decay exactly for $\operatorname{Re}z<0$, the open left half-plane. The proposition says that for a system the condition must hold for $z=h\lambda$ for \emph{every} eigenvalue $\lambda$ at once: Euler's method reproduces the decay of the solutions of a stable system only if the step brings all the numbers $h\lambda$ into the disc. For a real negative eigenvalue this is the condition $h<2/|\lambda|$ of \cref{sec:euler}. For an eigenvalue on the imaginary axis, $|1+ih\omega|^2=1+h^2\omega^2>1$ for every $h>0$: Euler's method is never stable for an undamped oscillation, and the factor $1+h^2\omega^2$ is the factor by which \cref{prop:discrete-energy} found the energy to grow in each step, while the backward Euler method, with $|1-ih\omega|^{-2}=1/(1+h^2\omega^2)$, damps it. The recurrences $\mathbf y_{k+1}=(I+hA)\mathbf y_k$ are linear difference equations, whose theory is the subject of \cref{ch:difference-equations}.

\begin{figure}
\sdpanel{A}{\begin{tikzpicture}
\begin{axis}[sdplot, width=0.47\textwidth, height=0.33\textwidth, xmin=-3, xmax=3, ymin=-2, ymax=2, axis equal image,
  xtick={-2,-1,0,1,2}, ytick={-1,0,1}, xlabel={$\operatorname{Re}z$}, ylabel={$\operatorname{Im}z$}]
\fill[sdfill] (axis cs:-1,0) circle[radius=1];
\draw[black, line width=0.5pt] (axis cs:-1,0) circle[radius=1];
\draw[black!45, line width=0.4pt] (axis cs:-3,0) -- (axis cs:3,0) (axis cs:0,-2) -- (axis cs:0,2);
\addplot[only marks, mark=x, mark size=2.2pt, black] coordinates {(-2.128,0) (-0.0021,0)};
\end{axis}
\end{tikzpicture}}\hfill
\sdpanel{B}{\begin{tikzpicture}
\begin{axis}[sdplot, width=0.47\textwidth, height=0.33\textwidth, xmin=-3, xmax=3, ymin=-2, ymax=2, axis equal image,
  xtick={-2,-1,0,1,2}, ytick={-1,0,1}, xlabel={$\operatorname{Re}z$}, ylabel={$\operatorname{Im}z$}]
\fill[sdfill] (axis cs:-3,-2) rectangle (axis cs:3,2);
\fill[white] (axis cs:1,0) circle[radius=1];
\draw[black, line width=0.5pt] (axis cs:1,0) circle[radius=1];
\draw[black!45, line width=0.4pt] (axis cs:-3,0) -- (axis cs:3,0) (axis cs:0,-2) -- (axis cs:0,2);
\addplot[only marks, mark=x, mark size=2.2pt, black] coordinates {(-2.128,0) (-0.0021,0)};
\end{axis}
\end{tikzpicture}}
\caption{Regions of stability (shaded) in the plane of $z=h\lambda$. (A) Euler's method, the disc $|1+z|<1$. (B) The backward Euler method, the exterior of the disc $|1-z|\le1$, which contains the whole left half-plane. The crosses mark the numbers $h\lambda$ for the stiff system of \cref{tab:stiff} with $h=1/470$ and the eigenvalues $-1$ and $-1000$: the second lies outside the region of Euler's method.\label{fig:stability-regions}}
\end{figure}

\subsection{Stiff Systems}

A system is stiff when its matrix has eigenvalues of very different sizes, all with negative real parts, so that rapidly decaying components accompany slowly varying ones. The system $\mathbf y'=A\mathbf y$ with
\[
A=\begin{bmatrix}-500.5&499.5\\499.5&-500.5\end{bmatrix},
\]
whose eigenvalues are $-1$ and $-1000$ with the eigenvectors $(1,1)$ and $(1,-1)$, and the initial value $\mathbf y_0=\mathbf e_1$, has the solution
\[
\mathbf y(t)=\frac{e^{-t}}2\begin{bmatrix}1\\1\end{bmatrix}+\frac{e^{-1000t}}2\begin{bmatrix}1\\-1\end{bmatrix},
\]
by \cref{thm:diagonalisable-systems}. After a hundredth of a unit of time the second term is negligible, and the solution varies on the time scale $1$. Nevertheless Euler's method is stable only for $h<2/1000=0.002$, by \cref{prop:euler-stability}. \Cref{tab:stiff} shows the relative error at $t=1$ of three methods with $n$ steps of length $h=1/n$, computed in binary64 with plain Python on the machine described in \cref{sec:systems-practice}: Euler's method, the backward Euler method, whose $2\times2$ systems were solved by Cramer's rule, and the repeated multiplication by the matrix $e^{hA}$, computed once by \code{linalg.expm}, which advances the exact solution by the time $h$. For $h>0.002$, the errors of Euler's method agree to four digits with the value $|1-1000h|^n\,\|(1,-1)/2\|/\|\mathbf y(1)\|$, the growth of the fast mode multiplied by the factor $1-1000h$ in each step; for $h<0.002$ its errors and those of the backward Euler method are both close to $h/2$, and decrease in proportion to $h$. The backward Euler method gave decaying approximations for every step, with the error $0.048$ already for $h=0.1$. The exponential steps have errors between $3.0\times10^{-14}$ and $3.7\times10^{-13}$, the largest for $10\,000$ steps.

\begin{table}
\sdfloat{%
\begin{sdtable}{colspec={rllll}, row{1}={font=\small\bfseries}}
\toprule
$n$ & $h$ & Euler & backward Euler & $e^{hA}$ \\
\midrule
$10$ & $0.1$ & $2.5\times10^{20}$ & $4.8\times10^{-2}$ & $4.1\times10^{-14}$ \\
$100$ & $0.01$ & $7.2\times10^{95}$ & $5.0\times10^{-3}$ & $3.0\times10^{-14}$ \\
$470$ & $0.00213$ & $9.1\times10^{24}$ & $1.1\times10^{-3}$ & $6.4\times10^{-14}$ \\
$530$ & $0.00189$ & $9.4\times10^{-4}$ & $9.4\times10^{-4}$ & $6.5\times10^{-14}$ \\
$1000$ & $0.001$ & $5.0\times10^{-4}$ & $5.0\times10^{-4}$ & $8.0\times10^{-14}$ \\
$10\,000$ & $0.0001$ & $5.0\times10^{-5}$ & $5.0\times10^{-5}$ & $3.7\times10^{-13}$ \\
\bottomrule
\end{sdtable}}{The stiff system with the eigenvalues $-1$ and $-1000$ of \cref{sec:systems-numerics}, $\mathbf y(0)=\mathbf e_1$: the relative error $\|\mathbf y_n-\mathbf y(1)\|/\|\mathbf y(1)\|$ after $n$ steps of length $h=1/n$ of Euler's method, the backward Euler method, and repeated multiplication by $e^{hA}$ computed with \code{linalg.expm}. Euler's method is stable only for $h<0.002$.\label{tab:stiff}}
\end{table}

The table shows the dilemma of stiff systems. The fast mode has decayed after a few hundredths of a unit of time, yet it dictates the step of Euler's method for the whole computation, $1000$ steps or more where $10$ would describe the slow mode. The backward Euler method has no such restriction, at the price of solving a linear system in each step, and for linear systems with constant coefficients the exponential $e^{hA}$ removes the error of the method altogether, leaving only rounding errors. For nonlinear systems the exponential is not available, and the methods for stiff systems combine implicit steps with the linearisations of \cref{sec:nonlinear-systems}; Hairer and Wanner~\cite{hairer1996stiff} describe them.

\subsection{Measuring the Period of the Pendulum}

\Cref{prop:pendulum} gives the period of the pendulum as an elliptic integral, and a numerical solution can measure it as an experimenter would, by timing the swings. We integrated the system \eqref{eq:pendulum} from $(\theta_0,0)$ with Heun's method of \cref{sec:euler}, applied to the system, with the steps $h=0.01$ and $h=0.005$, located the times at which $\theta$ changes from positive to negative by linear interpolation between the steps, and took one tenth of the time between the first such crossing and the eleventh as the measured period. The exact period $4K(\sin(\theta_0/2))$ was computed with the ball arithmetic of python-flint (version 0.9.0), with an error bound below $10^{-10}$. \Cref{tab:pendulum-period} shows the results. The period increases slowly at first, by $1.7$ per cent at the amplitude of $30$ degrees, where the approximation $2\pi(1+\theta_0^2/16)$ is accurate to $0.03$ per cent; at $90$ degrees it is $18$ per cent longer than $2\pi$, and at $179$ degrees almost four times as long. The measured periods agree with the exact ones to within $2\times10^{-5}$ for amplitudes up to $150$ degrees with $h=0.01$. Halving the step reduced the deviations by factors between $3.5$ and $4.0$ for amplitudes up to $90$ degrees and by factors between $5.8$ and $8.3$ for the larger ones; at $179$ degrees the deviations, $5.5\times10^{-3}$ and $6.6\times10^{-4}$, are the largest in the table.

\begin{table}
\sdfloat{%
\begin{sdtable}{colspec={rlllll}, row{1}={font=\small\bfseries}}
\toprule
$\theta_0$ & $T$ & $T/(2\pi)$ & $1+\theta_0^2/16$ & deviation, $h=0.01$ & deviation, $h=0.005$ \\
\midrule
$10^\circ$ & $6.29517$ & $1.00191$ & $1.00190$ & $-1.6\times10^{-5}$ & $-4.1\times10^{-6}$ \\
$30^\circ$ & $6.39257$ & $1.01741$ & $1.01713$ & $-1.5\times10^{-5}$ & $-3.8\times10^{-6}$ \\
$60^\circ$ & $6.74300$ & $1.07318$ & $1.06854$ & $-1.1\times10^{-5}$ & $-2.8\times10^{-6}$ \\
$90^\circ$ & $7.41630$ & $1.18034$ & $1.15421$ & $-4.5\times10^{-6}$ & $-1.3\times10^{-6}$ \\
$120^\circ$ & $8.62606$ & $1.37288$ & $1.27416$ & $3.4\times10^{-6}$ & $4.9\times10^{-7}$ \\
$150^\circ$ & $11.0723$ & $1.76220$ & $1.42837$ & $1.7\times10^{-5}$ & $3.0\times10^{-6}$ \\
$170^\circ$ & $15.3270$ & $2.43936$ & $1.55022$ & $9.1\times10^{-5}$ & $1.3\times10^{-5}$ \\
$179^\circ$ & $24.5111$ & $3.90107$ & $1.61002$ & $5.5\times10^{-3}$ & $6.6\times10^{-4}$ \\
\bottomrule
\end{sdtable}}{The period $T=4K(\sin(\theta_0/2))$ of the pendulum $\theta''+\sin\theta=0$ with the amplitude $\theta_0$ (\cref{prop:pendulum}), compared with the period $2\pi$ of the linearisation and with the approximation $2\pi(1+\theta_0^2/16)$, and the relative deviation of the period measured from a numerical solution by Heun's method with the step $h$.\label{tab:pendulum-period}}
\end{table}

The numerical solution of differential equations, with the methods of Runge and Kutta, multistep methods, adaptive choice of the step and the methods for stiff systems, is the subject of the books of Hairer, N{\o}rsett and Wanner~\cite{hairer1993}, of Hairer and Wanner~\cite{hairer1996stiff} and of Iserles~\cite{iserles2008}, and Hairer, Lubich and Wanner~\cite{hairer2006geometric} treat the methods that respect conserved quantities such as the energy of the pendulum.

\bigskip
A linear system $\mathbf y'=A(t)\mathbf y+\mathbf b(t)$ has a unique solution for each initial value on the whole interval of continuity of its coefficients; the solutions of the homogeneous system form a space of dimension $n$, and those of the inhomogeneous system are one of them plus this space. Fundamental matrices store a basis of solutions, the transition matrix $\Phi(t,s)$ carries states from one time to another, and Liouville's formula expresses $\det\Phi(t,s)$ by the integral of the trace. For constant coefficients the transition matrix is $e^{(t-s)A}$, and the solutions are built from eigenvectors, $e^{\lambda t}\mathbf v$, from the real and imaginary parts of such solutions for complex eigenvalues, and from Jordan chains, which bring factors $t^j$ for eigenvalues that are not semisimple. In the plane the equilibrium of a linear system is a saddle, a node, a spiral, a centre, a star or a degenerate node, according to the trace and the determinant of the matrix. The equilibrium is asymptotically stable exactly when every eigenvalue has negative real part, and stable when moreover eigenvalues on the imaginary axis are allowed but must be semisimple; solutions of stable systems can grow transiently before they decay, unless the matrix is normal, and a quadratic Lyapunov function, the solution of a Lyapunov equation, certifies stability without eigenvalues. Duhamel's formula $e^{(t-t_0)A}\mathbf y_0+\int e^{(t-s)A}\mathbf b(s)\,ds$ solves the inhomogeneous system, a stable system responds to bounded inputs with bounded outputs and to periodic inputs through its frequency response $(i\omega I-A)^{-1}$. The vibrations $M\mathbf y''+K\mathbf y=\mathbf 0$ decouple into normal modes through the generalised eigenvalue problem. A nonlinear system inherits asymptotic stability and instability from its linearisation, and near a hyperbolic equilibrium its whole phase portrait, but centres must be decided by other means, such as the energy of the pendulum or the conserved quantity of the Lotka--Volterra system. Euler's method for a system is stable only when the step brings every $h\lambda$ into its region of stability, which makes stiff systems expensive for explicit methods.

The next chapter turns to linear equations with variable coefficients, for which, as \cref{sec:fundamental-matrices} showed, there is in general no formula in terms of the coefficients, and finds their solutions as power series, among them the functions of Airy, Legendre and Bessel. \Cref{ch:difference-equations} develops the discrete counterpart of this chapter, the systems of difference equations $\mathbf x_{k+1}=A\mathbf x_k$, whose stability is governed by the moduli of the eigenvalues, as that of $\mathbf y'=A\mathbf y$ is governed by their real parts.

\chapternotes

Lyapunov's thesis of 1892 on the general problem of the stability of motion introduced both methods of this chapter, the decision by the linearisation and the method of functions that decrease along the solutions; it is available in an English translation~\cite{lyapunov1992}. Routh~\cite{routh1877} and Hurwitz~\cite{hurwitz1895} found the criteria for the roots of a polynomial to lie in the left half-plane, and Gantmacher~\cite[Chap.~XV]{gantmacher1959vol2} presents both and their connection. Floquet's theory of periodic systems appeared in 1883~\cite{floquet1883}. The qualitative theory of differential equations, which studies the phase portrait rather than formulas for the solutions, began with Poincar\'e's memoirs on curves defined by differential equations, of which the first appeared in 1881~\cite{poincare1881}, and Bendixson completed the theorem that carries both their names in 1901~\cite{bendixson1901}. The predator--prey equations were proposed by Lotka~\cite{lotka1925} and, independently, by Volterra~\cite{volterra1926}. Van der Pol studied the relaxation oscillations of his equation in 1926~\cite{vanderpol1926}, the theorem on linearisation near hyperbolic equilibria is due to Grobman~\cite{grobman1959} and Hartman~\cite{hartman1960}, the example of unstable solutions with stable frozen eigenvalues to Markus and Yamabe~\cite{markus1960}, and Lorenz described his system in 1963~\cite{lorenz1963}.

Boyce, DiPrima and Meade~\cite[Chaps.~7 and~9]{boyce2017} treat linear systems and the phase plane at the level of this chapter, with many applications. Hirsch, Smale and Devaney~\cite{hirsch2013} develop linear systems, their canonical forms and phase portraits, stability, and the nonlinear theory up to chaos, and Strogatz~\cite{strogatz2015} gives a lively introduction to nonlinear systems in the plane, including the competition model of \cref{sec:nonlinear-systems} and the theory of limit cycles. Perko~\cite{perko2001} proves the local theory near equilibria, with the stable manifold and Hartman--Grobman theorems, and the Poincar\'e--Bendixson theory. Hartman~\cite{hartman2002}, Coddington and Levinson~\cite{coddington1955} and Teschl~\cite{teschl2012} are standard for the general theory, and Horn and Johnson~\cite[Chap.~2]{horn1991topics} treat stable matrices and the Lyapunov equation. Trefethen and Embree~\cite{trefethen2005} study the transient behaviour of nonnormal systems.

\chapterreferences
